% !TeX program = lualatex
% One-source manuscript. The original Nous PNG follows the end of the document.
% Run LuaLaTeX twice in a writable directory. The release uses TeX Live 2026.
% For stable timestamps, set SOURCE_DATE_EPOCH=1795582800 and FORCE_SOURCE_DATE=1.
\documentclass[12pt,letterpaper]{article}
\usepackage{iftex}\RequireLuaTeX
% Suppress the path-dependent PDF trailer ID for reproducible release builds.
\pdfvariable suppressoptionalinfo 512\relax
\usepackage[margin=.96in,headheight=16pt,headsep=14pt,footskip=28pt]{geometry}
\usepackage{fontspec}\setmainfont{Latin Modern Roman}\setsansfont{Latin Modern Sans}
\usepackage{amsmath,amssymb,mathtools,booktabs,tabularx,array,longtable,graphicx,microtype,xcolor,fancyhdr,enumitem,hyperref,bookmark}
\definecolor{ink}{HTML}{183831}\definecolor{accent}{HTML}{095D58}\definecolor{mute}{HTML}{455E55}
\hypersetup{colorlinks=true,urlcolor=accent,linkcolor=accent,pdftitle={Abundance of growth witnesses},pdfauthor={Justin Candler and Uriel}}
\setlength{\parindent}{0pt}\setlength{\parskip}{5pt plus .5pt}\setlength{\emergencystretch}{3em}
\linespread{1.045}\raggedbottom\pagestyle{fancy}\fancyhf{}\fancyhead[L]{\small\sffamily\color{mute}Nous Independent Research}\fancyhead[R]{\small\sffamily\color{mute}Abundance of growth witnesses}\fancyfoot[C]{\small\sffamily\thepage}
\renewcommand{\headrulewidth}{.3pt}\setcounter{tocdepth}{1}\setlist{itemsep=2pt,topsep=4pt}
\newcolumntype{Y}{>{\raggedright\arraybackslash}X}

\newcommand{\passmark}{\begingroup\color{accent}\rule{7pt}{7pt}\endgroup}
\newcommand{\failmark}{\begingroup\color{mute}\setlength{\fboxsep}{0pt}\fbox{\rule{0pt}{6.2pt}\rule{6.2pt}{0pt}}\endgroup}
\newcommand{\keyterm}[1]{\textbf{\color{ink}#1}}
\newcommand{\code}[1]{\texttt{\detokenize{#1}}}
\newcommand{\dossier}[2]{\clearpage\subsection*{#1}\phantomsection\label{#2}}
\newcommand{\reading}[2]{\par\smallskip\textsf{\color{accent}\bfseries #1}\enspace #2\par}
\newcommand{\ruleline}{\par\vspace{1pt}{\color{accent}\rule{\linewidth}{.35pt}}\par}
\newcommand{\yes}{\mathsf Y}\newcommand{\no}{\mathsf N}\newcommand{\unk}{\mathsf U}
\directlua{
 local source = assert(io.open(status.filename, "rb")); local rows={}; local take=false
 for line in source:lines() do
 if line:match("^BEGIN FIGURE nouslogo") then take=true
 elseif line:match("^END FIGURE") then take=false
 elseif take then table.insert(rows,line:match("^[0-9a-f]+")) end
 end
 source:close()
 function nousLogo()
 local hex=table.concat(rows); assert(string.len(hex)>100000)
 local bytes=hex:gsub("..",function(pair) return string.char(tonumber(pair,16)) end)
 local path=tex.jobname.."-embedded-nous.png"; local existing=io.open(path,"rb")
 if existing then local previous=existing:read("*a"); existing:close(); assert(previous==bytes)
 else local file=assert(io.open(path,"wb")); file:write(bytes); file:close()
 function cleanNousLogo() os.remove(path) end end
 node.write(img.node(img.scan{filename=path})) end
}

\begin{document}
\AddToHook{enddocument/afterlastpage}{\directlua{if cleanNousLogo then cleanNousLogo() end}}

\pagenumbering{roman}\thispagestyle{empty}
\begin{center}\vspace*{.52in}{\sffamily\small\bfseries\color{accent}NOUS INDEPENDENT RESEARCH}\par\vspace{.72in}
{\fontsize{31}{37}\selectfont\bfseries\color{ink}How often is the local product large?}\par\vspace{.28in}
\begin{minipage}{.86\linewidth}\centering\Large Bounded gaps, exceptional valuations, and rank-one Siegel--Weil normalization\end{minipage}
\par\vspace{.55in}{\large Justin Candler}\par With Uriel, project research assistant\par 25 November 2026
\par\vspace{.48in}\resizebox{!}{1.18in}{\directlua{nousLogo()}}\end{center}\vfill
\begin{center}A formal and experimental study of the pinned Lean source\\
\hyperref[collab]{Human and AI collaboration disclosure}\end{center}
\clearpage\phantomsection\addcontentsline{toc}{section}{Abstract}\section*{Abstract}
A rank-one Siegel--Weil argument uses an arbitrarily large dilation at which a product of local Gram-fiber masses grows faster than an analytic error. We show that the original measure hypotheses supply a whole arithmetic progression of such dilations. For a nonzero scalar target and nonzero local fiber measures, fixed positive integers $N,Q$ and a positive real $\kappa$ work at every $a=N(1+Qj)$. A cofinite estimate with summable deviations controls the moving local factors. Integral denominators and positive local patches handle the finite exceptional set. Every forward interval $(k,k+NQ]$ contains a certified dilation.

We then extend the lower estimate to all positive multipliers with bounded valuation at those exceptional places. A single constant works across the class for each fixed vector of caps. In a neighboring rank-two model the stabilized four-square density at $4^k$ is $3/4^k$, giving a kernel-checked obstruction to a cap-independent conclusion from dilation and monotonicity alone. The unrestricted assertion for rank at least three is left open here.

On the nonzero branch the scalar in the Euler-product comparison is identifiable exactly when the mass-to-error ratio is unbounded on the common test set, in the abstract comparison specified below. The original normalization needs cofinality. Supplying it with our new progression yields the same elaborated rank-one theorem type. We check the dependency path against a separate control using the original growth theorem. These results concern this intermediate arithmetic and measure-theoretic argument.
\clearpage{\small\setlength{\parskip}{0pt}\tableofcontents}\clearpage\pagenumbering{arabic}
\section{What a single local mass leaves undetermined}
The earlier calculations at targets one and four pinned down two values of the same regular dyadic fiber family. They also gave a warning about scale. The all-even part at four is $1/16$ of the mass at one. Focusing on an exact fiber was enough to identify that fraction. It left a second question untouched. As the dilation moves, every finite-place measure in an infinite product is evaluated on a set that moves with it. A calculation at one place cannot determine how often this product stays large.

Fix a number field $F$, a quadratic CM extension $E/F$ with its involution, a nondegenerate Hermitian form of source rank $M\ge3$, and a nonzero scalar $\beta\in F$. Set $D=[F:\mathbb Q]$. At a finite place $v$, take the standard additive Haar measure and an admissible rank-one Gram-fiber family for the prescribed coordinates. The growth branch assumes that the measure at the local image of $\beta$ is nonzero at each finite place. Write $L_v$ for the integral coordinate lattice. In the source, a dilation means
\[
 L_v(a)=\{X:aX\in L_v\},\qquad
 P(a)=\prod_{v<\infty} m_{v,\beta}(L_v(a))^{\rm real}.
\]
Here the superscript denotes conversion of an extended nonnegative mass to a real number. The actual arguments supply the necessary finiteness. All products and comparisons below use a fixed family of measures until their hypotheses explicitly quantify over a new family.

The prior growth proof provides $\kappa>0$ and, above each cutoff, a positive integer $a$ with $P(a)\ge\kappa a^{D(2M-2)}$. Its explicit witness sequence is geometric. Positivity of each finite local factor alone cannot give the desired conclusion: the partial products of $1/2$ tend to zero. The extra information in the original good-place theorem is precisely the rate at which each local factor approaches one. Preserving that information changes the distribution of the certified integers.

\subsection{Three questions with different failure conditions}
\begin{enumerate}
\item Can one choose $N,Q,\kappa$ before specifying the progression index? A witness depending on that index would prove a weaker assertion. An infinite product controlled only after the index was fixed could exhibit precisely that failure.
\item Can the arithmetic input be substituted into the original rank-one normalization? The analytic test and arithmetic estimate must concern the very same dilation and the same product of local masses. Cofinality of two independently selected sets would leave a gap.
\item Do exceptional places really have to see a unit multiplier? We test a wider class by keeping track of the exact loss at each exceptional place. The adjacent rank-two calculation checks what can happen when these losses have no fixed upper bound.
\end{enumerate}

The following claims were proved in Lean. The constants in the arithmetic rows are existential. The rank-two formula uses a finite counting model and holds at every specified depth.
\begin{center}\small\begin{tabularx}{\linewidth}{@{}lY@{}}\toprule
Object & Conclusion\\\midrule
Regular finite-place product & One coefficient works at all $N(1+Qj)$.\\
Gaps of the chosen progression & Every $(k,k+NQ]$ is hit, and a cutoff at a progression point waits exactly $NQ$ for its successor.\\
Exceptional valuations & For each fixed cap profile one positive coefficient works for all allowed multipliers.\\
Rank-one measure identity & Replacing the growth argument yields the original theorem type, including the zero branch.\\
Four-square control & At target $4^k$ the normalized count is $3/4^k$ once $r\ge2k+1$, for every $k\ge1$.\\\bottomrule
\end{tabularx}\end{center}

\section{From sparse witnesses to a whole progression}

Fix a number field $F$, a CM extension $E/F$ with involution $\sigma$, a nondegenerate Hermitian form $J$ in dimension $M\geq 3$, and an $F$-basis $b$ of $E$. At each finite place $v$, fix the source formalization's standard scalar Haar measure and a rank-one Gram-fiber family for the resulting coordinate measures. Write $m_v(s)$ for the member of that family at the scalar Hermitian target associated with $s$. The notation suppresses the embedding into the completion and the scalar-to-Hermitian map. Both appear explicitly in the Lean conclusion. Choose a nonzero $\beta\in F$ and suppose that $m_v(\beta)$ is a nonzero measure at every finite place.

Let $L_v$ be the original coordinate lattice. For a positive rational integer $a$, write $L_v(a)$ for the set of source vectors $X$ such that $aX\in L_v$. Put
\[
 D=[F:\mathbb Q],\qquad e=D(2M-2),\qquad
 G(a)=\prod_v m_v(\beta)(L_v(a))_{\mathbb R}.
\]
The subscript denotes the source's conversion of an extended nonnegative measure value to a real number. The product is over finite places of $F$. The original theorem supplies $c>0$ such that every cutoff $k$ admits some $a>k$ with $G(a)\geq ca^e$. Our theorem supplies positive integers $N,Q$ and a positive real $c$ such that
\[
       G\bigl(N(1+Qj)\bigr)\ \geq\ c\bigl(N(1+Qj)\bigr)^e
       \qquad \text{for every }j\in\mathbb N .
\]
The order of the quantifiers is the improvement. The integer inputs and the lower constant are fixed before the index varies. A weaker statement allowing $c$ to depend on $j$ would permit the ratio $G(a)/a^e$ to tend to zero even along our selected integers.

\subsection{Where the original proof spends its information}

The original argument selects a finite exceptional set $B$, clears a denominator, and finds a positive local patch at each exceptional place. It combines these choices in an integer $N$. Let $Q$ be the product of the ideal norms of the places in $B$ and put $p=Q+1$. Here $p$ need not be prime. It is a unit at every exceptional place. The original witnesses are $Np^k$. A multiplier $p^k$ has prime ideal divisors only among the prime ideals dividing $p$, a fixed finite set $P$ disjoint from $B$.

This restriction on prime divisors makes a simple lower estimate sufficient. Away from $P$, multiplication by $p^k$ preserves the normalized local mass at the baseline target $N^2\beta$. On $P$, a good-place theorem gives each relevant factor a lower bound of $1/2$. The unchanged factors contribute a single positive baseline product $C$, and the factors that may change cost at most $2^{-|P|}$ in this estimate. The same constant $2^{-|P|}C$ therefore works for every $k$. The construction gives the original caller an arbitrarily large witness with fixed constants. Its successive witness gaps grow with $k$.

A multiplier $1+Qj$ still avoids $B$, but its set of prime divisors now depends on $j$. This is the precise point at which changing the witness in the original argument would leave a gap. The factor $2^{-|P|}$ would have a varying exponent. For a direct test of that reasoning, imagine that in row $j$ the first $j$ factors are exactly $1/2$ and all later factors equal one. Every row contains finitely many deviations from one. All factors are positive and bounded below by $1/2$. Yet the row products are $2^{-j}$. Rowwise convergence and a half-bound can coexist with the disappearance of any common positive lower constant. This example tests the proposed inference. It is a family of real factors rather than an example of Gram-fiber measures.

The pinned argument had already established more information before deducing its half-bound. Outside a finite set, for every nonzero integral scalar target $s$, the estimate is
\[
 \big|m_v(s)(L_v)_{\mathbb R}-1\big|
       \leq \frac{2}{q_v^{M-1}}, \qquad
 q_v=\#(\mathcal O_F/\mathfrak p_v).
\]
The quantifier over $s$ lies inside a single good-place statement. We may consequently select $B$ once for this full estimate, before specifying $j$ or its associated target. Reusing the set selected only for the weaker bound would require a separate proof that the full estimate holds outside that set. The Lean construction chooses its exceptional set from the stronger theorem directly.

\subsection{A lower product independent of the row}

For $v\notin B$ put $d_v=2/q_v^2$. Every residue field has at least two elements and $M-1\geq 2$. It follows that $0<d_v\leq 1/2$ and the full error is at most $d_v$. The number-field summability result used in the certificate gives $\sum_v q_v^{-2}<\infty$. Its restriction to the complement of $B$ remains summable. The comparison factors $1-d_v$ thus have summable deviations from one and are strictly positive. The absolute-product theorem supplies
\[
              K=\prod_{v\notin B}(1-d_v)>0.
\]
In particular, comparing each factor only with $1/2$ would discard a property of the index $v$ on which this proof depends. A prime ideal of large norm has a much smaller permitted loss than one of small norm.

The general comparison lemma in our proof accepts a nonnegative row of real factors $u_j(v)$ together with positive baseline factors $b_v$ at the finitely many exceptional places. It requires $u_j(v)\geq b_v$ on $B$, $|u_j(v)-1|\leq d_v$ elsewhere, and summability of the deviations of each row. It gives the common bound
\[
       \prod_v u_j(v)\ \geq\
              \Bigl(\prod_{v\in B} b_v\Bigr)K>0 .
\]
The application will actually prove equality with the baseline on $B$. Allowing a lower bound here states the comparison at its natural level of generality. To prove it, split the infinite product into its finite exceptional part and its complement. For every finite subset of the complement, the product of the row factors dominates the product of the comparison factors. Both finite-subset products converge in the net indexed by finite subsets, so the inequality passes to their limits. The proof uses the absolute-product theorem once for the fixed comparison and again to justify splitting each row. No interchange of a limit in $j$ with a product in $v$ is involved.

\subsection{Making the exceptional local masses positive}

It remains to construct rows to which this lemma applies. For a positive integer $a$ define
\[
                 u_a(v)=m_v(a^2\beta)(L_v)_{\mathbb R}.
\]
The normalized targets here vary with $a$, while the displayed growth theorem measures the fixed target $\beta$ on varying lattices $L_v(a)$. The local dilation theorem connects these two descriptions. Before using it we must choose a baseline for which every exceptional factor is positive.

Choose $D_0>0$ with $D_0\beta\in\mathcal O_F$. The sets $L_v(n)$ for positive $n$ cover the local source space. Since $m_v(\beta)$ is a nonzero measure, for each exceptional $v$ there exists $n_v>0$ such that $m_v(\beta)(L_v(n_v))$ is nonzero. The original proof supplies this statement as a positive-patch lemma. Form
\[
             N=D_0\prod_{v\in B}n_v,\qquad
             Q=\prod_{v\in B}q_v .
\]
Empty products have value one. When $B$ is empty, $Q=1$ is permitted and the progression consists of all positive multiples of $N$. In general $N,Q$ are positive because they are finite products of positive integers.

For $v\in B$, divisibility $n_v\mid N$ gives $L_v(n_v)\subseteq L_v(N)$. The fixed-target measure of the latter set is therefore nonzero. The local scaling identity writes this measure as an explicit power of $q_v$ times the lattice measure at $N^2\beta$. Thus the latter measure is also nonzero. It is finite by the original regular-scalar local volume theorem, since $N^2\beta\ne0$. Both steps are needed in the formal argument. Converting infinity to a real number in the extended-nonnegative-real convention would give zero. Nonvanishing alone on the extended scale would not imply that the real baseline $u_N(v)$ is strictly positive.

For each positive multiple $a$ of $D_0$, the scalar $a^2\beta$ is integral at every finite place. Indeed, writing $a=D_0c$ expresses this scalar as $(D_0c^2)(D_0\beta)$, a product of integers of $F$. It is nonzero because $a$ and $\beta$ are nonzero in characteristic zero. These facts place all our varying targets inside the range of the good-prime estimate. They also give the finiteness and absolute summability needed for the normalized product in each individual row.

\subsection{Why the baseline survives at the same places for all indices}

Fix $j$ and put $t=1+Qj$, $a=Nt$. For every $v\in B$, the integer $q_v$, embedded in $\mathcal O_F$, belongs to $\mathfrak p_v$. Its divisibility into $Q$ gives $Q\in\mathfrak p_v$. If $t$ also belonged to $\mathfrak p_v$, subtracting $Qj$ would put one in this proper ideal. Thus $t$ is a unit at every exceptional place simultaneously. This argument uses prime ideals of $F$ itself and remains valid when several such ideals lie above the same rational prime.

Locally at these places $L_v(t)=L_v$. Apply the scaling identity at the baseline target $N^2\beta$ using the multiplier $t$, whose ideal exponent there is zero. The scaling factor becomes one. It follows that
\[
                      u_{Nt}(v)=u_N(v)\qquad(v\in B).
\]
All other places satisfy the same row-independent estimate $|u_{Nt}(v)-1|\leq d_v$ because $D_0$ divides $Nt$. The comparison lemma now gives one positive $c$ such that $\prod_v u_{Nt}(v)\geq c$ for all $j$. Prime divisors of $t$ may vary from row to row outside $B$. The comparison already allows every good factor to incur its full permitted loss at once.

\subsection{Recovering the original dilated-lattice product}

For any positive integer $a$, let $e_v(a)$ denote the exponent of $\mathfrak p_v$ in the principal ideal generated by $a$. The original dilation theorem expresses each factor of the fixed-target product as
\[
 m_v(\beta)(L_v(a))_{\mathbb R}
     =\big(q_v^{e_v(a)}\big)^{2M-2}\,u_a(v).
\]
Only finitely many of the norm factors differ from one, since $a$ is nonzero. Multiplying them and using the ideal norm of $a\mathcal O_F$ gives $\prod_v q_v^{e_v(a)}=a^D$. The row product is convergent by the preceding argument. Multiplicativity for products therefore yields the exact identity $G(a)=a^{D(2M-2)}\prod_v u_a(v)$. For $a=N(1+Qj)$ this is the requested lower growth bound on the original product. Its measure parameters, test lattices, target and power are those of the pinned growth theorem.

The spacing is $NQ$, regardless of $j$. This is a quantitative gain in the form of the conclusion, although the proof does not compute the numerical gap for an unspecified input. $B$ is obtained by converting an eventually-true assertion into its finite exceptional set. The denominator $D_0$ and the local positive patches are supplied by existence theorems. The positive real $c$ also involves an infinite comparison product and the actual exceptional masses. The proof gives the order in which these choices can be made and the relations they must satisfy. It gives neither a universal numerical gap across all input data nor an optimal spacing for a fixed datum. There may be many other successful dilations between the guaranteed terms. What has been established is that the growth theorem's local information suffices to place a successful dilation in every integer interval of the form $(k,k+NQ]$, with $NQ$ fixed for the input.

\subsection{What the strengthened estimate reveals about the local examples}

Two different failures can defeat an attempted product bound. On the infinite complement of $B$, the half-factor example accumulates individually moderate losses at a growing number of places. On $B$ itself, even a single unrestricted factor can destroy uniformity. For instance, let a single exceptional place have factor $1/(j+1)$ on row $j$ and let all other factors equal one. Every row is positive and its deviations are summable. The good-place comparison now holds with zero error everywhere outside that one place, yet the overall product still tends to zero. The integral denominator and the unit congruence address this second failure. The denominator first puts all good-place targets in range and the positive patches fix the exceptional baseline. The congruence then keeps that baseline available for all rows. No assertion about the number of prime divisors of $1+Qj$ can replace the latter step at a genuinely exceptional place.

This clarifies the relation to the earlier Gaussian tests. For $F=\mathbb Q$, $E=\mathbb Q(i)$ and $M=3$, the measured masses $3/4$ at one and $63/64$ at four are two possible values of a normalized local lattice factor at the place above two. If the fixed scalar is one, the integer two changes that normalized target to four. The original local dilation identity then gives a different value for the mass of the expanded lattice at the original target: its scale factor is $2^{2M-2}=16$, so the corresponding value is $16(63/64)=63/4$. Dividing by the growth power recovers $63/64$. This is a hand specialization of the identity, using the usual exponent of two in $(2)$ and the previous fiber value. The generic progression theorem does not choose two as its global $N$. In particular, inserting the exact dyadic value does not determine the values of the other local factors.

Even without knowing the numerical baseline at a particular exceptional place, its invariance under unit multipliers is enough for our family of rows. Conversely, knowing the strongest possible exact value at a single place gives no control over the accumulating losses in the infinite complement. The two earlier targets therefore measured a different part of the problem from this progression. They determined what a particular local integral means after the coordinate and measure normalizations are fixed. The progression shows which features of all those local evaluations survive the passage to a uniformly positive product. It requires finiteness, positivity, simultaneous control of a finite set, and a summable bound over the remaining places. Each requirement enters a distinct step of the proof.

\section{What a uniform spacing gives at a prescribed cutoff}
Write $d=NQ$. Because $N,Q$ are positive, $N\le d$. The index of the first progression point strictly after a cutoff $k$ is given by
\[
 j(k)=\begin{cases}0& k<N,\\1+\lfloor(k-N)/d\rfloor&k\ge N.\end{cases}
\]
Consequently $k<N+dj(k)\le k+d$. At $k=N+di$ the next point is exactly $N+d(i+1)$. This supplies both a usable inequality and an example of equality for the waiting time within the chosen progression. There could be additional admissible dilations between the progression points. No bound on their number or location is needed here.

For $K\ge N$ the interval $[0,K]$ contains $1+\lfloor(K-N)/d\rfloor$ such integers. The finite counting formula follows by solving $N+dj\le K$ for $j$. It gives the progression natural density $1/d$ and the admissible witness set lower density at least $1/d$. The density deduction is ordinary arithmetic from the certified index range. Since $d$ depends on the original input and the chosen witnesses, it supplies no numerical uniform lower density over all fields or forms.

For illustration choose $N=2$, $Q=6$. This is an integer example chosen by the authors. We did not compute the exceptional ideals of a specific Gram family to obtain these values. Both constructions start at 2. The geometric sequence multiplies the previous point by seven, whereas the evenly spaced sequence adds twelve. The two columns name the first point strictly beyond the same cutoff.
\begin{center}\begin{tabular}{rrr}\toprule
Cutoff & $2(1+6j)$ & $2\cdot7^j$\\\midrule
0&2&2\\ 2&14&14\\ 14&26&98\\ 98&110&686\\ 686&698&4,802\\ 4,802&4,814&33,614\\\bottomrule
\end{tabular}\end{center}
The finite control compared a direct search against the formula at $37{,}008$ cutoffs across $1\le N,Q\le12$. The theorem covers all positive $N,Q$. Once the numerical constants are supplied, locating the next point is integer arithmetic. The argument for the existence of $N$ itself uses nonconstructive arithmetic and measure witnesses. Separating these two operations prevents an effective search claim from slipping into the conclusion.

\section{What identifies the Euler multiplier}

Suppose the arithmetic construction has produced a positive lower bound on a progression. Where, in the original proof, does that information go? The next theorem determines a scalar that multiplies an Euler product of measures. A scalar determined at one lattice or one place would not settle this question. The same scalar must compare the adelic theta measure to the product of the fiber measures for the fixed global data. Here we follow the rank-one normalization to its conclusion. We also characterize exactly what can be learned about that scalar from the comparison inequalities alone.

\subsection{Four quantities that must refer to the same data}

Fix a nonzero scalar target $\beta\in F$ and an admissible family of local fiber measures. Write $D=[F:\mathbb Q]$ and let $M$ be the dimension of the source Hermitian space. For each positive integer $a$, let $P(a)$ be the real product over finite places of the local measure at $\beta$ of the set of vectors $X$ for which $aX$ belongs to the integral lattice. At infinity, use a strictly positive test function $\phi$ and write $w$ for the real value of its integral against the archimedean fiber measure. Fix the fundamental domain and the Haar measure defining the adelic theta measure. Finally, denote by $I(a)$ the real value of the theta integral against $\phi$ using those same finite-place constraints on $X$.

At the start of the original proof, invariance and support give an Euler-product description with some scalar $c\ge0$. There are two branches. If the archimedean or any finite fiber measure is zero, the source uses its zero-product theorem to identify the theta measure with the zero product for multiplier one. Even if two candidate multipliers give the same zero measure, their numerical values need not agree. The proof obtains the required equality of measures in this branch without identifying the initially chosen $c$. The conclusion $c=1$ belongs to the branch in which every relevant local measure is nonzero. This distinction matters whenever uniqueness is said to follow from the normalization formula. \href{https://github.com/nousentllc/riemann-hypothesis/blob/40e56bab026552e089d7afcfeff08d2defb9c005/src/CompanionTheory/AutomorphicTheory/ThetaDilations/SwOne.lean#L23-L36}{Pinned rank-one proof, lines 23--36.}

In the nonzero branch the chosen test has finite integral. Its integral has strictly positive real value. A vanishing integral of the nonnegative test would force it to vanish almost everywhere, in conflict with its strict positivity and the nonzero measure. Finiteness is needed for conversion from the extended nonnegative reals. The value infinity converts to zero in this conversion. The formal proof establishes both requirements before multiplying or dividing by $w$.

\subsection{A comparison at a common dilation}

The analytic estimate supplies a finite constant $C$ and, at every positive $a$, an observed integral with discrepancy
\[
              |I(a)-wP(a)|\le C a^{DM}.
\]
The Euler integral identity gives $I(a)=cwP(a)$. It applies to these local test sets because they are measurable, have finite fiber mass, and equal the standard lattice at almost every place. The source obtains compactness of each set by taking the inverse image of the lattice under multiplication by the nonzero local scalar $a$. Local finiteness gives finite mass. The almost-everywhere condition is a separate arithmetic property of multiplication by $a$. These steps are retained by the new proof before it invokes the scalar comparison. \href{https://github.com/nousentllc/riemann-hypothesis/blob/40e56bab026552e089d7afcfeff08d2defb9c005/src/CompanionTheory/AutomorphicTheory/ThetaDilations/SwOne.lean#L42-L80}{Pinned test and product identities, lines 42--80.}

For an arithmetic witness $a$, one also has $\kappa a^{D(2M-2)}\le P(a)$ with $\kappa>0$. In particular $P(a)\ge0$, so taking the absolute value of the Euler discrepancy gives
\[
            |c-1|\, w P(a)\le C a^{DM}.
\]
It is essential here that the analytic test, its Euler evaluation and the arithmetic estimate all refer to the same $a$. The source first chooses $a$ above an arbitrary threshold from the finite-place theorem, then invokes its analytic estimate at that $a$. The coefficient $\kappa$ is fixed before the threshold. If $c\ne1$, division by the positive quantities would bound $a^{D(M-2)}$ uniformly over these witnesses. Since $D>0$ and $M\ge3$, arbitrarily large $a$ contradict this bound. The proof ends by replacing $c$ with one in the initially given Euler-product predicate.

The progression gives witnesses for this very argument. For $N,Q\ge1$ and a cutoff $k$, take $j=k$. Then $k<1+kQ\le N(1+kQ)$. The lower bound on that progression supplies the desired positive product at the integer just selected. This step makes no appeal to Dirichlet's theorem or to primality of $1+kQ$. The original downstream application does not use the distance to the next witness. Bounded gaps describe the arithmetic abundance that the comparison is free to discard.

\subsection{An exact scalar criterion}

One can ask for a converse at the level of the two real quantities in the last inequality. Fix a set $S$ of positive test integers, a weight $w>0$, a real profile $P(a)\ge0$, and a natural exponent $q$. Allow an arbitrary finite real error constant. The inequalities on $S$ identify the multiplier uniquely as one for every such constant if and only if
\[
                   \sup_{a\in S}\frac{P(a)}{a^q}=+\infty.
\]
The expression on the right denotes unboundedness, including when $P$ oscillates. In Lean the condition is that for every real $B$ some allowed positive $a$ satisfies $Ba^q<P(a)$. This formulation also handles empty or finite allowable sets without a special convention for their supremum.

For the forward implication by contraposition, suppose $P(a)\le B a^q$ on $S$. The multiplier $c=2$ and the finite constant $C=wB$ then obey the entire family of upper comparisons. Thus those inequalities cannot identify $c$ uniquely. For the reverse implication take any $c\ne1$, and set $B=C/(|c-1|w)$. A witness for unboundedness at this $B$ gives $Ca^q<|c-1|wP(a)$, contradicting the upper bound. This argument uses an arbitrary allowable finite error constant. Knowledge of a particular exact integral or stronger error constant might identify the multiplier in a bounded-ratio example. The converse describes the information in this comparison with a finite but unspecified $C$.

The power gap used in the original proof is a convenient sufficient condition for that unboundedness. It need not be a necessary feature of an individual profile. For example, multiplying $a^q$ by an unbounded logarithm on the positive integers suffices as well. The criterion can also be satisfied by a highly sparse set of very large ratios. In the rank-one application the proof has a fixed positive archimedean weight. A stronger polynomial lower bound holds on a specified set of integer dilations, and the analytic upper error applies at every integer. The source can therefore work with one extra power of $a$ even when its actual exponent advantage is larger.

\subsection{Two ways the comparison can fail}

For a control at the boundary, set $P(a)=a^p$, take $w=\kappa=C=1$ and take $c=2$. For every positive $a$ the lower bound at exponent $p$ holds with equality. If $p\le q$, the upper discrepancy comparison also holds since $a^p\le a^q$. At $p=q$ both bounds are attained simultaneously. Thus, in the family of claims based only on two power bounds with unspecified positive coefficients, removing the strict exponent advantage can destroy normalization even when every integer is available. At $M=2$ the particular exponent difference above would vanish. The countermodel tests the scalar inference and makes no assertion about a genuine rank-two theta measure.

There is a different failure even with a strict exponent advantage. Put $P(a)=a$ on even integers and $P(a)=0$ on odd integers, again with $c=2$ and $w=1$. Arbitrarily large even integers satisfy the linear lower bound. Arbitrarily large odd integers satisfy the constant upper bound $|c-1|wP(a)\le1$. The two statements each supply a witness beyond every cutoff. At most finitely many integers satisfy both inequalities at once. The kernel checks these two countermodels independently of the adelic infrastructure. They locate two logically distinct requirements. The relative gain must exceed every fixed multiple allowed by the error, and it must do so at integers where the upper estimate applies.

\subsection{What was carried back into the formalization}

There are three levels to the certificate. The first proves the exact real-inequality criterion, the strict-power sufficient condition and the two countermodels. The second expands $P$, $w$ and $I$ into the original finite product, archimedean integral and adelic theta integral, respectively. With the discrepancy and Euler integral premises stated on those exact objects, it proves both $c=1$ and the Euler-product conclusion with multiplier one on the fixed scalar fiber. The third proof supplies these geometric premises from the imported original theorems, retains the zero-measure branch, and reconstructs the full rank-one conclusion with the finite growth supplier exposed as a parameter.

The direction of the formal dependencies is important here. The proof with the growth supplier removed imports the original theta, Euler-product and local finiteness results it needs, then uses them on the fixed data. Its original-module import closure omits the pinned rank-one normalization proof and the implementation and statement of the pinned growth theorem. Consequently the conclusion is not obtained by looking it up after deriving a general scalar lemma. The proof must still construct the theta product with undetermined scalar, transform an arbitrary nonsingular one-by-one Hermitian matrix into its scalar coordinate, build measurable finite-mass test sets, and justify evaluation of their Euler product. Omitting any of these operations would leave a scalar comparison whose variables had not been connected to the geometric measures in the conclusion.

The version with the original growth theorem is a control for that interface. It verifies that the ordinary cofinal result has precisely the type the new proof expects. The stronger progression can supply the same interface by selecting one integer above each threshold. In both instantiations the comparison is performed against the same target and local families used to define the theta measure and its Euler product. The strength gained from the progression is the location of the witnesses. The rank-one multiplier was already settled by the weaker cofinal estimate. This division of labor explains how an upstream improvement can add an arithmetic theorem while reproducing the downstream conclusion unchanged.

\section{How far can the exceptional set be relaxed?}

Fix the original number fields, Hermitian form, admissible measures, and nonzero scalar target $\beta$. Write $q_v$ for the absolute ideal norm at a finite place $v$, $M$ for the Hermitian dimension, $d_F=[F:\mathbb Q]$, and $s=2M-2$. Let $L_v$ be the integral lattice in the source's chosen basis and put $D_v(z)=m_{v,T_z}(L_v)$, taking its real value when used in estimates. The original dilation theorem relates a growing test set at one target to a fixed test set at a moving target. Our question is how much control survives when the multiplier is divisible by an exceptional prime.

\subsection{One inclusion, two different quantities}
Write $e=\operatorname{ord}_v(t)$ for the ideal exponent of a positive integer multiplier $t$. The source defines its dilated set by requiring $tX$ to lie in $L_v$. Every integral vector satisfies this test, since the image of an integer lies in the local integer ring. The defining property of a measure therefore gives
\[
    m_{v,T_z}(L_v)\ \leq\ m_{v,T_z}(\{X:tX\in L_v\}).
\]
At a regular scalar target the source's exact dilation formula evaluates the larger expression as $q_v^{es} D_v(t^2z)$. The finite-volume theorem permits conversion to real numbers. Taken together, these identities give
\[
       q_v^{-es}D_v(z)\ \leq\ D_v(t^2z).
\]
The multiplier scales both the test lattice and the target. Dropping the factor $q_v^{-es}$ would turn the true monotonicity of measured sets at fixed $z$ into a stronger monotonicity claim for a function at different targets. These assertions can have visibly different numerical behavior.

A bound $e\leq H_v$ replaces the moving factor by $q_v^{-H_v s}$, which can be chosen before $t$. This observation has value only if it can be used at all exceptional places simultaneously. There may also be arbitrarily many good places dividing different multipliers. A uniform bound for those good places requires the full summable error estimate. Using the constant one half at each newly appearing place would introduce a factor depending on the number of divisors of the multiplier.

\subsection{A common bound on a larger class of integers}
Lean proves that there exist one positive integer $N$ and one finite set of places $B$ such that the following assertion holds. Give each $v\in B$ a nonnegative integer budget $H_v$. There is a constant $c_H>0$ such that every $t\geq1$ with $\operatorname{ord}_v(t)\leq H_v$ at each $v\in B$ satisfies
\[
 c_H (Nt)^{d_F s}
 \ \leq\ \prod_v m_{v,T_\beta}(\{X:NtX\in L_v\}).
\]
Here the product is the infinite product of the real masses used by the original theorem. The integer $N$, the exceptional set, and its chosen positive baseline patches are fixed before the budgets are known. The budgets and the constant are fixed before the multiplier is supplied. This order is essential: choosing the budget from the observed factorization of $t$ would only give a separate constant for each row.

At the exceptional places, the proof first produces an integral base target $N^2\beta$ with positive mass on $L_v$. One integer clears the denominator of $\beta$. Additional factors make the dilated lattice large enough to capture one positive patch at each of the finitely many exceptional places. At a given $v$ we can then use the same test patch for every allowed multiplier. In the real-valued product an admissible choice for the lower constant is
\[
  \left(\prod_{v\in B} D_v(N^2\beta)q_v^{-H_v s}\right)
  \left(\prod_{v\notin B}(1-2/q_v^2)\right).
\]
The second factor is positive because the negative errors are summable and $q_v\geq2$. This is a witnessed formula: the finitely many local masses and the exceptional set need not be computable from the statement of the theorem. Its role is to show exactly which choices can affect the constant. In particular it identifies the cost of increasing one budget by one. The exhibited lower constant is then multiplied by $q_v^{-s}$. The actual product can exceed that floor by a large factor.

The formal comparison lemma accepts an entire sequence of rows indexed by natural numbers. To feed it every admissible $t$ at once, define $f(t)=t$ when the caps hold and $t>0$, and $f(t)=1$ otherwise. The scalar one satisfies every budget, with order zero at every place. Thus $Nf(t)$ always defines a valid row. A single use of the product lemma returns a single constant. Evaluating this result on an admissible input recovers the original $Nt$. This device prevents a subtle quantifier error: repeated applications of an existential comparison lemma, once per desired integer, need not choose a common lower constant.

A small rational example makes the size of the enlargement visible. Suppose the illustrative exceptional set consists of the rational primes two and three. The congruence $t\equiv1\pmod6$ covers six integers in a block of length $36$. Requiring $t$ to be a unit at both primes covers twelve. Allowing order at most one at both primes excludes multiples of four and nine and covers twenty-four. Our code enumerates a complete period and checks the strict inclusions of all three sets. This example does not assign the exceptional set of an actual input. It separates a property of the multiplier from one convenient way of guaranteeing it. For these two rational primes, the fraction covered with a common cap $h$ is $(1-2^{-(h+1)})(1-3^{-(h+1)})$, obtained by counting residues in the two coprime prime-power periods.

\subsection{Two controls on a possible stronger conclusion}
One might hope that positivity of the original mass and monotonicity of the dilated sets suffice even for unbounded $\operatorname{ord}_v(t)$. We tested this in two ways. The first holds the Hermitian rank at three and uses the six-square model of the Gaussian lattice. At the prime two we compared the targets three and twelve, related by the square of the multiplier two. At every tested depth from four through eleven, the normalized counts were respectively $5/4$ and $65/64$. Their ratio is $13/16$. Restoring the rank-three raw scaling factor $2^4$ changes that ratio to $13$. The finite computations isolate the direction of change: the scaled quantity grows while the normalized one falls. They give no sequence of finite densities tending to zero in rank three. They also supply no original-fiber theorem at those two targets.

The second control lowers the Hermitian rank to two, so there are four rational Gaussian coordinates. This places it outside the $M\geq3$ theorem and gives a setting in which the proposed inference can be tested exactly. Let $R_r(b)$ denote the number of four-square representations of $b$ modulo $2^r$. Lean proves for every $k\geq1$ and every $r\geq2k+1$ that
\[
       \frac{R_r(4^k)}{(2^r)^3}=\frac{3}{4^k}.
\]
There is no positive lower bound on these normalized counts along the sequence. Multiplying by the rank-two scaling factor $(2^k)^2$ gives exactly three for each of them. One can therefore maintain a positive raw quantity at every step while the normalized density tends to zero.

The mechanism is visible modulo eight. An odd square contributes one modulo eight. Modulo four, a zero sum of four squares forces either zero or four odd terms. Four odd squares sum to four modulo eight; hence every solution at zero modulo eight has all four coordinates even. Exhaustive enumeration gives $128$ all-even solutions and zero primitive ones. For $k\geq2$ the target $4^k$ is zero modulo eight, so every solution in a sufficiently fine quotient belongs to the all-even sublattice. Halving its four coordinates produces sixteen lifts of each solution at the divided target two levels earlier. The normalization changes by $4^3=64$, leaving a factor $1/4$ at each step. The initial density $3/4$ at four is certified by finite histograms at moduli eight, sixteen, and thirty-two, and a unit-square action beyond thirty-two. This proves the full displayed family of exact counts. A second theorem turns it directly into the assertion that every candidate positive lower constant fails, with a suitable $k$ and all depths in the specified range.

For a by-hand check consider the rank-two target sixteen. At modulus thirty-two every solution can be written $x=2y$ with $y$ defined modulo sixteen. Its reduction modulo eight must have norm four. There are $384$ such reduced vectors: $256$ have four odd entries, and $128$ have all entries even with an odd number of squares equal to four modulo eight. All sixteen upper-bit choices for $y$ give distinct original vectors. Thus $R_5(16)=6{,}144$, or $3/16$ of $32^3$. At target sixty-four the same division recurs; the density is $3/64$, and the restoring raw factor is now sixty-four. This pair of computations detects two independent errors: discarding the free high bit in each halved coordinate and carrying over the denominator exponent five from the rank-three model. The Lean theorem uses dimension four and normalization exponent three throughout.

\subsection{What the controls leave open}
These conclusions distinguish three levels of evidence. The capped growth theorem and its local comparison refer to the actual measures, test sets, and exponent of the pinned rank-at-least-three proof. The rank-two counterexample is a checked theorem about finite four-square cardinalities and their stabilized normalized densities; our argument did not construct a rank-two original fiber measure. The six-square comparison at three and twelve is a finite experiment with exact integers. It falsifies a pointwise monotonicity prediction for those finite normalized counts.

For the original arbitrary Hermitian form in rank at least three, a cap-independent bound at an exceptional place would need another source of mass as $t^2N^2\beta$ moves towards zero. A natural test would identify a regular local representation of zero inside a fixed integral patch and prove a uniform positive lower bound on its contribution for nearby nonzero targets. Such an argument would have to apply to the original lattice and measure normalization, including the ramified dyadic cases. The current certificate proves no such near-zero estimate. Its loss factor shows precisely where a new local argument could improve the global conclusion. Increasing the caps one at a time can cover each given integer eventually, but the constants obtained in that way can tend to zero. This exchange of quantifiers is the remaining obstacle to asserting one common constant for all positive multipliers.

\section{What is gained by combining the three results?}
One can now hold the field, basis, form, scalar and admissible measures fixed from the local choices all the way to the rank-one equality. There are three reasons the substitution works. First the positive coefficient survives an infinite product uniformly in $j$. Second the analytic test is available at each of the positive integer dilations supplied by the arithmetic. Third the two powers differ by $D(M-2)>0$. A stronger result at any one of these steps would be useless if its witnesses did not match those at the others.

The term comparison provides a check on the structure of the argument. An import graph by itself would be ambiguous. The module containing the original growth proof also contains local scaling and finiteness lemmas used by the new argument. An import edge alone could suggest that our progression simply calls the sparse growth bound. We instead inspected the types and values of elaborated declarations, starting with the new normalization and separately with a control that uses the old provider. The new route uses the progression and its product comparison. The control uses the original growth theorem. Both compare to the very same type of the original rank-one measure identity. The graph audit checks the membership of these specific declarations. The number of other visited declarations has no interpretation as a measure of mathematical understanding.

\begin{center}\small
\begin{tabularx}{.95\linewidth}{@{}YY@{}}\toprule
\textbf{Progression proof} & \textbf{Original-growth control}\\\midrule
Summable good-place deviations and finitely many fixed positive local masses
& Original local growth theorem and a geometric witness sequence\\[2pt]
$\Downarrow$ & $\Downarrow$\\
A common lower bound on $N(1+Qj)$ & A common lower bound at cofinally many $a$\\[2pt]
$\Downarrow$ & $\Downarrow$\\
\multicolumn{2}{c}{Same positive archimedean test and same finite-place product at each chosen $a$}\\
\multicolumn{2}{c}{$\Downarrow$}\\
\multicolumn{2}{c}{Same original rank-one measure identity}\\\bottomrule
\end{tabularx}\end{center}

The zero-fiber branch precedes this diagram and needs no comparison of growth rates. On the other branch, the progression supplies more dilations than the comparison requires. Our two proof terms agree in their conclusion. Exactly which lower estimate each term consumes is visible in its elaborated dependencies.

There is also a difference between proving existence and finding examples effectively. If all local factors were explicit at infinitely many places but one exceptional positive patch remained unknown, the generic argument would still produce $N$ by choice. Conversely, knowing $N,Q$ would tell us an arithmetic progression without giving the best value of $\kappa$. The quantitative six-square calculation at two does not remove this distinction. It evaluates particular tests in one chosen local field. The growth theorem quantifies over a family of fields and forms and determines what those tests have to accomplish collectively.

\subsection{How much of the positive-integer axis do the caps allow?}
To see the difference between a congruence restriction and a bound on valuation, take an illustrative exceptional set $\{2,3\}$ over $\mathbb Q$. For uniform cap $h$, the first forbidden powers are $2^{h+1}$ and $3^{h+1}$. They are coprime. In a complete period of length $6^{h+1}$ the allowed fraction is
\[
 \left(1-2^{-h-1}\right)\left(1-3^{-h-1}\right).
\]
For caps $0,1,2,3,4$ the exact fractions are $1/3,2/3,91/108,25/27,3751/3888$. We checked every residue in each period. The class $1\bmod6$ alone has fraction $1/6$. This example shows why bounded valuations can substantially enlarge the class of multipliers. The theorem supplies a different permissible coefficient for each profile. Letting the allowed fraction tend to one in this example says nothing about a positive lower bound for those coefficients.

\subsection{The next mathematical test}
At an exceptional place, removing the cap allows the scaled target to approach zero. In an admissible rank-three example a promising test would be to exhibit an integral nonzero isotropic vector, find a compact open neighborhood on which the Gram map covers a neighborhood of zero with a uniform positive density, and connect that test to the original regular fiber family for nearby \emph{nonzero} targets. Such a result would address the point at which the cap-dependent proof loses its coefficient. A statement about the already assigned measure at the singular target would not supply that comparison. This is an open verification question in this study. The rank-two control is a diagnosis of a missing hypothesis in a general scaling argument.

\section{Reading the formal certificates} The following extracts give the exact module names and source line numbers in the package. Mathematical symbols in the extracts are transliterated. The original Unicode sources are included. The proof obligations range from uniform product estimates to the final type and proof-dependency checks.
\smallskip\phantomsection\pdfbookmark[2]{A floor for an infinite family of products}{plate-ProgressionGrowth13}\subsection*{A floor for an infinite family of products}\label{extract-ProgressionGrowth13}
The local data form a doubly indexed array. The first index moves the dilation and the second moves the finite place. Lower bounds on a fixed row allow a constant to depend on that row. This lemma quantifies one coefficient before every row. Its three inputs distinguish the finite exceptional set, the uniform estimate on its complement, and the convergence of each actual product.
\smallskip{\small\sffamily\color{accent}Proof extract \quad ProgressionGrowth.lean \quad lines 13--31}\par
{\fontsize{9.2}{11.2}\selectfont\begin{verbatim}
 13      (hbad : forall  j v, v  in  B  ->  base v  <=  u j v)
 14      (hpos : forall  v  in  B, 0 < base v)
 15      (hgood : forall  j v, v  notin  B  ->
 16        |u j v - 1|  <=  2 / (Ideal.absNorm v.asIdeal : Real) ^ 2)
 17      (hsum : forall  j, Summable fun v => |u j v - 1|) :
 18      exists  kappa : Real, 0 < kappa  /\  forall  j, kappa  <=  Prod' v, u j v := by
 19    classical
 20    let V := {v : HeightOneSpectrum (O F) // v  notin  B}
 21    let d : V  ->  Real := fun v => 2 / (Ideal.absNorm v.1.asIdeal : Real) ^ 2
 22    have hd : Summable d := by
 23      have h := ((Primes.summable_one_div_absNorm_sq F).mul_left 2).comp_injective
 24        (i := fun v : V => v.1) Subtype.val_injective
 25      exact h.congr fun v => by simp [d, Function.comp_apply, div_eq_mul_inv]
 26    have hd0 (v : V) : 0  <=  d v := by dsimp [d]; positivity
 27    have hd1 (v : V) : d v < 1 := by
 28      have hq : (2 : Real)  <=  Ideal.absNorm v.1.asIdeal := by exact_mod_cast
     Growth.two_le_norm v.1
 29      dsimp [d]
 30      rw [div_lt_iff0 (by positivity)]
 31      nlinarith
\end{verbatim}}
The comparison factors depend only on the place. Since $q_v\ge2$, each factor $1-2/q_v^2$ is strictly positive even at the smallest possible residue-field size. Summability controls the infinite tail. Looking only at the infimum of the individual factors would conceal the possibility that their product vanishes.
\clearpage\phantomsection\pdfbookmark[2]{The finite set is chosen before the index}{plate-ProgressionGrowth56}\subsection*{The finite set is chosen before the index}\label{extract-ProgressionGrowth56}
The source quantifies over the finite-place Haar measures and the Gram-fiber families with their prescribed normalizations. For a nonzero global scalar, the needed positivity hypothesis concerns the whole local measure at its image. From the original good-place theorem the proof extracts a finite set outside which the full deviation estimate holds for every relevant integral target.
\smallskip{\small\sffamily\color{accent}Proof extract \quad ProgressionGrowth.lean \quad lines 56--81}\par
{\fontsize{9.2}{11.2}\selectfont\begin{verbatim}
 56  theorem progression_growth (hsig : IsCMExt F E sig) (hJ : IsHermND sig J) (hM : 3
     <=  M) (hlam : forall  v, IsStdV v (lam v))
 57      (hm : forall  v, IsGramFib (Fv F v) sig J 1 (coordW (Fv F v) b (lam v) 1 M)
     (coordHerm (Fv F v) sig b (lam v) 1) (m v))
 58      (beta : F) (hbeta : beta  !=  0) (hne : forall  v, m v (hermScal (Fv F v) sig
     (algebraMap F _ beta))  !=  0) :
 59      exists  N Q : Nat, 0 < N  /\  0 < Q  /\  exists  c : Real, 0 < c  /\  forall  j
     : Nat,
 60        let a := N * (1 + Q*j)
 61        c * (a : Real) ^ (Module.finrank Rat F * (2*M-2))  <=
 62          Prod' v, (m v (hermScal (Fv F v) sig (algebraMap F _ beta)) (dilSet b v
     a)).toReal := by
 63    classical
 64    -- Choose the exceptional places for the full summable estimate.
 65    have hgood := GlobalTheory.T.rank_one_local_density_almost_all F E sig hsig M J
     hJ hM d b
 66    let B := (Filter.eventually_cofinite.1 hgood).toFinset
 67    have hBg (v : HeightOneSpectrum (O F)) (hv : v  notin  B) (s : Fv F v)
 68        (hs : s  in  v.adicCompletionIntegers F) (hs0 : s  !=  0) :
 69        |(m v (hermScal (Fv F v) sig s) (latV b v)).toReal - 1|  <=
 70          2 / (Ideal.absNorm v.asIdeal : Real)^2 := by
 71      simp only [B,Set.Finite.mem_toFinset,Set.mem_ofPred_eq,not_not] at hv
 72      have hfull := hv (lam v) (hlam v) (m v) (hm v) s hs hs0
 73      have hq : (2 : Real)  <=  Ideal.absNorm v.asIdeal := by exact_mod_cast
     two_le_norm v
 74      have hqpos : (0 : Real) < Ideal.absNorm v.asIdeal := by linarith
 75      apply le_trans hfull
 76      apply div_le_div_of_nonneg_left (by norm_num) (by positivity)
 77      apply pow_le_pow_right0 _ (by omega : 2  <=  M-1)
 78      linarith
 79    choose nv hnv0 hnv using fun v => exists_pos_mass b v _ (hne v)
 80    obtain <D0, hD0, x, hx> := exists_den beta
 81    set N := D0 * Prod v  in  B, nv v with hNdef
\end{verbatim}}
If that finite set could vary with the progression index, multiplying its norms once would not protect later rows. Here it comes from a single cofinite assertion for all the admissible local targets. This quantifier is the first piece of the proof that makes a regularly spaced family possible.
\clearpage\phantomsection\pdfbookmark[2]{Keeping an exceptional factor fixed}{plate-ProgressionGrowth104}\subsection*{Keeping an exceptional factor fixed}\label{extract-ProgressionGrowth104}
At an exceptional place the integer $N$ is divisible by a local positive-patch dilation. Positivity is transferred through inclusion of the corresponding expanded lattices. Local scaling then places this positivity at target $N^2\beta$. Finite mass allows the proof to pass to a positive real local density.
\smallskip{\small\sffamily\color{accent}Proof extract \quad ProgressionGrowth.lean \quad lines 104--128}\par
{\fontsize{9.2}{11.2}\selectfont\begin{verbatim}
104
105    have hDN : D0  divides  N := Dvd.intro _ rfl
106    have hpos (v : HeightOneSpectrum (O F)) (hv : v  in  B) :
107        0 < (mL sig b m ((N : F)^2 * beta) v).toReal := by
108      obtain <e, he, -, -> := exists_exponents (F := F) N hN
109      have h := dil_eq (b := b) hsig hJ hM hlam hm v beta N (e v) hN (he v).1 (he
     v).2
110      have hnz : m v (hermScal (Fv F v) sig (algebraMap F _ beta)) (dilSet b v N)  !=
     0 := by
111        obtain <c, hc> : nv v  divides  N := (Finset.dvd_prod_of_mem nv hv).mul_left
     _
112        intro h0
113        refine hnv v (measure_mono_null ?_ h0)
114        rw [hc]; exact dilSet_mono b v _ _
115      have hz : mL sig b m ((N : F) ^ 2 * beta) v  !=  0 := by
116        intro h0
117        rw [h, h0, mul_zero] at hnz
118        exact hnz rfl
119      exact ENNReal.toReal_pos hz ((hfin N hN).1 v).ne
120    let A : Nat  ->  Nat := fun j => N*(1+Q*j)
121    have hA (j : Nat) : 0 < A j := by dsimp [A]; positivity
122    let u (j : Nat) (v : HeightOneSpectrum (O F)) := (mL sig b m ((A j : F)^2 * beta)
     v).toReal
123    have hfixed (j : Nat) (v : HeightOneSpectrum (O F)) (hv : v  in  B) :
124        (mL sig b m ((N : F)^2 * beta) v).toReal  <=  u j v := by
125      let t := 1+Q*j
126      have ht : 0 < t := by dsimp [t]; omega
127      have htunit : (t : O F)  notin  v.asIdeal := hunit v hv j
128      have h := dil_eq (b := b) hsig hJ hM hlam hm v ((N : F)^2 * beta) t 0 ht (by
     simp)
\end{verbatim}}
The multiplier $1+Qj$ is a unit at this place. It leaves the coordinate lattice unchanged, and the exponent in the local dilation identity is zero. Consequently the normalized factor has its baseline value for all $j$. This is stronger than finding a positive mass after choosing each new target separately.
\clearpage\phantomsection\pdfbookmark[2]{The exact forward waiting time}{plate-ProgressionGeometry5}\subsection*{The exact forward waiting time}\label{extract-ProgressionGeometry5}
At a cutoff below the start, index zero suffices because the start is at most the increment. At a later cutoff, Euclidean division of the displacement from the start gives the index just before or at the cutoff. The proof advances one index so that the lower inequality is strict even when the cutoff is a progression point.
\smallskip{\small\sffamily\color{accent}Proof extract \quad ProgressionGeometry.lean \quad lines 5--31}\par
{\fontsize{9.2}{11.2}\selectfont\begin{verbatim}
  5  /-- An affine progression with starting point no larger than its increment meets
  6  every forward interval of length that increment, including the initial interval. -/
  7  theorem bounded_wait (n d k : Nat) (hd : 0 < d) (hn : n  <=  d) :
  8      exists  j : Nat, k < n + d*j  /\  n + d*j  <=  k+d := by
  9    by_cases hk : k < n
 10    * exact <0, by simpa, by simp; omega>
 11    let t := k - n
 12    have hr := Nat.mod_lt t hd
 13    have hdiv := Nat.mod_add_div t d
 14    have ht : t + n = k := Nat.sub_add_cancel (by omega)
 15    refine <t/d + 1, ?_, ?_> <;> rw [Nat.mul_add, Nat.mul_one] <;> omega
 16
 17  /-- For cutoffs equal to progression points the forward delay is exactly d. -/
 18  theorem sharp_at_member (n d i j : Nat) (h : n+d*i < n+d*j) :
 19      n+d*i+d  <=  n+d*j := by
 20    have hi : i < j := by
 21      by_contra! hc
 22      exact (not_le_of_gt h) (Nat.add_le_add_left (Nat.mul_le_mul_left d hc) n)
 23    calc n+d*i+d = n+d*(i+1) := by ring
 24         _  <=  n+d*j := Nat.add_le_add_left (Nat.mul_le_mul_left d hi) n
 25
 26  /-- All good indices below a cutoff, expressed without a search. -/
 27  theorem cutoff_indices (n d k j : Nat) (hd : 0 < d) (h : n  <=  k) :
 28      n + d*j  <=  k  <->  j  <=  (k-n)/d := by
 29    rw [Nat.le_div_iff_mul_le hd, Nat.mul_comm j d]
 30    omega
 31
\end{verbatim}}
Sharpness here has a specific meaning. A progression member followed by the next progression member realizes the whole increment. An argument that also certifies additional integers can reduce the wait. The separate capped-valuation theorem is an example of a larger supply of integers, although it does not identify the optimal set of witnesses.
\clearpage\phantomsection\pdfbookmark[2]{The comparison characterized as an implication}{plate-GrowthConsumer7}\subsection*{The comparison characterized as an implication}\label{extract-GrowthConsumer7}
The abstract question quantifies over every permissible multiplier and error constant. From an unbounded product-to-error ratio, any nonzero discrepancy would produce a finite upper bound for that same ratio. Conversely, if the ratio is bounded, a multiplier displaced from one can satisfy a suitably chosen finite error bound. This second direction explains the quantifier on the error constant.
\smallskip{\small\sffamily\color{accent}Proof extract \quad GrowthConsumer.lean \quad lines 7--33}\par
{\fontsize{9.2}{11.2}\selectfont\begin{verbatim}
  7  /-- At positive allowed integers, the lower profile exceeds every fixed multiple
  8  of the analytic error's power. -/
  9  def RelativeGrowth (allowed : Nat  ->  Prop) (P : Nat  ->  Real) (q : Nat) : Prop
     :=
 10    forall  B : Real, exists  a : Nat, 0 < a  /\  allowed a  /\  B * (a : Real)^q < P
     a
 11
 12  /-- The matched scalar inequality identifies the multiplier for every finite error
     constant. -/
 13  def ForcesOne (allowed : Nat  ->  Prop) (P : Nat  ->  Real) (q : Nat) (w : Real) :
     Prop :=
 14    forall  c C : Real, (forall  a : Nat, 0 < a  ->  allowed a  ->
 15      |c-1| * (w * P a)  <=  C * (a : Real)^q)  ->  c = 1
 16
 17  /-- Exact criterion at the level of real inequalities. If the relative growth
     fails,
 18  the multiplier 2 and a finite error constant already give a counterexample. -/
 19  theorem forces_iff_relativeGrowth {allowed : Nat  ->  Prop} {P : Nat  ->  Real} {q
     : Nat}
 20      {w : Real} (hw : 0 < w) : ForcesOne allowed P q w  <->  RelativeGrowth allowed
     P q := by
 21    constructor
 22    * intro force
 23      unfold RelativeGrowth
 24      by_contra bad
 25      push Not at bad
 26      obtain <B, hB> := bad
 27      have absurd := force 2 (w * B) (by
 28        intro a ha haS
 29        have h := mul_le_mul_of_nonneg_left (hB a ha haS) hw.le
 30        norm_num [mul_assoc] at h |-
 31        exact h)
 32      norm_num at absurd
 33    * intro growth c C hbound
\end{verbatim}}
The equivalence concerns what follows from this inequality with a fixed positive archimedean weight. It does not state which arbitrary functions can arise as local Gram products. The literal measure proof invokes the implication only after constructing the actual test integrals and the appropriate estimates from the original data.
\clearpage\phantomsection\pdfbookmark[2]{Two countermodels with different failure modes}{plate-GrowthConsumer89}\subsection*{Two countermodels with different failure modes}\label{extract-GrowthConsumer89}
At equal or reversed exponents the normalized growth can remain bounded forever. In a second example there are arbitrarily good growth witnesses and arbitrarily many analytic tests, but they occupy disjoint classes of integers. Either situation allows a scalar different from one to survive the purported comparison.
\smallskip{\small\sffamily\color{accent}Proof extract \quad GrowthConsumer.lean \quad lines 89--116}\par
{\fontsize{9.2}{11.2}\selectfont\begin{verbatim}
 89    exact_mod_cast heq
 90
 91  /-- With a nonpositive exponent gap, the admissible multiplier 2 obeys both bounds
 92  at every positive integer in the unit-weight model. -/
 93  theorem boundary_countermodel (p q : Nat) (hpq : p  <=  q) :
 94      (forall  k : Nat, exists  a : Nat, k<a  /\  1 * (a : Real)^p  <=  (a : Real)^p)
     /\
 95      (forall  a : Nat, 0<a  ->  |(2:Real)-1| * (1 * (a : Real)^p)  <=  1*(a :
     Real)^q)  /\  (2:Real)  !=  1 := by
 96    refine <fun k => <k+1, by omega, by simp>, ?_, by norm_num>
 97    intro a ha
 98    have h : (1:Real)  <=  a := by exact_mod_cast ha
 99    norm_num
100    exact pow_le_pow_right0 h hpq
101
102  /-- A linear lower witness on evens and a bounded upper witness on odds give no
103  information about a shared multiplier. -/
104  theorem unmatched_countermodel :
105      exists  P : Nat  ->  Real,
106        (forall  k : Nat, exists  a : Nat, k<a  /\  (a : Real)  <=  P a)  /\
107        (forall  k : Nat, exists  a : Nat, k<a  /\  |(2:Real)-1| * (1 * P a)  <=  1)
     /\  (2:Real)  !=  1 := by
108    let P : Nat  ->  Real := fun a => if a % 2 = 0 then (a : Real) else 0
109    refine <P, ?_, ?_, by norm_num>
110    * intro k
111      refine <2*(k+1),by omega,?_>
112      simp [P]
113    * intro k
114      refine <2*(k+1)+1,by omega,?_>
115      simp [P]
116
\end{verbatim}}
These controls change different hypotheses. Improving the exponent cannot repair the disjoint-witness example when the analytic inequality is never available where the growth occurs. Adding more tests cannot create a positive exponent difference if the mass is exactly of the same order as the allowed analytic error.
\clearpage\phantomsection\pdfbookmark[2]{The two branches of the literal measure identity}{plate-RankOneConsumer52}\subsection*{The two branches of the literal measure identity}\label{extract-RankOneConsumer52}
The proof begins with the source's Euler-product measure, characterized up to a multiplier. If any required local measure is zero, the zero-product theorem proves the target identity. In the other branch a rank-one Hermitian target is converted into its nonzero scalar parameter and the provider is applied to that same parameter.
\smallskip{\small\sffamily\color{accent}Proof extract \quad RankOneConsumer.lean \quad lines 52--72}\par
{\fontsize{9.2}{11.2}\selectfont\begin{verbatim}
 52  theorem sw_one_of_growth (provider : RankOneGrowthProvider) : RankOneConclusion :=
     by
 53    intro F E _ _ _ _ _ sig hcm M J hJ hM T hT hd d b li lam hli hlam mu FD hmu hFD
     mi m hmi hm
 54    obtain <t1, t2> := theta_inv F E sig hcm M J hJ mu FD hmu hFD 1 T
 55    obtain <c, hc> :=
     GlobalTheory.T.exists_isEulerProd_of_invariant_of_fiber_supported F E sig hcm M J
     hJ 1 (Or.inr hM) T hT hd d b li lam hli hlam mi hmi m hm _
 56      (theta_measure_locally_finite F E sig hcm M J hJ hM mu FD hmu hFD T hT hd) t1
     t2
 57    by_cases h0 : mi (hermOf (Ainf F) sig T hT) = 0  \/  exists  v, m v (hermOf (Fv F
     v) sig T hT) = 0
 58    * rw [(GlobalTheory.T.euler_product_measure_zero F E sig hcm d b 1 M _ c _ _).2.1
     hc h0]
 59      exact (GlobalTheory.T.euler_product_measure_zero F E sig hcm d b 1 M 0 1 _
     _).2.2 hc.1 hc.2.1 h0
 60    suffices h : c = 1 from h  transport  hc
 61    obtain <beta, rfl, hbeta, he> := GlobalTheory.T.hermitian_one_by_one_is_scalar F
     E sig hcm T hT
 62    simp only [he] at hc h0
 63    push Not at h0
 64    obtain <phi, C, phic, phip, phif, hb> :=
     GlobalTheory.T.thetaMeas_rank_one_sub_density_product_le F E sig hcm M J hJ hM d b
     li lam hli hlam mu FD hmu hFD mi hmi m hm beta (hbeta hd)
 65    obtain <kappa, hkappa, hg> := provider F E sig hcm M J hJ hM d b lam hlam m hm
     beta (hbeta hd) h0.2
 66    have : BorelSpace (Fin 1  ->  Fin M  ->  Ainf F  tensor [F] E) := by
 67      have : BorelSpace (Ainf F  tensor [F] E) := <rfl>
 68      obtain <_, _, _, _, -> :=
     GlobalTheory.T.infinite_adeles_local_ring_totally_real F hcm.1
 69      exact (GlobalTheory.T.tensor_locally_compact_ring F E sig hcm (Ainf
     F)).2.2.2.2.2.2.2.1 1 M
 70    have hphim : Measurable fun X => ENNReal.ofReal (phi X) :=
     phic.measurable.ennreal_ofReal
 71    refine (original_scalar_comparison sig J b mu FD mi m beta phi hM hc h0.1 hphim
     phip phif
 72      hkappa hg hb (fun a ha0 => ?_)).1
\end{verbatim}}
This separation explains an apparent paradox. The conclusion is stated everywhere under the original hypotheses, although the positivity needed for uniqueness can fail on one branch. What holds on the zero branch is the measure identity written with multiplier one. Several scalar multipliers can produce that same zero measure.
\clearpage\phantomsection\pdfbookmark[2]{Where one theorem is substituted for another}{plate-ProgressionEndpoint10}\subsection*{Where one theorem is substituted for another}\label{extract-ProgressionEndpoint10}
The measured product, rank, degree, and target appear in the type of the provider. The progression theorem supplies $N,Q,c$ for those very inputs. The waiting-time result selects a progression point after the requested cutoff. At this boundary the extra upper bound on the wait is deliberately discarded because the receiving theorem asks only for a point above the cutoff.
\smallskip{\small\sffamily\color{accent}Proof extract \quad ProgressionEndpoint.lean \quad lines 10--24}\par
{\fontsize{9.2}{11.2}\selectfont\begin{verbatim}
 10  /-- Supply the actual finite-place input of the original rank-one comparison
 11  with a uniformly spaced sequence. The comparison asks only for cofinality. -/
 12  theorem provider : GrowthConsumer.RankOneGrowthProvider := by
 13    intro F E _ _ _ _ _ sig hsig M J hJ hM d b lam hlam m hm beta hbeta hne
 14    obtain <N,Q,hN,hQ,c,hc,hl> := ProgressionGrowth.progression_growth (b := b)
 15      hsig hJ hM hlam hm beta hbeta hne
 16    refine <c,hc,fun k => ?_>
 17    obtain <a,ha,_,hp> := ProgressionGeometry.from_progression hQ hN
 18      (P := fun a => c*(a:Real)^(Module.finrank Rat F*(2*M-2))  <=
     GrowthConsumer.finiteProduct sig b m beta a) hl k
 19    exact <a,ha,hp>
 20
 21  theorem sw_one_via_progression : GrowthConsumer.RankOneConclusion :=
 22    GrowthConsumer.sw_one_of_growth provider
 23
 24  end ProgressionEndpoint
\end{verbatim}}
The final declaration feeds that provider into the original-type normalization argument. It is possible to inspect the statement of an intermediate result and still miss an error in the types of its maps or measures. Here the substitution is an elaborated term. The endpoint type and the declarations that its proof actually uses are checked separately.
\clearpage\phantomsection\pdfbookmark[2]{Accounting for the cost of a nonunit}{plate-ExceptionalDilation27}\subsection*{Accounting for the cost of a nonunit}\label{extract-ExceptionalDilation27}
An integral vector remains integral after multiplication by a positive rational integer, so the original lattice sits inside its expanded counterpart. This set inclusion gives a monotonicity inequality for the raw measure. To compare normalized local factors one must divide out the exact power of the residue-field size in the dilation identity.
\smallskip{\small\sffamily\color{accent}Proof extract \quad ExceptionalDilation.lean \quad lines 27--48}\par
{\fontsize{9.2}{11.2}\selectfont\begin{verbatim}
 27        (mL sig b m ((a : F)^2*s) v).toReal := by
 28    let mu := m v (hermScal (Fv F v) sig (algebraMap F _ s))
 29    have hsub : (latV b v : Set (Fin 1  ->  Fin M  ->  Fv F v  tensor [F] E))  subset
     dilSet b v a := by
 30      intro X hX
 31      apply (mem_dilSet b v).mpr
 32      intro i j k
 33      exact mul_mem (natCast_mem (v.adicCompletionIntegers F) a) (hX i j k)
 34    have hscale := dil_eq (b := b) hsig hJ hM hlam hm v s a e ha h1 h2
 35    have hn : ((a : F)^2*s)  !=  0 := mul_ne_zero (pow_ne_zero _ (by exact_mod_cast
     ha.ne')) hs
 36    have hfinite := ((scal_fin (b := b) hsig hJ hM hlam hm hn).1 v).ne
 37    have hraw : mu (dilSet b v a)  !=  top := by
 38      dsimp [mu]; rw [hscale]; exact ENNReal.mul_ne_top (by finiteness) hfinite
 39    have hh := ENNReal.toReal_mono hraw (measure_mono (mu := mu) hsub)
 40    change (mL sig b m s v).toReal  <=  (mu (dilSet b v a)).toReal at hh
 41    dsimp only [mu] at hh
 42    rw [hscale,ENNReal.toReal_mul,ENNReal.toReal_pow,ENNReal.toReal_natCast] at hh
 43    have hq : (2:Real)  <=  (Ideal.absNorm v.asIdeal : Real) := by exact_mod_cast
     two_le_norm v
 44    have hp : 0 < (Ideal.absNorm v.asIdeal : Real)^(H*(2*M-2)) := pow_pos (by
     linarith) _
 45    rw [div_le_iff0 hp,mul_comm]
 46    exact hh.trans (mul_le_mul_of_nonneg_right
 47      (pow_le_pow_right0 (by linarith : 1  <=  (Ideal.absNorm v.asIdeal : Real))
 48        (Nat.mul_le_mul_right _ he)) ENNReal.toReal_nonneg)
\end{verbatim}}
An upper bound on the local order turns the possible loss into a fixed positive coefficient. The cost depends on the place, its cap and the source rank. It applies equally to all allowed multipliers at that place. This is the step that fails to give a coefficient independent of an arbitrarily increasing cap.
\clearpage\phantomsection\pdfbookmark[2]{Quantifying which multipliers are admitted}{plate-CappedGrowth21}\subsection*{Quantifying which multipliers are admitted}\label{extract-CappedGrowth21}
The theorem chooses the baseline dilation and finite set first. A cap function is then given. Only then does the proof select the common coefficient and quantify over the multipliers. This order permits the admissible set to grow with the chosen profile, and the guaranteed coefficient to deteriorate.
\smallskip{\small\sffamily\color{accent}Proof extract \quad CappedGrowth.lean \quad lines 21--39}\par
{\fontsize{9.2}{11.2}\selectfont\begin{verbatim}
 21
 22  theorem bounded_order_growth (hsig : IsCMExt F E sig) (hJ : IsHermND sig J) (hM : 3
     <=  M) (hlam : forall  v, IsStdV v (lam v))
 23      (hm : forall  v, IsGramFib (Fv F v) sig J 1 (coordW (Fv F v) b (lam v) 1 M)
     (coordHerm (Fv F v) sig b (lam v) 1) (m v))
 24      (beta : F) (hbeta : beta  !=  0) (hne : forall  v, m v (hermScal (Fv F v) sig
     (algebraMap F _ beta))  !=  0) :
 25      exists  (N : Nat) (B : Finset (HeightOneSpectrum (O F))), 0 < N  /\
 26        forall  H : HeightOneSpectrum (O F)  ->  Nat, exists  c : Real, 0 < c  /\
 27          forall  t : Nat, 0<t  ->  (forall  v in B, WithinCap v t (H v))  ->
 28            let a := N*t
 29            c*(a:Real)^(Module.finrank Rat F*(2*M-2))  <=
 30              Prod' v, (m v (hermScal (Fv F v) sig (algebraMap F _ beta)) (dilSet b v
     a)).toReal := by
 31    classical
 32    -- Choose the exceptional places for the full summable estimate.
 33    have hgood := GlobalTheory.T.rank_one_local_density_almost_all F E sig hsig M J
     hJ hM d b
 34    let B := (Filter.eventually_cofinite.1 hgood).toFinset
 35    have hBg (v : HeightOneSpectrum (O F)) (hv : v  notin  B) (s : Fv F v)
 36        (hs : s  in  v.adicCompletionIntegers F) (hs0 : s  !=  0) :
 37        |(m v (hermScal (Fv F v) sig s) (latV b v)).toReal - 1|  <=
 38          2 / (Ideal.absNorm v.asIdeal : Real)^2 := by
 39      simp only [B,Set.Finite.mem_toFinset,Set.mem_ofPred_eq,not_not] at hv
\end{verbatim}}
The local condition is stated using ideal powers, so the argument does not silently identify a rational-prime exponent with the exponent of a place above that prime in a ramified number field. Caps apply to the added multiplier. The baseline may already contribute its own exceptional-place factors.
\clearpage\phantomsection\pdfbookmark[2]{Reusing the infinite product estimate}{plate-CappedGrowth89}\subsection*{Reusing the infinite product estimate}\label{extract-CappedGrowth89}
The capped proof replaces the fixed baseline at each exceptional place by a baseline discounted by its permissible valuation loss. Each discounted factor remains positive. The good-place comparison and the absolute convergence for each actual row supply the same hypotheses to the common product lemma as before.
\smallskip{\small\sffamily\color{accent}Proof extract \quad CappedGrowth.lean \quad lines 89--112}\par
{\fontsize{9.2}{11.2}\selectfont\begin{verbatim}
 89    let u (j : Nat) (v : HeightOneSpectrum (O F)) := (mL sig b m ((A j : F)^2*beta)
     v).toReal
 90    let base (v : HeightOneSpectrum (O F)) :=
 91      (mL sig b m ((N : F)^2*beta) v).toReal / (Ideal.absNorm v.asIdeal : Real)^(H
     v*(2*M-2))
 92    have hbase (v : HeightOneSpectrum (O F)) (hv : v in B) : 0 < base v := by
 93      have hq : (0:Real) < Ideal.absNorm v.asIdeal := by exact_mod_cast (by have :=
     two_le_norm v; omega : 0 < Ideal.absNorm v.asIdeal)
 94      exact div_pos (hpos v hv) (pow_pos hq _)
 95    have hfixed (j : Nat) (v : HeightOneSpectrum (O F)) (hv : v in B) : base v  <=  u
     j v := by
 96      obtain <e,he,h1,h2> := (hf j).2 v hv
 97      have h := ExceptionalDilation.lower_of_bounded_order (b := b) hsig hJ hM hlam
     hm v
 98        ((N:F)^2*beta) (hbetaF N hN) (f j) e (H v) (hf j).1 he h1 h2
 99      have eqn : (A j:F)^2*beta = (f j:F)^2*((N:F)^2*beta) := by dsimp [A];
     push_cast; ring
100      simpa only [base,u,eqn] using h
101    have hsafe (j : Nat) (v : HeightOneSpectrum (O F)) (hv : v  notin  B) :
102        |u j v - 1|  <=  2 / (Ideal.absNorm v.asIdeal : Real)^2 :=
103      hBg v hv _ (hint (A j) (hDN.mul_right _) v) ((map_ne_zero _).2 (hbetaF (A j)
     (hA j)))
104    obtain <c,hc,hlow> := ProgressionGrowth.lower_product B u base hfixed hbase hsafe
105      (fun j => (hfin (A j) (hA j)).2)
106    refine <c,hc,fun t ht hcap => ?_>
107    have htgood : good t := <ht,hcap>
108    have hft : f t = t := if_pos htgood
109    let a := A t
110    have ha : 0 < a := hA t
111    have eqn : N*t=a := by dsimp [a,A]; rw [hft]
112    change c * ((N*t:Nat):Real)^_  <=  Prod' v, (m v _ (dilSet b v (N*t))).toReal
\end{verbatim}}
The auxiliary indexing function sends inadmissible multipliers to an admissible baseline. For any admissible multiplier it is the identity, as proved when specializing the bound. Thus the reasoning over a total family of rows does not enlarge the conclusion to include forbidden multipliers.
\clearpage\phantomsection\pdfbookmark[2]{A neighboring rank with an exact obstruction}{plate-FourSquareControl234}\subsection*{A neighboring rank with an exact obstruction}\label{extract-FourSquareControl234}
When all four entries are even the square target admits an exact division by four. Starting with the verified base for target four, repeated division yields the scaled count for $4^k$ at every sufficiently deep power of two. The bound on depth moves with $k$, which matters when considering a sequence of different targets.
\smallskip{\small\sffamily\color{accent}Proof extract \quad FourSquareControl.lean \quad lines 234--259}\par
{\fontsize{9.2}{11.2}\selectfont\begin{verbatim}
234  theorem target_powers (k j : Nat) :
235      4^(k+1)*count 4 ((4^(k+1):Nat):ZMod (2^(2*k+3+j)))=3*(2^(2*k+3+j))^3 := by
236    induction k with
237    | zero =>
238      have e : 2*0+3+j=j+3 := by omega
239      rw [e]; exact base_all j
240    | succ k ih =>
241      let m := 2^(2*k+3+j)
242      let t := 4^(k+1)
243      have hm : 2  divides  m := by dsimp [m]; exact dvd_pow_self _ (by omega)
244      have ht : 2  divides  t := by dsimp [t]; exact dvd_trans (by decide : 2 divides
     4) (dvd_pow_self _ (by omega))
245      have em : 2^(2*(k+1)+3+j)=4*m := by dsimp [m]; rw [show 2*(k+1)+3+j=2+(2*k+3+j)
     by omega,pow_add]; rfl
246      have et : 4^(k+1+1)=4*t := by dsimp [t]; ring
247      rw [em,et,even_scaling m t hm ht]
248      change t*count 4 ((t:Nat):ZMod m) = 3*m^3 at ih
249      linear_combination (norm := ring_nf) 64*ih
250
251  /-- Even along powers of the fixed scalar multiplier two, the family of
252  positive stabilized normalized counts has zero as its infimum. -/
253  theorem no_uniform_lower (c : Real) (hc : 0 < c) :
254      exists  k : Nat, forall  j : Nat,
255        ((count 4 ((4^(k+1):Nat):ZMod (2^(2*k+3+j)))):Real) /
256          ((2:Real)^(2*k+3+j))^3 < c := by
257    obtain <k,hk> := exists_nat_gt (3/c)
258    refine <k,fun j => ?_>
259    have eqn := target_powers k j
\end{verbatim}}
For each fixed $k$ the stabilized density is already visible at finite depth. Letting $k$ grow drives this density to zero. This is a mathematical control on what scaling and monotonicity imply alone. It is outside the rank condition of the original growth theorem, and the finite-count result introduces no new claim about disintegration on a singular fiber.
\clearpage\phantomsection\pdfbookmark[2]{Checking the route taken by the proof term}{plate-EndpointAudit42}\subsection*{Checking the route taken by the proof term}\label{extract-EndpointAudit42}
A module import can make thousands of declarations available. It cannot distinguish between a theorem being imported and its proof actually being used. This driver compares the endpoint types by definitional equality, traverses references in the elaborated terms, and checks a different set of required and excluded declarations for each route.
\smallskip{\small\sffamily\color{accent}Proof extract \quad EndpointAudit.lean \quad lines 42--68}\par
{\fontsize{9.2}{11.2}\selectfont\begin{verbatim}
 42  private def checkMembership (refs : Std.HashSet Name) (yes no : Array Name) : CoreM
     Unit := do
 43    for name in yes do
 44      unless refs.contains name do throwError "Expected reachable declaration {name}"
 45    for name in no do
 46      if refs.contains name then throwError "Unexpected reachable declaration {name}"
 47
 48  run_cmd liftTermElabM do
 49    let original := `CompanionTheory.PlacePermOperatorBicommutantSpan.sw_one
 50    let pinned := `GrowthConsumer.sw_one_via_pinned_growth
 51    let alternative := `ProgressionEndpoint.sw_one_via_progression
 52    let originalType := ( <-  getConstInfo original).type
 53    for name in #[pinned, alternative] do
 54      let otherType := ( <-  getConstInfo name).type
 55      unless  <-  isDefEq otherType originalType do throwError "Endpoint type differs
     for {name}"
 56      logInfo m!"AUDIT Elaborated type matches pinned sw_one: {name}"
 57    let growth := `GlobalTheory.CompanionTheory.SW.Growth.growth_main
 58    let progression := `ProgressionGrowth.progression_growth
 59    let shared := #[`GrowthConsumer.sw_one_of_growth,
     `GrowthConsumer.original_scalar_comparison,
 60                    `GrowthConsumer.forces_iff_relativeGrowth]
 61    let alternateRefs  <-  projectClosure alternative
 62    checkMembership alternateRefs (shared ++ #[`ProgressionEndpoint.provider,
     progression])
 63      #[original, growth,
     `GlobalTheory.T.exists_frequently_mul_pow_le_tprod_rank_one_local_density]
 64    let controlRefs  <-  projectClosure pinned
 65    checkMembership controlRefs (shared ++ #[growth]) #[original, progression]
 66    logInfo m!"AUDIT Progression proof traversal passed: {alternateRefs.size}
     project/study constants"
 67    logInfo m!"AUDIT Pinned growth control traversal passed: {controlRefs.size}
     project/study constants"
 68
\end{verbatim}}
The original normalization is imported here only for the after-the-fact comparison of its type. Its proof does not occur in either tested endpoint's dependency path. The progression route contains the progression result and omits the original sparse growth theorem. The control contains the original growth theorem and omits the new progression. This paired test guards against both a renamed conclusion and an accidentally circular proof.
\clearpage\section{Reproduction, provenance, and collaboration}\label{collab}
The archive supplies the original-source anchors, the four inherited finite-count modules, the new proofs, two independent arithmetic controls and the verification receipts. The dependency inventory identifies 1,863 input modules by their source Git blobs and compiled object hashes. The original normalization itself is recompiled separately for the endpoint comparison. The archive instructions give the exact commands for source checks, restoration of the selected cache, isolated compilation, kernel replay, and comparison of proof dependencies.

The run used Lean 4.33.1 and Mathlib at \code{0df444a360eaa60ab8c11dca51a86af692955474}. The original project is pinned at \code{40e56bab026552e089d7afcfeff08d2defb9c005}. The dependency cache comes from the public preservation project. Of its selected compiled objects, 1,017 also match the independent MIT manifest byte for byte. The remaining objects have a distinct byte hash in that manifest. We used their recorded preservation-cache hashes, checked all selected source identities against the original pin, and explicitly replayed project declarations against the pinned Mathlib base. Lean's replay routine accepted 43,467 project declarations over this base. It skips classes marked unsafe or partial, accounting for 31 other records. The merged inventory contained 390 shared declaration-name occurrences. The cause of the binary differences remains undetermined.

\textbf{Verification boundary.} This run concerns the selected rank-one arithmetic and normalization argument. It uses the built Mathlib library. The four-square control is an all-depth finite count. The rank-three comparison at targets three and twelve is an independent finite control. An effective exceptional set for an arbitrary original input, a cap-independent lower constant, and the full global RH build lie outside the verified endpoints of this study.

\textbf{Human and AI collaboration.} Justin Candler directed the research, selected the progression question, and requested an independent parallel investigation. Uriel coordinated three AI research workers on the progression, the normalization, and exceptional-place dilations, integrated their Lean proofs, ran the cross-checks, and prepared the manuscript. The original formalization, the verification archives, and Mathlib are prior work. The parallel workers share a model environment. Their agreement is not independent human peer review.

\href{https://github.com/nousentllc/riemann-hypothesis/tree/40e56bab026552e089d7afcfeff08d2defb9c005/src}{Pinned original source}\quad
\href{https://github.com/homeowmorphism/riemann-hypothesis-original/blob/main/verification/CACHE.md}{Preservation cache}\quad
\href{https://math.mit.edu/~drew/grh/}{MIT verification archive}\quad
\href{https://github.com/nousentllc/riemann-hypothesis/tree/b4d12e0825d2191f38fdfb5e97b327c8c8ad1b3d/research/uriel/experiments/gaussian-target-four}{Previous Gaussian-fiber study}.
\end{document}

BEGIN FIGURE nouslogo
89504e470d0a1a0a0000000d4948445200000296000004960806000001e725dca6000000097048597300002e2300002e230178a53f7600
00200049444154789cecbd3f8c16c913f737fdd3459665c0720e64ce00a796059b387002a44e0ed23739909c1fc889250707927320b0e4
c0d241e0c801acf4869660035b7a238ed8b2616559b21cb8ad5e6a96d9e93fd37faabbab67be1f6975dcccf3cc9f7e7abe53dd555da5b4
d6532a4aa9575aeb67c95f148c52eaadd6fa49c9156635e6e59795325f3ed35adfb5760e8069c0699a7ed75a2b8eab2d6acccb83fc6cd4
0bb82eac1635aff55fd6960ccc45cd17662e76f1f7b647832d3192b4bca66975bdace7e2e899d64117bfbe8313adf5277b33cb798ddc7c
b67610b59f9a2a8d79e504e186f5f1729aa675833f98a6e94fcfe7bdb4949daa8da994faa1b5bebef8ffbabf1cb16e4073de168d5aad31
636f80a1816f6badffb1b63ace53bb41595e40252c5e064fe77fbb5e10ae7df4efcd86243e585b98a9f602cae905ebeff9cc18aee373d3
bd67ee89e2c6544a55317346a4a831e931bcef7a89d436d895520fac8d69dfbf62c8b35c13c3d8fc0297d6a5e893eb3b358fef3b76096c
c3496b6382d993db3b0a8f7fdbda524876632aa51e591b09572ff291d343388e9f60524593fd98fb2ed2f79908ee69adbff8beef3a47ca
f137be7faab52ed2e089a3315d17e9fa5c08df31b87eb0c077cd8f7727f49914aadb997491f7ac1d3f79577a13f3e8c9daf193e7a1e373
4f6a57ef99b9b87a5c953948c6fbc8ea990b1bef86b5f3c0e43ee61fa79fbfe60f6b0f03ae5e3905b697b0b03b8b0719b9dec9e68ff88c
e447bde40574666d6120d49053c309e61cb27b668f5e3953ab7776eb99dc9304335b3754f3312f25ab31d7ae5deb03e5bcf61ce1b1b5a5
80458738e1f891588310387b8deb47aa70fcf3a5c3aff8989c6e0b6e2dad314d56d312611d4e9a0b3406bd52eabdb55300f38f5dede5d9
dbd51b82ab677244b8459da74144c7f5dc911247632aa56ed598bb7451bd314b1829ba6e82ab9717342623684c46d0988ca03119416332
52ad319552ff9bb571e7541d014de52317d899ae199f23c0de98eb863c52c3e205c4086b63fa7aa16ffbde606b4ca5d40b6be3c1607b9b
d7f02a1ef26d2efd316eb586937d556f08f4cc0d527ae5003dd838b1aeaffe5fc7ae28295d6dd16509b4ef47312e0a6be3f677d6ab2ebe
98359fabcf3b57945898c73cf7cf7c3df52fe55ccb63bbb6fb3eef3b966bbbeb5c8bbff7abff7fe2faeee5315c1b536f34f1efbdeb78ae
bf8cc6fc14d8676d731dd3f5b9d0bee51fcb028154625f26be17906ffb149860d9723dc71cd3b56f496ee470ef17c966cc912b10a27684
5d51b06b2ea53d73b5efdc9720c0d59b7de70eedf71d734d72cf1468de5c5bfdff37eb138d486acca50d5602c74a605f0fd15a7bcda3da
243de69cbdd2d718bef3b93eef7b34e947ff2ef63137f981ac8d9df1dd586a6cd3b2b11c93db3f56ffef4ddb16dd33b9b5d2d710be7372
45d4397ab1f3be283cd258040f5dfbac2fc436a652ca4491ddb47614e0bb20df8dd66a4c4e621bb3f94280f5795b34a6ef3e43df59b2a9
99be13ec8df50fe77a791537e6c11af27cb5c03fa941838d7934fff7d6ca8b2d8b26a899351b53e2dbdc778e623b93dee020016f63729b
423d692557cec63c980ffc7cfe476874138353335bfc9221ed715d075334dd07adb5954228770e608db36736e0a4c7495d0d49bc9bffe1
189bcfffef4bdaf2ebb3eb9e29a557aeafa56051d545faf0ad639067f3abb563e37b578e7180c64c3ec6fcced05a27bd3bae3466abb79e
c4c6a4746c7ffb3e13b316f4376bcb8e588e587c0d11ea40a17d2e2e7ba6e45e99f23ddf3142df5fce5bba3e27fd6d9ecae69b7483d3f0
ee2938e68e6588c65c672f8c61d52b8b3310c670d198d21ef151f997c4a163e98fdb2baef35fa9b6542eb9375512201b48055905719a59
d22bd741125aebd8f851af6b38e57afe35b57914b29225675cd7f79cef2e5f701b3ef4e08bac899dd9c2b62cd549dffd93ff3ceabafeb5
f58152121af2caa3d6b2217ddf736da35a454e7e5b7fd9f70bd5842a4cada3d936e11825c57cdfb77d8df502e2eca19113c04fd6a5ba52
278e63bf531b6b0ace77b1a9e434484ca338bef38d338cb0e447f2ce1a153ef2c1acafbe8958df85fbaec3f7f95c5ce7498951f2f6ccad
1b09e13b31931edfa891387a716d57ca88a534a6a5996be820efac1d01d4d5da93dca562bed3f1aa24ca2bc91b971a39fc23e7ad5b1b96
acac577f6ca717738ba4e1a489c5a9997fd2d870da2e3cb7e9c964eef9868739c7ac9e3dc647a11ffcfa72e8b886c35f9473cc1a094fd6
8b3897dce0e8d9e605143a4ee0fc4142c78ca1468e0e4b53178f2dfb5bd8d700f4985aabd47c78825dcf16c7db9caae45abcef7cec5c37
591b57af8c1d40b83eb738de66c1e6e2c6745dfcd47978e7baa6d891952f7eb3ba66fa2ebaf73899ce7fa5bc8eeb5a5d786ce3283b9b75
89b44487d9ea3a838faa6b989b3415c8b5aa57aae771ade735af336755af15102ad985bbb6205c1d818be49e29711e318616d79dba447a
c8869c1a5d6bcaaade611b72c165cc518dc73d7b55efa8a12e35a33da27ae65e1a92b8f4e173f7cee455bd7b08beaa753f875fd5cb9577
644af501ed2924b0c67d61893423c1c63cda12e952828db964cf51bf5cba79e4c77ce9a8b326b67308be80f6faf299e1be3fbc8018f136
a6526ab95cc41b9308166de67bccf7fe884f78ccf9a0005bde631eb567d698bcc10be8275b6b2ba30835e6615e3a5c6b2bbd8ff974f551
38d35ab36b4c4f6ac8d8a18d766e428f3948048dc9c856635ea6e9e69c91de2b41cd9ca09b49e0316724a6312fed4dccbc8739a4abb716
7ba86d2186a8c6446f8c23a567423b3728291472be953dfa681405bbe2f1bf4a8e9df97afe071ef7abb02c10400ffd49d608c8b7f0e8e8
640f275d0d5ac34935125596fb1df5b12f9ee870351cf5d22e757c7bc2592cde9b72c2d5e07b843d1382ebb13f4aa3b2cf670656f53eb7
b6ec8c6a393a6652d66b8f4ef5c63c12d98fb971b085f47104c8ea60cb94cd61672e6b52de2ec958559bda6bcf394d239361e0fe729b88
348cf6d3536deab08a663a6ee0929a0d1caadc2a6a89740a0bf3e8c3fa6bae847babbf4f2669deac65e6bf26fb026d777d7e999c64dd90
ff6fc05463c79b3f93891fcb1b09f5d805f767b9502aab0dd6291b6fb5d2f1aab9e0227209593a9bc8694c6ca519eab670b154d3cc9ccc
ffae69bdf9dfeb1ece717c6e101ec3488d2c85de244c1b9cc6686a61b986aa838ca2c7dcccacaf6bf5e48ec55ddf733de6aecfe51e9f9b
d2c6b42ed0b5adf4585385c6acd1b8256373dfaf708ff6b394d272f0928e9f33a60ea63c2f8525b19eef2d9c50d7c2f979df397c9ff7e1
fabcefd82570bc80bcc5e26864e2cde64a239ba2174aa874b8f143b59cd92aee99be5f35f5265cc709f51ecee3bbf6e550cdce4cb9c09c
9ba97dfc1cb2c6e68b5ef1cddae9b8095f2f0adda4ef3baeeffb3e1b3a7e0d8a628d2a4fa7358967e2bc1791c349574f736d9306c6e68c
b0056e7111ea81a17d055813d8b9e4440e5f84c1d4684c5ff2fc2595cecb323d971d865deba6ac8d0ed8bd8a4cf794fc986f992325c4dc
4cad86643916f7ac11cb45796e506a8f9c297a9bd345444dea661c77735b2eb33793db73c91984a0b91dfca1b179c9316b59226caedea5
960aad716166973ed5bc3676bff9a2511f69ada38b1bd5a4d50f5c73d6a87b439a1f746a38e1213a3e73b4a585189b3382c664048dc908
1a9311342623684c46aa35e611974d5769cca3ae3faffa981fad516bc4671e76c95b8d2ad29bdbf60adee68cb03566a80786f6ed0996c6
44f6989f347bcce7b9c53d53ad58bc8bdc38f4a9703ed30437ac0b20d7a0b458fc2889f582512213c5e02f7f3cb34c505d2d18bf8de999
b97fe6eb897f8f52ceb53cf66afb5bdfe753b6bbeec1b58dfe5eb9bebffc63adbc1ff9e3c557b2f73ce63e0759aebbd9732fdfd6ab845d
e75c527b556f357c0deacbad14f8bc856f55c616599a29dd142ab93e47839f581ff2d07c046496ab581bf38f25ea47ed910dfb2f6b0b13
8e5ed59421c7e6318da6947a616dac4c52630e326c9c97d3fcb9dea194aa1a6512dd98d28683bedea9b5be656dfcc5436b0b23293df36f
6b4b67020daa626bc0ade5c0b1d6f3e3625ff8c97459f28e51c227c78820fbcf750ed79f6f0494f3e73b8ee71e9e05f659c798ff920b85
30716f9d41c1856f0494438b55759b8fb9a3db7320723ace9578ca4c78581ff410a399cee159216c990139301d66d97357bdf76bec9319
6ccc8aa650f4afcd4184cd79d161968de858dfbe2d4b21cdacd898b199b2b82687bdc709ed4bbd066fcfaca49533b9b98f4a61a9c9ebc3
db336b8e76521395a47cc745e84d5ebd671664cdda23d1bdd9d9336b8fc10fd33339131a4b20255955a983d0ea992d66865af6cc9863f8
3eb32c31e1fbee12b13ea0964f8869a8b9415d9d29f687bcf298bb0ed4918f9ca7de6a9079e433cf38adfe3f0a4b336b93f9b8464da7ad
49e91c6a5176476bbd9e967b1f95db73d6cc56bd52925ebaee39e4e6751d63c921e2337d8db0f51887f6b9b868ccdabe9105cdcad8b87a
5df5735ecc100b7ec453bee73bc6d6f7439f0ded5b23fe31afdd909cfc4b5aaf2c8533622415454ea2eab47ac4737a65e83ba17d6bfed5
a8c7dcb6b644d02bd23897269a195ba8a334fff0e27fcfad0f941f73938bc6acf92be61e3b372896be9b3afb73a510be2ffc9ab63fb1be
4d2c7be6636b6f435c171f83a32173f2175b0eb79c6c5c57a6e05aa4217391db20ad52e9c6724533992f266aba3fa7411c8fe1a984543d
d67ce6726eaf8454572ef1cefa50f8f3ecbdb1a4b73bdfe6a51718337c7334cc3badb525ee731500d74dd6c85ce8da6e7dd8872b9a6bb9
dec6150516fa731d6723e28cedf8dc51721413151fb5e7da687dc87153bdff5cd759d0902f5c8db9b8f717aeefadff9c8ff91aee478a83
853dc85194fe324ec0232796e9e4c4d5c211bfe4930a3dd378026f39ce151b68fbc575ada98f7989ac58ae5ea65ef360cb21c61588d5e2
b8d1762f7763fa2e28e5a238ce57e0377ae48adf8fb27fb91a73e9b05f51add0f0eafcd68d248ea46eacd7a4a7ce4271a695b01a92c4bb
c99a74bad9d78eeb8ac5b5263d2904b14a2684de6f7ed71bd9fa90e3f32513c35369634a6cc899940675dd47ec779770964a7c27c9160d
f532d7673dd77ec3b3dd4951a52adf854ba2e5b5162f91969ea5dad143abbd104b57f57a6b4e4a62d5a0aeb7360b29ab7aadf9c998257b
1271740a16a23573b4c7db456d577054cfdc43431297c94b6af4ce1cd3283a9b8a34b4d65597e4c4acea5df7cad1d7085dfa99b87be7a6
66f60e39a941ad7b0af6cc9895ac20b231a769bab3f8f7536befb82c1f7536d90a3ee67b7cc4676adc1b123833e26dcc75f889f50160e1
7dccf7fc884f78cce583c664048dc9081a9311342623684c460ed998b5b2be1ed2cea48ad29739d8b9eeefb0467b0d428ff9376b0b08e2
ed99137a6732a19e0912416332126ccc50261560136c4c90464c635eaedd46ef0c137c9bcfe0ad1e47ec638ede194156e0167aa79be817
10deecdb70c6b41f9ea4c6748434a3411724f74c34a89facc71c0deaa66404746571001ab4a0316971001a7401cbaa5e57231ed116e55c
22ed3cd0911a9575f1beaf41a783342aeb145c6041e721a8329fb968d4d7d6ce1d5325e1490892820f5a6b9145e98e08a533fa5d92aa34
ef98574efeb3bacc1f8b4d4d72a31c155a89b45c40f36da31a6f37ba76cc35eb08a01551f5008f0ed582f096771ec5d612e5e7354fefaa
08cf7249dc7d65d77d5dfefda8992f420a46f51cf7becc38bdec94af97ed39d400c095694ff21f6376c4b7ae8c882dffa669320fd22bc7
b565ddcf68bf65e84fd4ab3c1552c8ef946f7a6dafee952b03475a90fc4a6bddaa9c531386ec98f4ca3ad35adfb5767ad8b05fa5713bb6
12c8a24d4c5b7cdecb7cdd50b184264395e994642f4577cac961bffafe165f39d9cabe1f8159537eb29e8bf49d7bf197d429e9985fe64a
05f4100ecd308a392b5e0b45f02525757d668def3b31c7e4a2e5b96a319262367f0d732594514a75992bec5993af942114b375786dccf9
521453e2f54ba79b623ae6e0926c4650d4f63f7c458fa46095fa6a81a7c4e0e7e96755d620f3e027f499126aff58bdaf7fbdbff6f5e4d2
e555ee6abcd82a09b19fe7b8aed0395cf710fa8e94ebf75c77d2d45b0b8619fcb822e73d8d9c8ce7582f190eddf3fa9de5a55d9d7515d8
2182e61dd3d180bec6b2707dcee12b8eb986b7a1ef24559d4b20e2fa7fc41c6ded2ff75c7f682ef3b9b54518bde23197984083e4690dc7
7138b8175b52c2777e57e773e1fb7e2156a9321faef34bb2359b2aa6a7b648d65cdbc29b525a86fddbc2e312db298be73717d75f6a32bc
5b5c7f54a79c0698466aaa982dcb8bd0e4b20976b86f8cfb699ade73bc9e0305202f60cb22f733d5df03ba7ee3dafcc46d5e385453ccca
82dd76cc1a387e481fc9411852ee49ca6fd2ec557ea04e69f85a2bb969054a4da12a3453ccd50ffb546bcd51cfbd090503952196834814
8d261d938218beceff3f925a728c9e47b85f69fef556aff2c3764acee3b4825603f4bd86468ab93ec9302b1e1dd79ecc208a7965b6a1f7
3537514cc72abde58a47d1318325ab0ba5af4c344b7d17dea38b4e29e59a9bbb24579dd42c5ff86bd149452ebe9fd215e44c6a8734915d
aba5beef1cc2d11d69090f8608700dbdde4586900d38f89496f0e0ca93eb0b52e8cd2803b845fb5d1f2de981e8b0b755277d7ff4149c31
2ca39452630024316cc283de91d7be87a4c735498d422f61d81a35ab920562074db558ce35ee3129e9d029627a821cca7541552f201274
4c2012744c2012744c2012744c2012744c2092e13a26721d1d83613aa6c36ffe992a34801d325c1a42075d6a03f9ae0913ec3c88574c5f
0758702d36b50a1807b11dd3a4b64e8826bac695fd17c840e4abbc20bced71abb222be6bc4ab9c07711dd3f783c7d2aa63f8ae131d9307
31af72534ccaf763271e07e1523b4044c724fbf0bbb523ff78e89c83d3bd635227ba6fed283feeb0a54472d9d303d9b563566ec8bfac2d
1df114446027359adf513dc4fa5b7cf6966bffea8f657d51d7722ad6c601cf91c01bfa01a39c01ab6404b13c5f2ed38dc4e46aff10f8e8
c9e23a628e7d8765ac20a56a452d6a8d927df7e03b5f6a46b5dc0c6cf4bdac04acbe7bdaba06dff77c9f8fa175aaeb5bbe9ba8784e9176
57c4753d4dfcfc9287d6960222d6a39f5a5b0a6999b8f56dc66b66d7843a9b2b7f68e8f38ecf72a5ac76966559c19ea4b64965b4deaa25
61ddf55cb2d9da9148c2bdfc1d53692ee2babba4ecaeae989853bcc2ebf58650fb1cd98b54b563861abd3512922250e998736b470057e7
4c28b435440e5217d53aa640a51461df52ece83b6b470057e78ce463ddbba947958e89d77718adf593d4ce469f3fb376ec14f68e894e19
4f6a5a405321973b9560a826e57aff325fd2b29ea54f99536b7c5ef92ee704bbf44ec9fd835a1b050c5846afa734c3a698be1f4a12adfc
d57b62991a7bf567cdb372c2a29803bdbed9aa65f8ee7974c5a48a6e7f5a3b22e0bcf762c5f4fd40602c16199b979df2de3245f62a5db6
3510e34c4d5ee4f941a79445eeefe1fade96fa998198efbb1c9eb66cc5745d10188f5a6bac4a8f9bd53107ee94a354c4e5a2b4487f3792
073f232ba5efe9cec1d70ebd063f942ee74ab85bccb570dc87c901304dd3cd9263ac49524cdf4d0011b0c66026c21e8114ad987be8943b
574ceb7a12a3de2d4abedf64f0e33af1802445f58c84ebf74974753a3feb3a6ecc677cc74b61b363ba4e3c223d32c28dc4d6dc24fd3da0
bf2f9e39cba75c6f8ce03ca6e3c4e057820611787ea3e7d69608e6b9c9e9d75298df57df72056b3c752d0329c56b63fa465a83c29a6ccb
d3192e686963d2ba9ebf7b5e432d9c1d93821dde583b0685fb87f275ccd61dc2751d89d33c3fe602fa3e7cc7739d3bf73a5cf86cccdd74
cac8557ed15090437738cc2cb2bb4fac1d71df553527f02dc5dc9b5dd94a2d6b9ccb87ef1a4ace4feb832c1b3277929e55315d2718990a
9dd25b2903793979b9ec98be1b1e987b152efdb3b5a521bedf688fcb7c2f3ae61e53f669ad59b28ecdf83ac5d4a863f8cebfd7b5e7b362
8a4ad9574a85577828ddca0d6b0b33474cab3d77ccac9199442afd58d65c21616a0c552be56226f28f9aebfda2636aad3f8d1cbb375341
298379e16bba39039996efd5e894d40792985d93dcd732ada78b288dca9019d92a744a2bbeb1e6f916e7ed66cb86ce9d43d1f495c7f363
6f144ccbb9caa95207099db3b19bd37b1d6be6eb8af94eea3d383d3f03d92f1f98632c9f6d34f209f3f91e78a2742ee0ceba11836335e4
4bfa3b71ecbb729da1bfd4eb702ae6950f781aad37cc1de4eed61c25d7f928077bb0740c8a58791473498fa7760bc64ef2831ebc50a72c
56c9f93c742e67a77429d191895ef323a4d16e3074924f8b4e128aac39a37b4e1aadba52aaf8ce23b9333a52c26cfdf1c6a89a5779ce9f
f96acbbf82ebfcaf12aff381eb38ab63be777c6ff3cf752c697fd334bd75dcc7dd9836e0bc17f1b98b1216441917e41d6b8700467a3d17
967d39c9990f75115c5a11cbf2e2b93be91e824b02f770a3a6e7281557ace956ba172a7a603aa3990337a1732c1db37a01aabd45c357e6
bcf7a2b9c043b4a47a5df8aef5ca7dc1a90d314b7a6fd5522d06f382edd59842c65b8afd81125748df47c92b3df42aea41a6eb37ab0c5f
09f48a76f9eb7df0e51f95dc31733ae3a8f38089ab52bbd9a65ba619db1cb3b48e99da19f73a211ddb0eb5d2de1424e47aa9b5b60651a9
742fa43f2dc2a73c37bae6f111bc24b1f7c899c57775dc9857b2957687a3534e02063fb127af3e0a1c8198198ec29592aedfe39bd6fad6
ea734e7f3fab7a4bae570ebfb19f8d36b43a13d3717db0dbbc4d3b66e44ddfe35e48b667426d3af283ddc4c65c44f17859d854e8940950
e77326d1aa657fb6a085e7277802bcaef908b5f568ed5cb5fa6ea8a14c055a744a5eb6d4d3da28982a8a0995ec4fe037e8e2e64ca575f5
dd2a4b4f814da09d3f8e5053934d315bb9aa401a01a1103dfbc115288c4e29185fe794fcbb14774c9f1760069d5206a3754e0e1b139d72
0c9c99957d1db637458a19ba29744a79f87e2f89bf5549f55d6f50053aa54c7cbf8bafc3f6245b310337735b6bcd5e5b10f0e0fbdd7c9d
b617b965a1bd13b4e894e239755d60e837ed4196621e3599e85e18e1f71311c10ec09ae48ee95a140f0037c9af72bcc6c707af72003241
c73c18a1f96749e0557e307cbfdfb48357f999b505ec01a72fbd17391dd3b9107ea37a1810c0865a8a728c2477ccc0fa615ff530208455
768f6f8bab72ae13ea093c3f402459a3725f07544abdb5360290017b7491afd3029042f63ca6af03fa3a2c004d3a66088cd041291c8bd1
9c07f0292a0031142ba6af03fa3a2c004d3ae684ce092ac06663a273024e58073fa1ce296d4d09904df36c6fbece0bc0922ad345d4f9ac
8a06d32ff5bc6bed0060d94f7a6614867a021f4d8a03843ae7840e0a1c88aa5a810e0a669aaef9a18e77c3da418c5c6501f0d27c319a09
34a60e7acfda49cc1d149df4b8742f72ba95f875095ef5c7a1fbf2dd594163abbd2efe5e591f00bb4164bdf214155dd1bcd83ca883e842
fa33a8617e3c86e8986b421d159d721f0c992266b90c75619f3eb53e0886a599621aff382aebcac6d8f681bc01cd68a696a6432aa59ed1
88fabaf501d08d39245142879c7ad9962bdb1045043ab0fc0d765110a094d55ce5d7c57c24a6772a61cca7b5172d76ceb8355d073b8e46
f97bd17043e46d948cc984b2e8849f17977a7bd7752339094df98c5253bb274aa95be6cd13b884214c25a91e9e988b3abc2d9ae009bb21
65101383c8b94ac76bddc5d7957f7cf7153316b317f32b39d8211773b9c374c849fa04faa273beb676dafcb9eaa443077298876c7d3fd3
34fd657dd0e654ea002696e1dc8d4cf1972621e90bad75d7948834db6014fe8eb533913db95987f481cf340a10363515e701d6e5406b1e
7429a59629bdcdbf8d9d675677deb78ec4cfbd3d7ac986ef94b3421c2882fdcaa065d9067b61d8ba3deb1f6315a0f1d2fac2b87c5805a0
5c19b4986d7b73dd0ed929b7d4416bfd62d549c525b20ff07ad50937bd5c748fdff7e2111b72a013ea9009c779324dd31b6b475b5e9a07
88eb8c64c20c3527e96228a5e4b41bcdc8db1197f978faa53c2f7d85e32339a563dc5e4dd19c2fcec93db7fa2e7359892886514acaf676
bfb651bf0c56b076fefa8ce94c7f5a3b02df8b392e07adce53939194b2c514cbb42aaa043a3044a75cb810abd79fd45adf9a984d05a5d4
eca36f312b707b1a7c8a6cb4958d4d268bb75e81a9afefd6afd4d15fe1430d74b0c6e718a0e87d009a36028de9da294d7479ca0f6f42b7
ac8d7561ed94adbd2ea33e54dd3a25d93d0fcd047682511e13ba55c42ac082fb78d5839257e17ab7ac0f5ccdd3e4dcdf9b2e9d52f8c8f0
e3e2dfc5cb2f564b38ae591fe0e78fd0f5afda3eb474a21b626ccad864fe2d5f49153c2e4d5fe1eb354dae72331205424ca74c185957f3
5737fa81aab90123aedfea941219364ac8da587ecc660aac94620f9830cb69ad8d2b6a287f0d9a774ad72b840209bcb8268195526c7396
f44a5d2b305bd22cc7f55f8be944b190e9f3fbeae3dedcf3d2e9a194960f5b6b9da35277381493a699ac576a83f53bbf732826256df8bc
de1e085fb37cfbd2a68e7eb3b608658eb05e5f5d497ca5eb7884b780402e9eebbf56e3fa43c733be7dc7f7cc5ba2eb22ba25126c4a6779
3d0fce8086d45c445b55292aba339d01258bebdf9c8130a6c6c6f5a7b4a7489a0764ac1b33f454bba0575eb5f9bec82201c9f1948bef56
6df0d8220739df6b4553a5e4485aa5b5be5e4b0d5afc3035cf91706cd183a0d6afef87d6960ca86372a6943e4fec2c450bb4e85cc11987
44ce52aedf350892b4e8ac77fdc6e224558e632691a35ca173669823de63c5c0394892f20aef9d9fb2384061b1282be9959e9b6f671145
ce42ee7530e40b125bbc6037f194e695bef8a1cc92800f8bdda794df5231fc9837ad2d0cacaeedde6a2525e7f55fd03b8f528896d521de
aebc0ecfb5d643654573bdf21c0cb3eeba7426a4162d3ba5c8064821d6fe1be5dee665cb8b4d22b225a35346923820192615b6c4dfa557
a7fc362f651d85bdd68a94d8297b952cd975879c2243c9809b269df2a0f5bad7a1642291a8e84d5edf23db938593dbe7e47d12cdea1ebb
67076ede29773ec0b118e17ea58906921104600a221e21a7cf89b5a523504a0f11d5bba2194d2d77af94aba24b5628be5438ab99a19474
1a2d5edf775a8c6400002000494441544b97dbcdc152d471bcd65e8f96eeb977e68c56af6fd7abd0fc58ad7303255318e92e7ef44dab42
af1497eafdfa6e1d4f692d231871741a8bf47b73dc97084f5bd3d1f75c4a6419fb488ba058934a71931bef686d14827166b8a681a478da
ba66f275358cf521413894c589f00e29becd7b479e5f09581da072d6e67a70a91dd2bc8d561df286d46b153179be88169fa87296c8113a
b9df9c6bb72543ed39a7387cea2aa727097189f8978321c1aae36d3449d74c6f9d3925cd30e182e2dc8c8bc1d0377a9d4b4cbeefccd421
099acabae8909206313188cd25b4ac67b3c7f2c13541c992cacc8321ea9c43e457ec855973b3281a3aec433c5ad6b5eb504d37d419cd20
46f49c6f0cc374cae957ba11458af91ec59e2e07860ff6f4a00ed52967e634c9265bc568eb7d38d9eb5b63e820dfa376c8398fe55ecd18
449e0fc8decd16744a200e744a200e744a200e744a200e744a200e744a200e744a200e744a200e744a200e744a200e744a200e744a208e
e13a259245ed9f1195d22a180ff6c5509d725e10155ae20ac667984ea9947ab6fa7f2c22db292329e55fabff77168107e33344a7f4bdae
b92bca021988ef941bafe92a1565415f4650cae06b9ad293801d21ba53fa5edb2b72533f03a188ed944aa927d6460fb02df78564a57c63
6df103db724788ec949955638f5894749788eb9405beed3fac2d6048242a65b66f1bc11afb4054a76498de41b0c60e90a69498de01723a
2557e48fd01ce92001119d5229f5c8da98cf1dc663810e4851cabfad2d20893d4d8975ef943502760f1a04bc9b29312c1cdb0f1ff67227
5d3b259635f0a1b57ec431c833b5d9e7da45bedf4729f57ef919fa634bf5dd2d117fed72caa327a93701295aebb7d68e3049833c5fa7f3
b1f1f9af4a29963cec3d95f2a3b565a7382ac8c6f0462995da29ab117bfd1c6fbf2e9df28081b9738995d41fecf7d4d762ca289c54edb1
b5e3eaf1ae67286a5128612fa56ce2b911344d727ffe47c60ff6d5da12266914aeb57e3f1768b576fec4e5ba3da53f1f45a184cd3b65e3
c18dc46992e41fac519b6d0555cf75c2cddf03faab62b30f59716c404e576a39d4202c74ad7331576b47014d95f2c053409fac2d8944b6
ddc52b957980147a4d57a159a74c5973c37cdeee23d8b996e4924a0fe87bfaefefd69e4c7a54c56da994296b6e3861fb817ab3d591b5d6
bbf07f37e994586de86643c54fac2d07a1955262b5a11baf8a6bad8bedd051a9de29e1dffe4968049b428267c5b26347e11051424aa951
55e7b9b5259e60ba1bc954ed948254f2beb5a50f96d724a4687b19b8a48278ca8668ad5d7eec2d453bb7b6ec9c6a9d12b6a497a44ea6b5
3edc5af62a9d522975d7da082ec8ec64afad2d3ba696527eb6b6802537525a436bfdccdab863d83b65c870ef09f332de22b4d63f5c839e
102347d1a7524329a54e4588529bc5a0e7a5b5d343ad8e49f6bf732500adbf71c62d98edbeb143688dcf16aca16b1b6eb3de489916ba24
a793a57cc72cc3087986123a8d599af16679eed4e51129d7cdad945eb719e8826b0a4a3c6c4a89a00b91043b65893950d3c6e5544a045d
0016583a25ca86c8a0c65a7a93e0c091788035f9c01a2ea5445e49193847c9a9cca36a1aa4f8121c7c2d19618728b629e14e1445f14033
b30882f9ce3dad354b6e500464ec9ba4c1a7a7433e5e2cad558179d5cf5c5382459dd27313400ed19dd2f55b52477cbfdaf62230f2fe9d
23ee014ab9634213e74b7c1dd2fa60dcfee2b887ec4ee9ba1100a6817309012178c4e5b6b5c541402ddbe712f2dc08e80867612bad7557
ef1c94723fbcef7427ceb42e2523f1e44e09ef8d5892a3a038522506d2ba64cf99e62825bc37fbc1d7a1ba92d4291109b43b44bef55295
1291403f11a9301938e7317d91e6ad88ee94c2a3ca0f8d2762c7e70ebcc495a290e0c890979da02b4529478f2a8f5e0b332096e2053a5c
134a127419dfa6b5d1fad04f39ef955f9285c0446f32a179da1eab0e5dd7137b1daeefa67cdf778c927688ed94c34f96a353fa717d3ff6
18aeef96b6c1e6ebdb63af8003d06b1c116353a6d67139324eef464d38923f0494cd84a279d7cbd750c9096e46767ae4c174257f085611
7311e84c7fd1b2871754ceef856f1944e01849046d4ad78907e54c6bcd96742bd02e27add342bbaea56890e1385e04e79cd9e10ed129b9
071fbe76197190e3c3755c1735eed9bb700ce169f2a9f91bad52b43c5878b13ed57e1b7895724f9d72af4ae9b98ec7eb7535a3e11ce8c0
a5e847524a4117a51dd2142df00d6442cc35cde9af68a9ad532977f6ea66558e50dbb4544adf75704c9a6f1d8b626a83218cbeefc66029
256758bd04467f9525f2a1d179821db214ab537a8a8e836d9ae525f779d9b4d649a605a959f275d3f74e6a55aeb05edf3b7b75b3ce9f4d
81f619edd51d73cc48dfb7f1f8fcc5792d5794d275612353a1434aae7af1d4dad2801a05a8ae28e50e3b6593a9a01ae7f2e1bb86d2f3bb
8e5b12fec6a294d2a73ac071b8544a576f1f991aca156823b63478217ce7afe556ecae94204c2884ab45873c12179d5229f5ffedec9edf
595bcab146982d71a9d1b4d3a24fbb544aad75cb25a267d61666f6e6d0d8e237df13087ea194f27a8538e33403381d1a7b2d8db73ba5ac
f4433db4b6342290bba9ba42f7e25f472a4499c3c6abb3aa6b91ceedf4333752e84d6a64d3b89812dad12bfc0655886523d436b51f68df
b91b4e776d4e7579bec73225d4c545c50d7787ec89efc7ae1504e181dd8518c39e26cf6f7367a00db5494d950cc52b562eaf6c113a9fef
3ba5d7b91ce80cad96153a6497cc6a14f5dfb4438670ad4230db421db218a394f39ff9df41ffae2fef83e32fd40edce75ab4bf193458e7
9b2adde3eadc8f1ce78cf97bebfa8ceb1cd1d7626d709c40fa9febc60a7fa0e03dbbbe53b943be757da7f57d3bfe5ed5682fd712db73df
eb432215c2d39c51dd339546be6f03a916cf1a7aa86ef826ead7d43425ac4e6902630feee509e54e62cf15141ad450e47cb3f9489abd50
1bbfff6920f9bef5d99c5436d67288cb1de10b1341cb20dea9fdf9d89772a44231b6170f45cb24aca14e692ee86f6b871c586318373a08
6b87244f4de835f9ae71508928bc9d72a2aaf88122e45d6995049560eb2411e7629f6f1d8d60a79c84bec61b774896f3c5bc791087f013
6ba0b3c63494b08e79c3da92498b0e19f1aa9e24d88f92880a5d13f4049f73f8b74d7c64ed0e39e7d6d9ea9054e85d6487344b4016f981
5c7f96b787e5bc5bafefd54586a62faa3342dec584b7cacbde654542a4be1d39852b29c8979ee81aeb5f62ce5dbcaeb9468734afe7a582
581fb07947eab89b0e99fb1d1f494ab9b880183b898d96e992b7ce45811a1fad1ddb344f3d9d83abbd5c6d5233e779d6720863d735b433
93cba9cd512cae868bf8aecb765a2a614a87fc40aaa806e990968de8fb9d7ddb597039c4139df85f5c4e78a6bf4fae733aaee199e0a091
17ae6b96f8e7b8f66030c53ab2c8f5999cbfacd7b70bac8a8c43f25ca4eb378c306774ec6763615bcd481774dbda01aec09582b915110b
c3ee7107aab029e59583fe0cff0a45db800552d4d3a594334df36fd6e894574e10b85160c1be1a3385c8dfaabaf7a97a3282c5e853355e
8957c2294de1a8d01fb93c5f32ded7f79e0f71a41a5e5b9820552299aa2b65888d88eb1654f3aa7084fe49a960b6057b9c69cf4e990a83
8a7473ed153c80afb5d6de348435c87427b3991e43744a5370689aa6fbd68e38c455e0524a9978c99bd68e0023a8a6b829a15a50c3a476
c86f0b5b505c1d1dadf5add42934b2e19a2ae6622c1045283b5d0a629532f31532847fd9458a2af59c426ab16c4464a7ccb01dbb4ea570
9132bf5b73115b64fd1ce76fc4715da25edf19e940eed12b661789adccda1cfa51373d24bde77f6baab5984e498d1c3b3a9da36f769900
dfacab8e55ab8dfc99295c3e08ae68a19688e894a9f6546a0dc251a18eb95504f47b854a683de78efb77ca840e79ef88abfdcc031871df
9f193ae69501e2d6efa29472e5ae6459d4d76da0438df8d9dae1004b4f7f12f100670ff80211f5d6317df3ac5cbf53974e493ed337d60e
9b0f477955c752734aa664f0d46de1180764446f76c823d98e29a404dd66f0b2c635a5d2b45392426e1ad1785d87a9d5314d5cc0c28bb3
b56af543aac7279666afef581b121d329ed659e25ad15229d121998950cc2133b73551ca5649a48e4aa87d476cd7ea4a196ab01974c832
42ed37e27294aa9d121db229f77c27e30a296b45b5d737cdf8ff61ed58800ec94b4804466aeb9a9d32786074c83a84da7d9436aff2fa0e
35cc840e599550db8e52cc9ebd53925f3484d7f601d5699629af841a4a6939ea179cef350652121b6a297eee92b55346bcb691d7bb1dcf
3d678a0984e90a5ba7ec513e0ef8d15abbe21d2f90ae96ad5201a6944e034c8416a249160916a5dc7af2d021fb306a9128aed7b7d74ec1
6bbb3b4f5d17b0f166eb4af1eb3b10467f013a657f7c1d50ea6fc3a194e890f271a62aecbd94d60787527a0f804e2907dfef24f1372a52
4adf8d4ee890a000f159d7001bdf5c07f2acdfee4af6eb7b63cd8db5561808f8b13d6f36696fb5924ee9fd225edd32f1fd66d27eaf1aaf
6fe76b028058d83ba5c9526b6d045238735d8734bb32b78aad3397cc8457b778467885e72aa5b343fa266901482157299d5f824aca67cf
4a094035923b2555d202a01ac9af6f9ffc4f787d0f41e0f71353ee85f3f57d626d01232126109bb308fd904595c025623a25dbeb1bafee
31f0fd7e93a0df10a36f200e744a200e744a200e744a30b3590fb215e8946046ccec093a259811d329312574301090014006e894076294
84fc787d1f08df6f37e1f50d4018744a30e556afad055edf0761a404aa50ca836012a82e4a218baed091d3299d09de47a9d1022e3ae897
45077d2cad493857339e69adef5a5b0148ed5f58620ba4019b1288039d128823b753deb0b6fc7cada3c41d28863d3f25ec4a500a5edf40
1c259dd2375f89573828a2a864095ee1a006787d03719476cadbd6968082021043b58a637885835c385edfcef5c2bece0ac0162c45e87d
1d106a0972a83ad0f175560042b0744a2822e0a4fa9410d412a4c299c9d7ab964a29315962817c58063a33e462bc63edd8e8b4005ce947
9c9d72da785da36382186ad8944e2fcff4ab46380041d89572825a8242aa8cbe431d2fd46101982a4f0979d713a3638210555edf331b9d
ef546b8da922605175f23cf41a9fa6e9be52ea89b5151c9eaa4a39b3a198f74c1a116b2b382cad22cf9f5a5b7ef1195345604913a59c7e
aae58f699aae593b7e01c50417345ba3a3b5decaca6614f3adb5151c8e664a39b3615f1ace233a30d831cd3be514394fb93172073ba6cb
12db980e8709f6e3d26ddd776cc7a4011238105d931144bea2af51e7849d7910ba67c8a08e796eedb0f98e57fa311091b68546dba109f6
4b4835d139774c97d177888c0ef7586b3d44cd411087b84e39958dbc6f9b7a31d656301422b3ae2d6abca4f2152b27c747742a40ea98ef
ac1d01b4d6622afc833c44bebe5d28a5cc6bf9a663d715e0091a9f6192a66aad6f51877b6ded04bb6218a574e11810bdd65a3f737c140c
c4d09d7289714722ba681feca6532e2197e47753c4749aa647982602201e8a6bfe7d9aa6134c1cfc629762b9c6917bebc3344d4fb4d608
42028765218a4b6ee0b97073886a3aa644fec25764f2773c9c638f567f5fb0be12ec0933ca524abd77f4754d42f96e7e36e80f42e9e110
966508a5941966dc0f7c64e6a5d6fa85b515000150da993711576282986f4114d339bc58ae89f5b7af78aeb57e656d0580b76f3e9aa6e9
5546ff4478270310cb0d32c5738989c47b81ce0ab620c7a419bdfcb1f1d12d208e15805826a2947ac5d099971867d32b781df70fcd873f
7338554a403ae246402c0b216be04ba1f5b98509817a8ba1be5c68cef049e4fc770908e7e904c4b2024aa95bd3347daa2ca03efe1f636d
4cd3f4dfe0a1ca8756cb3e2001ecf13b1a9e6aad91cb550810cb8678e2daa4714a426fbca55f46145c12ba5bf4f7a081b5570a164f0c00
c4b2334a293387f5d7a11be1389cd162882f34faf88a39c771384450ba348c938802838de5f66a1514ac5601f4df8ede5e03f2dcf57bd2
e2888bba3dc68aa46d0f4c4810d25acb0762d990c5ca89ebf3c3133abbc9cd39a7b47108e90da4b9e9c629395aacdf85fe921c71f43b9b
dff4d662854db06f80f66018de80656aa496c9e368a8f784c25542950ac12fba8772adc2d35045520810cb8aac1278885f2eb9f0008fe0
14496538d1e9f592056e2096955825da1c360464751f4582af9432dffdd3da112047241629fa6686cda4b36a7f58991dc19c6505561dfc
f5c8b1722bb1fa93ac65b1d05cdf52284f464e1ab16affcff4c2011d805832b34eddbe93b4ed278b7fdfa17851719045f979715de73b09
cc5f3af2fea4d542a031104b461c352e9e0f7d438443707e27e79134beafae47e23526e378e1c6a46203cc402c9970bded77be96fbabb5
a52326c1edfaec7bced98898ccf6402cf9d8fbdbfe74bd8186bd527828e85a5a8050b0c6402c2bb2b3c062572811ca687604a5f1db7238
b1a49a24563d12fa7b647da18cbd678c11f1b0fa44632f9ee38005efbc6f5087438925cdf3ac9d004bfe7638694ab8c378ac6e480f1772
38a06692623a05e3ebb3befbbec438e21c46c1f2cf9a6b076e8e665946cdf3f82c951c98c5b739f430f9445fbce5b683f6f78e4e022f89
f9bb4f221c716fa8681fd8e0306299d2213263f33e585b7e9d7bc80796a625bc8eab51621847f51c33e43f8db51a5df3d160c5912ccbd8
0e91151ba9b50ece77d2b0c737f7240e1a7aff1db82e51198f3696455e1bed85b5a8ebede3c4b3fd92943c9994571504388458a6cccb14
c6463eb6b65ce5bb742b6776800586de173802a525f02e740df4c2126d0d9bbe1a21ecdf12ac7aef88670512506f7014cbd23b945c71c3
da9280c94b68924d6c7ce39ac444af26cc891e529f33e1920d2bae1b5aeb2711c992ef53fb5b41ec3d5988e4665f35394ead8d1eb6463c
201ec459fee28c63c50765e58919ca5fab18b2148d116d7a483fc77c477aaa3012122b80dec1c345fb775b16b948081df3423fcf6cffad
11cf7c2d70f404d8bd58c67680adace529d0503ec54afd7bf1e0562d5a355b308b87347625c8d928391569ae6e734e6fc1d7459b54b538
4dece7aafd6331e9f1b2e6bc69c413031c3d018e6059c67480a7d696428c954ae2729678a49b8e58384d757ba22d20aaebf2657d9ccc65
9937385f262d30737a99e2fed0d1f69a442e5aac8cc3c4bcf81ced9f1cfb49a52a8ac2b462dba2e728473abb4efe4b3f7cc8a37b410b8b
69d0f021d664b3ad92ffba18b4fd599316c78622212bbb9bbd5b96228472fa651de4589a3db841d7bb9bacdc8bf63fb7760a6351f88cd5
09480e3090c9d11d3cb161156c50395425f0edfdbcd643baa0fbf24333ef476d7fdbdad9979346fde29eb565051c3d6e763b0ca789fa60
da2e49824541c12d63ddbea584a094923b0c6ef51b31ac9649e53425689c9398df0243719b3d5b9641a18c0da76885f1a03b6a833f8e88
1b8ce1b5a3b67513a15c04b98bc60c511dedff9469d8fedcd1fedd3206c50821d2bfd9ecd2b2a41ffaa3b56301de9cf5715459cca1a905
7c1456b5c99de019b9ca5ec5327853e804f5a13027aed213e7b93186c00f9e93348ee8e009ae1f06e59065cf59a367b844188310747285
d2c31d91dd59965bc30bbc2deb42494baad523c2efc70be528f0aee2427bff628f96a5572863c226403e14745eb5701b2c4c5eb6a637e0
e8f9c5aec472ab40d89e82aca54143b6267194104c764299b2828ed223b1ab6178e821c270a21e14c4dc3c09037e533ef0ec6c7314070f
9c3a95a08cea5db2d58c967d5e38de2c59e407383cbbb12c43491af066ac03a593bb29e0524e46a9072499d0ef8967685f62e9bb91db5a
ebaa39228fc89617b5034fb5d608752924f01cb166a01a915d0cc303791ecf2194fc0814ca894aba62b8588e2fcbffe12df75d5896110e
86c3bf15b9615acac80986e24cf8accba30fc577615952528250c283cf8b6cd5700830b0c804dfdb79764689292094059878caad521781
11dc21d8ebda703119d28f42e821ab057ebf3212d6ef7f4095c89d9795988918a623510303b42820aa4a6421ace5168e44c2f409da78c5
21e22ccd307d91a3d0953c60599616739b999879616a638e1c9c2ede55cee4be5b16436c9f50ae736ea28d571cc2b20c118acfc4f0a30c
c6a1392cff0c36a21690273491a3d7e031d6d08b5576eca5a3685916b5a814e911a1f60cad3b8ef97d1484329e65f9638750de6b9d297f
4f1cdeb20c11c8b88e00e844a4d7e0199940bda977a8e8c8c7e12dcb1073a1fe85d5794a1f7fb3b03883998ec04fa8fd5ea3397830239d
8505390be5f9a28cb18250f202b14c60e5289a130f2c633831b40903a7410126246e2190f33cfbd385385e8763a61ebfedf5c66a439df2
728848739a5f95badc84d00b508c231ccb04e16334d30158964c381c45f3443b04132431970f260bf233258399fb1684b21310cb4a186f
e33c345a2c2543520fe06521902f56318fe83702c030bc01b46e793964379ef45b3d0bed0319cc230f9a6f84e75f30b02c3b60bc94b350
d2900b614807c28c30a80ae62c9288231d00886567287b0f423c76ce7251034dd1e0053918104b001a601c8068e7b1815802004004104b
000088006209000011402c01002002882500004400b104008008209600001001c412000022805802004004104b00008800620900001140
2c0100200288650328a9eb7fbffb1b0560c7402c2b62eaa72c4ac0fe974aa9f32e170200280699d22ba194fabfa669faf75747ff0f8c78
22233600e301b1acc0c29a7402c104603c300c67642e7c1f73c4d8cf01006400cb92891cf1838509c038c0b264a0c44a244f390a560120
1c88650129c3ee0dbe433001900d86e19950bde76b8c87fcbeac2d0e009005c432835ace19cc610220170cc313504a3da8edc586971c00
99402c23514a7d9aa6e963a373413001100686e111f4102f0cc90190052ccb0d7a5a794aa97fac8d00802e402c3db4989f8ce0a652ea59
e76b00e0f04c104b374aa9f7ade62723f80b319800f40773962b843a5710837910cc8b516bfde3e8ed201158960b247ba1e1213f0c0f8e
de00528158d2db7c0431c2fce5fed15abfc78b512687174ba5d4131ae68ec05f835c27284429f5056d288b438b25059abfb176080656c7
21783c4dd31da5d4ada33784240eebe0195974cc705c6bfdcada0176010dc5cdad7cade9d83359b3569be6f9d2fbd687fd9c6badbdd11a
3472cb3548ee69adc558d8669988b571efecc13a1b71750f3d9c7f5a3b02a4de27fdb6efb4d64fac9d4c343c87e183d6fa91f501de7314
b1fe8d94526fa769fa9df152cfb4d677adad8d39dc307c2fc358aceeb159584abf53526576315b38d9aa9d83f840ff7d68ede1e38ce148
afe77f507b6866a19c684aa2fb737b18b134f33f3b9befbb696d016b0bec4d85dfbcc5390c97d32cb5faadb1d68c55b8fc9ba6e9b6f5c1
30af16225995decfef21c492ac81afd68ec181b3c7e28eb5e597c5c315bfe89ccfa373b00d97b5d69f56c76f320cd55a9b11cb536b879f
d073f5742dc685c2dc7544b57bb1248ff76e436e5a3d443be023f585524e03dfff9be91c2e3e3bb6d5225b905682f8d6fac0d5cffe43a2
f9d2dae9a7db886ad76249a51f9c96c08e68f91049e7c3c6f5dd67b04cb6c490e31c4e84bf186fe43a1db5d66bafbc48762b963444e5ac
912316589797c48453dd2cb4fe82d6d2e21c1c82b9b6b844be18c98a2c5dcf9e625d7661976279c0b93c58978e79be00f773978ed29c5e
0c46304b63616b0de939e11239f1f7ba3bb184d3e3f084e6149794cc636f0df767fe2849afe712ff8a73a2598c3284e6603762394a328c
5a602df14fb4d6d15eefdcfe921824ce9d7760ef73f062d98558d29cdd28c9306ae10c9b3928ef626fbb2093d36b6b8b07c7b2c2223047
dd87e1c592565060ce0e29dc2e495c8698351cd75aa7b475d212cf089017a003438b25ad411d2a6b506590c28d480963c9b5d412cfc199
d41743f10e0c2b9634d1cdbd0615ec8b1b9177936da92508e6611c217b6548b13c48b07916dcf3632343b17f318259da975a9ce30af89d
db339c581e29d83c13eef9b1a13182593b9d5d8b7338a8961e0eb8194a2c114309726921668d051359a71a338c584228e3419d71371949
1b72cff1dcda01866708b18450268310220f66c589c3028c5df5137b8e57748ef3c5e66fd607e5f231f6ca181743bcb7b678e895a64dbc
584228b3c0bce5068b9c8ae78e84be5ce7b8be38077b3df0d2118489519e13f72effac0f86b9937a0ccfe753fc10375dc7a81dac2f5a2c
2194a0362468a5197382d0396a5843a5e2502d465929e5b4142b3fd35517a788adee08a10460930785d97abe2d1c45f354c4a7d57f2f30
493da834efba3cef83c57fafcfcb6e03ebe74f16c3fce039a7693251065f1c01fd0f56ff9dc3b2a297b9e620522c219400d4476b9d5497
9cace3b5859c24d6944929296ac0917da94be62551c3f0a3670ee2c4f13606e3fd86496206ea22462c69b2fae899833881588e8fc8c073
f32276385772fefe1929cc4dc4309cdea0a12a71201da4f11a1f112fbc8a868c992ffdaed4e5a8fcdc38c3ac4f09a1bb6509a1ac0602d3
c7a76bfe03134349d362a942991b537a6d619dc6d43a6a4a57b1845002504415478749544322b99550fab9ab1eb8711c39b6ad83f4b7f8
9d44534ceece6e6209a104c04f6480356b291173cec800f11312c024215b05e9c7f28714a76f17b1845002b0c9a6107106d3537ed8cda0
6e12c9228b96e6254fac1d0124086673b18450021045b3f94a12ca98f345d71dda8204f7f6c6c7d6d7d955309b8a2584120059500dab28
614eac3b1473bc7f523334f514cc666249e10710ca768c50a01f3820016b45d71a56a9f39e53c705174dc41201e70024112360c5eba095
52bef5dbad499abf4c49e7c64975b18450760396e5bee188431491f734c361d4a5ac4c0bcb1242d981528f25e843ec1093e9f745d1bf04
aa8a259262801558afbe4d749672065833c41722e95a9c54134b082500d5e01296a472ba8c25244aa99ab7d24715b184507647fc5b1ad8
24d40267a9199e3194df5afed804ad75976c4cec6209a114010af08f4954ed24e6f9e8a7d69600159fefd8f9d36e9533cdd2256b63f6c1
20942270542f14414eff907a2fdca42cd8e06e9384153c33df52b3ac87a064197f043e3273aab5ee36efcd6659422801282276c1466a4c
e226244029533737b99e770ac042e52ae200002000494441548f11ca973d8572e24afe6b523a591b412fd8d6ef0a60a45adb4da8151266
8488c296526a861bc17caab5ce8af98c155c29a38b62cb920a9e7709120536dceb773b738858d1042badaae3ce083109d399b5d3cf1bca
3b1955ea7759abdcda69f354d2344c916549d98c6f5a3b00b0fb4ace1c97b86cd9dca4d4a069350cd55a5fe4d2a4e7fb77eb036ed886e6
261b51a53aeb45645b964aa967090d09dad0cd53184172b8c7415621c5ae704b4998cb8209d1596439af1ddb78b2c8aa2e4e28a75c6f78
eadc06688364cf313ce136e4dc88cafa23ad2d68a4f03673c9e4bb5eb19225248b25126388a56b58c516104b9b8436610dd50179e48825
428404225d5832fa8d68f12f852248a21ca3478935954e928307422916d121364aa99cfc837bf2ea5f81f248c6469054992ba4a9b4e5cb
68fe77f4b03a55c47b9c939368cb124229971d5a95bbb6a652daa3563b703ccf1962d9fc9c9c4479c36363a840173e486ef694d09805cd
3dbfad48148c1bd6163e92d68433f1b86be317b2695926c65a81c60c6055fe93118b7bb2c7b0a1c4b6e8e231ee61f54ab0b463085a9634
c700a1944b0feb2095e4450b3b15caa4051c1d436bd8d79e4750d3826623289688a5944dee9adc569040a4b2a7b5ed17503c65b4d17134
efb7d67a88dc12de61381c3ab219e1818263277d0147effb4fb95ece6b1d21a986d3b284508a47fcf03b335c685719de3356baddb3b600
31586299397402ed38973efc261e5a5b36d853107a86509a0c3b526adc000757c45229f55fc0a1231bad754e284e53324726bb99aba45a
3a2942f9729017e0a159afe0f99f8ede209219649e326be5cd5ef270669468c84e9e0bda72299698a714cf28f3597f595bb6b9ddfa226b
90f10ced329e74affccbccad4028c5f37284f9accc7e7426357f610a19f77e03423916ffa21f0c75a6e562b2ef882f6d9bfbc29db3728f
8a997648bd774a708bba558371310ca76245b02ee5f16d040f716eee80d1632a339e99f3111c74ad514a0d110570e90d47ce3c798c90f0
3573edf734f23ca549b1962194cff72694b43289e31877ac1d02b156f0c0c294c1209eef5ca11cd6037c845549ad125ba42440e6385f29
56503a2cccfe0c22943f3285f2dd884269868a1942f9e100a53192a7604cdfa1b64c2ea14dc1fe5db0c4728260f6e47ca035df39b5e287
2b546556b4d1fd260d15c989f3c8da31062999f76fce75c0d731b64aa9bb26407fb15f17f49d996e6d6a0dc3afeccc1f668174c417a52a
2c56f77204affe4c4ae5c5158fb5d639ebe245d1703aeee2059a71beb3699a5e4dd3f44fab10aca0584ebf5664e4041a8378c45b5c9418
2379bd37318c801488e4d9e861506b2a0be6bd65ecf00825275083a73fe285a4f0b7bf31424c2185afe47865771d0e6486d2d3347db676
e4e11d5dec4a2c270ccb6b205a4832d6395f6180921725d30a878b99544a9969a22f29952963474c819795593063fae1dbde2bbd72ea86
97743040481612066b42f490b4f025b0bbe13688c3e90d0f612c217ad0abd4333e0067528592bc97ba50286f4b1493a5573653289f9287
1b427950922dcb353981a507e69ec484184cf352e286a5854ea90b104607668ac5f2f24070000591f8d095ce492e1031f74a736a5fad1d
e920751ab0481e86fba0210adec2362f25b58b1194c2e1e892d73d33e8d0bdfc58dc4f89503e9ffb3084321d5a2fff691d804e7f3f68df
33f2790c099b65b90696a61c6bb28253ae8b9383ea4371963d11392d221d2a9bf127f365be969e2dbf9a585e9ee098e146c64a79656d6d
4841707588262259e9daa711975b4aa175148cc4512adb30dc8759c247372ebe7c2b07348c6b2e94cbe13559f59c62f392db13ec5a335c
e1da5f2e86d60a429907fd26b142f97cd5e6f3b39fb2de5c24d52d4b173bf5a0377572303a67425cdc13591546281fd0dfd4e0dca9c06a
ac40cc745aaa1518137d21d221da432c2f4fcee7bdec49b5b9164a47f5cab3b2e1a8fc9fd334fdd7bda7398e400da15cb2717c7111095d
c57209f31ad41664cddf9100ce7fd2acb33d60867b4fe0d12e83fae966edf3520b302098e2c4725d37bc1be495bc6cf80186ea77023f34
e88771267e54ca7a86cf4944874f9fd688d835ddef0bf3769e7844f901ad09178318cb32c44e86eb4026c396b8a8494678d0eed7cc0f21
962e1a3938c0f1109f84b9158523a7dd59f2d543876a614ac43a42149e8f7a3f400c3757a14c5dd38275e6ace0f4660aed6f4768d83f34
521c8e61c5d285f190ae057421a42f1d5f0136ef6872dd6ac31e7f14a3776a5d653bd6e27998213b0dabb9dbdecc297f5db5e910b58a86
1d86d782de7a771771857b19ea1b2ff18b3dcecfd1c3f6aac34ab1c324006ee8cc149b101b6259818e9e7c787c5754584f1e6288121a25
b4aac985a0f49d121b93c64cf7f5e7a3d2e881dfbd704e4c39437d48134c8865268dab5e7a0b3d81726839e73f95460387f3aed3b4c85b
8ef61495de1062194f4381bcddbb38d391a9289e878de92c187d8959c9b32b6f780de6ba3434c15d4b286faf3ccd10ca8e509da9ebabac
391cbc9993e14abe7f47b8cfdcffb33182b768cf9490a407d6964e402c3dcc19b82bad577f0a711c07630d2e1ef41b0c177e6d2142c378
d3297f4331a3aef4c1307c41c561f6e16a4c1f01e684b862b2b68756c771cd2186ce51e37c1cc0b2bc6a45720ae5b785e508a1dc217359
687aa0ef15dee167b23425585dde3942c678cb98f8e5921544ec1c562c8d55b0988be19ac85f0a24d6171f0863152e84b364d58b04d10c
4e0b954c2198922191827b2e6db87eb861788dfa2ea86a095c70e468edd1b70aa717d62f8a9c157022a3060e2396154412950141340cde
e4a6a2d92957abe8d0aadd0fc317663f87509e2f86d9104a104d695dfdd6493c564968cc7cecb9f5211e96495b44c7a0eed6b264b62491
2016b05262b9895c37fd33e8fc16fdadf9b487321fbb134be6708e43acef05fd2810cddd672697c6aec4926b9e050e1bd092124710fa6a
3b763167a994fac22194a5f34a0064f6bb2f19cb002fa0b94ca4e46bc0d0962557213308249044ee8b1ffdb82ec35a96948ca054286fa0
830169509ffc907a59a3ad351f8de1c4d258930cab6ee6441670de0091985adc992ff2ef94e3003033d430dccc4d4ed374c7da110f3c88
603832cb94a0af33338c588eb60202004e94522653fe9fa98744bfe743fc307c4ebe6bed88e7353a0c181d2a2b929c4bb3d3b2c55d22da
b22cadcc0791047b244700f12c9423562c0bcbc99e6aadc5a4a307801b08667b448a65e1d001c5bec0218060b6459c58ee2dc100003581
60b6439483a74028cfd101c011c9e9f7148207121123960542798a1a37e0c86408e61d4a61081210310c2f10ca97145201c0e1497d8e30
1a4ba3bb65a994ca75c6402801b8ca89b525006230d3e82a964aa957d334ddb4766c03a104600565234f4af386f46ef1741b861764347f
a7b5c67c0b001e301caf434fcb324728bf412801d8e4714a1361381e4717cb92e6299387df78030210478600a2b4f306cdc53237bb3984
128078729e333c63617a0cc373b29bdfb6b60000bce42cf92581051e9a8a656620ec19d67a0390c5f3c42f15d7b3da334d87e158c70a40
5be019e7a39965a9947a646ddce669abeb03005c3ca7a8dfe3a1996509ab1280f6e494a3c073e746b2589ed08a040040c3670f62e9a6c9
303c67491584120020895673960fad2d615e07f70200aa91e95fd83d4d86e1180600d08f8c2930d4b07250ddb2544aa1d101180b24d376
d062189e9a4aedd4da020068c91db4b64df56138bce000f405617b3c88134bfc4800f002b1e4415475470000900ac41200b0067e030710
4b00764c66a62ff80c1c60ce12801da394fa314dd3b5943bc433e806962500fb264928811f8825003b050b4278815802b05f3e66dcd94b
6b0bb8a0c59ca5992cbe6fedf080f9120078407c252fe2963b2aa5ee5a1b01004964d6023fb7b6804baa8b65c6d245a48702a080dc2a8d
5a6b24d00880146d00ec181aa999e4db37b7ee12cf5d98560e1eac0800a0035aeb2f5aeb5b4608977fd334dd5b0ebb2194dbb4b22c8d79
ffdddae1e7b9d6fa95772f000034466cc132bce90000926819679914bf953b490d0000356866594eb02e010003d37a054f52d546c45c02
00a4d0d4b29c605d020006a5c7daf0c7d696004aa9f7febd0000d086e696e594b714eb86d6fa87b51500001ad1452c270cc7010083d133
45db536b4b80ccc4000000c04237b1d45abf9da6e99bb52300a57b030080e6744dfe6bd6ac5a1bc3dc574a25a57c0300000ebacd592ec9
18623f25cb1400009a20422ca73cc13cc9c8950900005948aac173dbda12e663664d64000048468c586aadffc910cc3798c30400b440cc
307c86b20d7db57684f9a0b546390a004035c489e5942f98e7a8210200a88548b19cf2b2ab5f80953e00801a4872f05cc1ac05cf113ee3
5547e26000003762c57226d352fc8ad53e00004ec40ec3d728a5fe8929e7b906c372000007e22dcb195a1a99940b73fa352cc76a1f0040
11c358964b723310c1ca0400e4328c65b984442f2963d1f4cbca441261004032438ae5f46b589ebae2c7708d44132b7f0000d10c390c5f
43d6e2356b471cf7b4d65f3a5e3e006000762196534110fb02882600c0cbb0c3f0358b20f677d6ce383ed3f01cb5ca010016bbb12cd730
d4ec41826100c025bb15cb1906d13cd35ac3da04e0e0ec5e2ca79f8269c4eeb3b52311c46902705c76336719c2386e48e8ee053eb609cd
699abff75d6f0800d09c4358966bb82c4de2b1d61ae209c0ce39a458ce30841bad79aeb57e656d05000ccfa1c57209832368cda9d6fa81
b5150030248798b38cc1cc69d2bce619d321ef2fe638b1bc1280c18165e9a1c2107d0d2c4f00060296a587794510599b2fdd9f2a626d79
9abf2f24d2000061c0b24c4429651c387f343e2dac50003a03b12c4029656a95ff5dfb3c088607a03f104b46cc307a9aa63bcc87fd46b9
3b01001d815856842a4cde2f3903ac4a006400b16c8852cacc3b7e4c3923c4120019c01bde10adf5a7d9c3befc9ba6e9b5e72a72ca6600
002ab02bcb92c26e5e68ad9f593b0100a9cfd393699a7e20f7c14f76390c5f84f760ad3600095084c75b1249381617ec7ace7255c80c99
cf0170b010c8f959b9adb5fec7fee4b1d9f59ca5d6da0ccb4fe87fdf2c56ca3cb13e0cc0813002698c094a20f33709e54b9a4b87503a38
8c375c2965de9cbf5b3b30540707c16141ce20963782c3850e29a5cc5bf3a6b5e327afe11c027b8246516f02b774c3e441b0b6028bc385
0ed11bf486b5e3277f2c86ea6f91d4028c887170cefd3820942734e486504672c838cb458df1136be72fcc90fdfb221b102a3c029128a5
6e99ba500b810c257a99e7253f597b4010ace0891baaacc13c27e80acd3fbe0a4c29ad79a7b58663b30088e5820cd1347c98a6e9193c88
a0263425f4cae3a40c0191640262e920533467cc300725244031d40f5f39bcd73140249981580628144dc3375a7e896078b0090dad5f14
a6f9834856026219414eb6200f104f70099338ce60445319886502c6eb384dd397cc61918b7333cc42273f064a2913c3fb2cc129130396
f1360262990947625f0f1f484011da3130146af62cc321138379c9de8553b12d10cb42a81ef89f954ff3ce2c538380ca8484f1c9467c23
071fb4d68f0edaccdd815832516188bec5e9344def4944b10aa311e4f47b526954e103436d01402c2bd0c8daf4f1bf4ed3f43f1a21d55a
7ff17c0604a017df23fa6b298a4b60450a036259110a24fe54a1e2632ea7743d9f8e3ea4a70887f9af9720ae3173918f8efedb480562d9
087a38df371ca6e760429bfe9905d5fc7b342702bda0ee92089affde12f4b2f281e5b30300b1ec0043b0bb244e17d7f2c9f3ef2597024c
c35d5f1ec5079e7f4bb1024bf99fb5d6fff9d8b7702c20969d514afddfd334fd7b876e84e370992fd564099aa6e921a54af3bd58802050
0ab71373cec1699afedbb92cee344df7a6693a3b6483ec97a78bb2c79789a5c97963f2aa7ea45a514038b02c3b4022396d15865a54a904
e3704a59a8a223111685f560650a069665434cf0f22c943185a18c25b2b03a4f60758ae49c1c34b3f5f82035648b0aeb9d919509b1140a
c4b211147bf97922a14c3dabb138b4d67717e2f918e2d985f345b671f3779dc3936d7e5bfa3def63582e130cc31bb0aa2c59a540148526
3d23a701e0e38cd6ea375941b318929f93c5098400b1accc6adeb1e99c14a5007bb6a3709bda9c512cecab9e4b481773da104c41402c2b
b2ca83794643added77477b1944f7ab0764dc42627a1f8d3aff4bf22fa0d80585665612164cd5376b8de478b2580bb12d211da7f89a928
baf80db0c2470010cb4a2c8572d4ac31ab7b98cb062f97135eef20aae794dd69b6089debdc57d75e659eb8367bb8873d01b1acc03a3e72
34ab66665d4ea3f43e560f7f0c59a5129452ff2cb2910f3befb7720c0edb8ff6024287eab00c247f3aea4dac2db611620069be6f59b6c1
b7f65c3cebc263945300740262c90c593597ec2069ebebc5bf47f0aa7f5dfecf0e86aedf16ffde4bf295218158f2b3b46abeb53a692d96
eb9927e1d62559954b5e5b1f1a8f2b0980697103e800c49211ca24b3e4598bf33646b275b9b62a876f7fc7d2c95e19f80f0fc492972bab
67b4d66bf11c9565cecab9a42be884c382060d80583241c1de7b653deffa97c0f6df731ce2e9eafff7f2121e0a88251f7baebe37c2c3b9
4e65b716989159cf131f79e5553720967cecb603bb3ccab4da4732c1f4778381b46d028058d6634f968d0b31f3969e39bc3d89a505e68d
db03b16460e7f3953e2479c58f18ac8d9ae28d8158f2e0eab8488bd68e078e33b9accd3d81fed51888250fbbb62ca916b7645cc2b127b1
74bd0c40630e2b96662584598d42ff2d1583bd27687559ced27109e8a8402c05f0dbd16e9896eb2d1f24f3ef3f9552c84aed07091cfa82
21b7000e65595242555fc7bb9691422cc800e135b1f8da0c0c008d9ef4ea0f81ed891c6d18be190bc99c060bd9ad3bb287553d9eb0a894
effff0ac277fc86d1cec9dc3882525528d21270d962f68f8a6b5653006119c736bcb4fd6ab7a46c4d7b736a10c45d7429f5ba714047e8e
6459fe6e6de1c35b547f07095b47109c50fb8f3e0f5df2c27559949cc73f14081d626023bbd0b0095ba9acc40884460dc35a4e1ba3a177
d696fcf3c0db1e01c4d286ad13ce0c6c5d7eb4b6fc424c62dd8d6cf4d706b62e43a3a1d03da7826887080e219689cb116bacb91dceba8c
c8c83d92f3e4bbb5453814b9e1c555d1b280902803e2289665f49bd395612792a0453ae0447a70be4b6b2ded7e7c4e9e0b46f28cd3cb9d
238b55b04d401a4711cb58274549cd9c2d8bf4e628c3f1889012890fe156fbff315086f1cfd696ab045fcc0b30bc6604739657c90e228f
b448df48cf504471795b8873086ccc5bce7cb5b6082326f6715d22d7c786e3f10a70f26c03b15ce0280e95caf388cf7f96ea70a079b260
5cdec4d34eb5d8cc212a39103bf245550ba47cdb60f762d9d292d35ac7ce8b7d9766619250c6cc93bdb4b608416b1d651dd1b23f512f2c
12f1cd17d5344db7ad2d61627faf3d04f057e5089665ec1bf3a9b5258fd8cef9594a0d687a50a31c0a5a6be975ab63e79dbf4b187a1ad1
4eb176531d6b03fc5ec37004b18c0a058a9cf38a394e4ae7fcb3e7d0cb8845e2b094eb85520dad758a13e7e356884e4de8659912d674c3
da029a7104b18c19da707392727d342c6c6a0190488782ce2db85e280d480998bf43eddfccca5c5893c1f0ac156705616d1fac2d0e0e5a
1e251a38787ec2ba6a87028653c390fea487b66ab807253c8e9d1fbb446badac8d42d15ae72c2cf8585b3449247fe404c96bad4b842cb6
4f21d428c0aec53236ae2e36142385c4e1e092378a39dfa0c9ab39e732cccc4d297ef8bda640dc67d1640b62372f406afbef99239d54a7
ce15122c52387902ecddb2ecfda62c99637ab848d6fa252591b0794998877d21907f5b1f8ae76ca0e1f79a94e990357f2cdaff538ac569
86b32609c6a2fd4b96bbbe645a2d55b2e0e2f01894d6fbcdff49439ead37f959e1102708cd036dadc8900a6ba98d8c18c797a5de5cb210
47b598d8fa2689fde61cf548d32dadd9bb651933e4a91a8c4b01dcc30d63b985b217347fc99e49aa01ac2ff1d8c41b032d096dcee11d3c
2d1242d030f69eb5432ebb2ade4673d2230966d5d1ce0670f278d8ad5846aed0689610822ccc11e2e4cef658e59204f3b1b5431e2f2b0a
65cc72dc1a290a77c19e2dcb981fbd69c7305e499a13929a3aeb71478ba63a945842f20beb76cd153791cb717bc4250fc19ec5727338d1
cbcb4b969ba435d666d8ad52b2d4a442f5dabbb3786145056a37e21bb53f8a870966cf6229ba1093b120e8a1ed1dd271d268d82daaf6b8
d6fa115999bdadfcdb0531b9396cceddc2c9e3e6c80e1e113564e841b9d141349f923553dde2939a259eaccceba541df999cf4b0262317
60c0c9e360976219e3dcc95c1257057a686f91a559d36b6baca87bf490b69c82906ee5ff436d527b786e5e88375abda40a80583ad8ab65
39ac47cfbcf9170f2e87f57b468e1b73ccebad13f7764e689b8c199e53dbdf6612ced3c50bea5641320c4ece368e855ae20e76b9822722
91ed079ab31a0a5a8531ff5d5fdde329d5c836f7fe5e82b3802cfc92ca8adf1acfe70559b5ffdd95e7f894dafe8bf4e5a131abcab092c7
66af62b97553b7e179ac4fe472d3207868eb10f18cdc1062058be1900e1e08657dc8a35a1cb3376009e151d88a024070fa8a238a25b2af
b481ab92e24da905de06674b0ce1e459b13bb1544a6d7502bc312b5321816ec9bc27701031af0a27cf8addcd596e39773007569f4ae566
ef092ec13b245bbf139e95abec7118ee154a509f8a653182de5b9085888519a370b439cb98ac2ba08c92ace0416ad7273a1a5b0b3322a6
b40ec51e87e1de1bc2b0a22ef470fd55f324f80d79093d2f9df36a8a635796252c8fee5415cae957ad6dc047689512a6b416ec6d181e12
cb114b0b0c4343114ba9b50db6c1503b925d0dc3378614589150918db6e7e6546b5dadbef7d108fd7698f6f8c5611c3c10ca7a98b2af8d
4f292a37e60ef02ed48093e7174711cbad2c2ba08cdf5bb71f9641b21212c4e112ced4623762b9e1dc09ed0365edde2b2f23964132b151
4e04563cb19b394b7a689d3f2ce65dea119aef6ac0ae4af6f624f43be2f9f9c99e86e14ea1145c497178040c85afa15e0c1b58b0b1c111
e62c43f33120131a024b48b6c095dde8d084cae46e4c711d86dd8ba5f4acd50323c6c15221cb11b8cae1c572da8b582aa5e0b16b085995
928af17fb4b6801c7c0b370eefe499766459fa86dac8aa520771f925b10c9205df737478a6bd78c37d9e3c78f1eae06befdee0f72ec7f7
dba26d773e674955ec003f123da7de55282009e7020e4c75eddcb2249066aa0211e5869b02cb87874099dcc3afc73f8258cea0fc2d3352
86e3104a1ec87afcdb77b0a3b7f35e86e1f7ac2d365f9552a1655d2011210fcf0d6b0b48861618788512ec442c4d212b7a707da10f330f
8d358435c5acf414ab7bc82655861976d308616b81c1636bcbc1d89583476bfd8444736bb2ff3b424d7820b17adae1d44f51edb10ccaa7
e09a9f5c63a6b00e3f2adb5d0d9e19b21eb7e20191888189c60e9f77e6c5686d0551443e1b13bd90b0028ed86de890b178c8ca3cb176fe
e21a0dcb0f1f16514ac388836f10ca7c2851f39650be34cf0e84f22abbb52cd744541e340f2132d81452db437e748f6c0911bfcde1c383
421ca9acc4ab0d27d04db2321193594635870f84320f63286c08e539599210ca0087b12cd750a884cf0388376c0194d2eb0df36151702e
03a5d48f8da427883f8ee43096e51a1a72dff02407be8f10a37c68aeeb94f19008114ac4ccc39335e913cac7644d42282339ac584ebf9c
40d70341ed26c4c89b1415f821cbdcf5224a052142896c0498cfce9bc38702a572d861b88bd0f011f365796ccc956df1efb4d6fff1c667
001158d73d616aa99c435b966bccf0d1e704a26139f2fda553e2f0f91fac2dc04920c01cce1b2620960e162b81d6e9aafea2097310df96
3fb054ae1e665e9dac775736f3db5874c107c43200055aaf9d4008644f84e6c72c6b1d9441f3e9ae00f313386ff881586eb07002dd5e7d
f26f01a5608701ab6e78216bf28fd54167e7cda7d1ef4f2210cb48cc5b9a86e6cba4110864074df104989f924822394c45209689789c40
9f69821d806ad07cf972c9ee37386fda01b1ccc4e10442203ba882a98bee0830bf815c066d815816e2700221901db041a9ef9675d167e7
0da2321a03b164c0e104fac331af0440344aa95bd487e61ca170de740662c9c8da098440769003d58afa4a5f85f3460810cb0a2c9c40af
11c80e625904983f84f3461e10cb8a68ad9f91687e212b131d1f38590598c3792390df8ede002d98ad0363619a40f6862518c0002ce6b7
4f3027291758960d2127d023b23261391c1c93e58a8412ce9b0180583666e104ba4b13f9e080d03cf623386fc60162d909935c426b6dac
cc1708643f0e14606e5e98d7cdef7ff4f61809cc59760656c571a000f30770de8c09c41280cad0fcf42d38f6c6066209406528af24d2f9
0d0ee62c01002002882500004400b104008008209600001001c412000022805802004004104b000088006209000011402c010020028825
00004400b104008008209600001001c412000022805802004004104b000088006209000011402c01002002882500004400b10400800820
9600001001c412000022805802004004104b000088006209000011402c010020028865039452ffbb52eafaee6f14801d03b1ac8c52eadf
4cd3f41f69ad7fecfa4601d83910cbfafc77d34fd1fc5f767e9f00ec1a8865459452dafc87cef09f60280ec0b8402c2ba1947ab53ab211
cdefbbb941000e06c4b2026441fee13ab252ea91b51100201e88651d4216e4dfd6160080782096cc28a53e6d1dd131440700080762c988
52eaee344df7238ee81ca20300e402b1e4e573ecd1622c5000801c20964c28a5be241e29c6020500080162c9000dbfefa41e4929f58fb5
1100201288250fd1c3ef153711a80ec018402c0bc9187eaf817509c000402c0bc81d7eafb806eb1200f9402ccbc81d7eaf817509807020
96993087fec0ba04403810cb0c9452b72a84fec0ba04403010cb3cbe5638e6356b0b00400c10cb449452ef2b1e1bd625004281582640f3
8a0f2b9ee2a6b50500200288651aa1d46b2c28a5de0abb67000ecf04b18c4729f5a2d1a97eb7b60000ba03b18ce7cf5627524a3db33602
00ba02b18c4029d5ba8ced5fd61600405720961b28a559dc64f100002000494441541ef408eba1584e0080102096db7cec74ded2041d00
00462096013a673347903a008280587aa098caaed9cc114604801c20967eaac75446803022008400b174d030a632e65a1e581b0100cd81
58ba69165319412f0713006001c47205925900005c402c17506ca3b864164aa957d646004053209657a991a792833f845e17600619f3e5
02b124a45b6f78880e43cfd85e100062f90be9d61b622e8f4169b550500988659f441939d44c3a0ce4f04e29f508bf873c0e2f9654fb7b
88a58548ae7108cc08e2efa33782440e2f968cb5bf5b50adfe0f9081d61a73964239b4580e189283f9ac838004d0f238ba65395c480e86
e2870109a0857158b11c78a50e86e2a029c6e164d2152aa5f4ea2f3ae7aab1941ddf77fd89b5a895d6dadab877c83a931a80be89d65a09
bf4427e6c1706d0ff0526b2d26a9492b288faa490ff85c6b5d65aa88e27697c264fedf383b9356b16df545324a7256c59d6aad452591f9
cdda720c8615ca893abad67a8470a7a6180b486b5dd5f23659a01a386166b13443f15af3ea2f6a4e4329a5de1786bbdda797eb0d297dfd
70c37049e9d70a4080ba9b1719d6abc473b4f0883fb1b630605ee4d43e5c71c1dfa5ac5e3be29ca5a4f46bb92040ddcd45b4406531bb5f
fb1c4bcb55295545d4a66962b7d62898be46d26c0989b88f25969d6bea808634b0fe9a9ca3e230fcaeb5a5002a81522d983ec599548bc3
88a5849a3a9c200e6f9b16cb581b9ca3caea32330f48ce9913e3489ba6e99bf5a14848c86a9740e91e637c24cb528429cf08e2f0b6b9d6
608efa5aedc50d35e7eccc90df441c68ad6f91782689260965132123a751370e2196a8637368fe6ce020f8a3f2399aad3433a2696d0c13
2b94a7d3347da0ffe6d275aefe2896e52eebd860354f342d4615dce7588aca88553ebf51d88fa23f1372f588feabc88a3db3be2598dd8b
e5ce4b321c2e603b97162163cc7d6d5467e43989e1adadf848adf5dd54c1ec394a3c8265b9e7920ca82b1e4f8b9031cebe76452c2995a0
74de69ad93a6234830538058d600951ac19216e1278c7d6e7dadd2471146287363423f585bfc402cb9a1097771951ab919c4e290420baf
2d4b9f730c61252f44382f10ca69b5465d2c7bb62cf7162ae403f39609b4186d1c6d44933af4767c7f88f6daa5581e2c54084b1fd36831
daa8720ea1153e4b42818662af96e52e4385c026e7311fa2a579b944796f2b2dad95385c3d4c71b5dd89e54eb20a258178cb4b6245b024
8a207615498da5b5e2c4d231b7ba5bf66859ee21ab502a5827fe93688bb160aa26e51cdc56d7105548f7caaec4f2c05985f61c4b1a8dd6
3a2534286b9d71a23302256d77c4de2ccbdd641502d519c14ab39c27c2a65cacebdb33bb11cb16e9b8c010443fc005c3e4e820ea0ac97b
6b2503061bec422ce96d7be8f91c6456ba2465fe36772d778b734c9ef5e187f13e4b632f96e5d005c8988093277dde322b1e3271de92fb
25de3d09ee51195e2c61515d82e0f45f44c55b16929d591c8cc91e2c4b04a08335d1f37a05738ad1c3e1821220ae6138e8c4d06279c400
74b04d62edf02cb14c1ceeb34e916011421f46b72c8f1880ee050fd11562bde225e166b15e71eeb5e270f2746058b1ec5dbc4828082b21
4cf9026b23ff396a8b96cf7a8558766064cb120e0d1b3c445769e1e8a9768ec0ba6b2cbee8c09062890ce85e1056729516d31298fa3808
c389e55132a08372c832dbb4fc4ac2cf5a9c03c86044cb12562548a1bae51799295c62e25e90c0506289658d2015b2fcb602c839acbe2d
ef3b6a250dce68962596356e80e19e8da9616d6de43f07da7de70c2396108168e01177f3d8b9959793563783e7a13d23599658d618071e
2207b4aac7e788619907d75a7f6ab8661cbf73638610cb0a3901f70cc2873c041c316c4ec3c0901febbc076714cbf28db505803c6eacbf
45162127d6392a447148b02c7d2f9fbd9df302f1628964198013f28edfabd9a8ae7324e6c08cc167c116919821e98e52aad8cb9f98b1fe
4eaff9da112c4b24cb00ac50c6a0a7355b95e31c1ba2c0b630c39c4729f54529a5a769facbfa4098cfe67ba6163b2d18893de72d93df81
ce995ad8ed239df33d8758c7225a2c0b8be103e0456bfdb6b6879cce51d58a2d81044793f3b474aedbd462ff4ec7f4fa18485435850196
e67778b8106bdf3a7a36a45b9625c5f00fcb86450208e321d75aab9aed612cccdae7c8a12021710c211f43ad67fa5a6d2b53ac581eb806
380748eeb00f6abef4f6b8fcb2ea3d49b62c91862a1f8825d86277c64885a8862bfc666d1180996c96785d0381613808628445a9cbd981
6f14da348b8d79feaecc01d2e7d7fdeac1eabfa906ce299d677edeaf889de79cb7e8ef3aadb7bfdb2a5f8448b144607531c870b30faabe
f452e7521d965bb225c770ce6e881b8623b12f0b78d900c08c28b144625f00ae80797b4148b32c615502004422462cc9aa44625f0006c1
3cb3739079e4df9750c0ba74245996b02a011880c532c5ef8941e6662efdcd42586b06c6b323422c615502201fe37c2591e42a43fdd7bc
aedcda23102996e5776b0b0040044aa9572492b59cafbf93688acef2df3dce928a900100044222d98abf9552e78124cd5d916059a20819
002b1c2b579a4229d45a0ae5ccb54ee7dda4ab58c2aa04c04bb7d2b994bd27d6883935ab727c7f9435de57fbc88b44c1ec6d59c2aa04c0
4d17b12403e6b3b5c3010964d0023659e369589d5cf9529a6076134b58950004e9e5ec883160ce33d778bb6a13056991d437969e9625ac
4a00fc340fa58bb5e4721d30aeda44119839cc5739e7e3a68b58c2aa04401609850193adc325549be8a5b523cc1fc1bd8de86559c2aa04
a08c6fcced175318f09cacc322b4d6c9155b2564236b2e96b02a0160814d3c12926d079d3989dc4efc7cf76c643d2c4b58950004e81063
1995ff9486d02ce4d451ef5d97aba958c2aa04208a18b164118eceebb253c389bae6f76c6d59c2aa04609b6662999035e8d4da52484ec9
88dae56e4334134b589500442331437aad21f099b5254cb730a2969625ac4a0098e028e49518bf584b2c53735a767b9134114bca570900
9045f7f84549d51bb7686559220b3a006373f85afed5c51259d00188a7776ab600357d0e4959897af93f5a5896b02a0188476a41af9a53
69a9e14bbb154b589500c49352006c2fbc1fe13eaa8aa584f59c072539d92a180af698c708aa4d0f8ce2e4a96d59765fcf79500e3f19bf
737a888bd4b9d4665413cbdeeb380f0eda7edff478194a0c946f4a4dcbf2f08ddb1131d9a5413c4aa958e7cea15f86bd86ed55c4125665
77300c1f9328b1e4c82909d2a96559c2aaecc848ab22c0155a3f37a9f18d879eb76417cbce299f0000f1a426a590500ba74724c005352c
cb23c68901301c19e51da29204a792983ba25bd03eab5826143d022086c34c27240c713f585b1a5269a961b400e66458e782dbb28c297a
0400b0893534b8a7b9de595bc2d4701ec6de7beab5b2a2b48e2a15bc09853dbce97933e08253adb5c889f8d8bad40b4e8ee2ac4aa8d9ad
ac8d8dce5deb1a7ade7b0a9c9625845206f084835492caea722e634e882ded6a554e5c9625cd632013ba0c445a633427f7d1da11a0b725
d10a72707c8f395dad36c9b0faef71547b1cc5aa9c182d4b08a51010633924b1a5156a2648796d6d09f339b83702a5546cb6a1d42a9055
28164b948c00911c3e11438058c768b51866ad756a2d9c1c6b74f95d53a5f1a1b5c3e6548a01c0615962691d88012fd572aa2ef8c819ea
e608a652ea51a4657a2ec959c9219648c32607c9792cbbd57bde0b1c738411dc48fd8211cc58470d3987feb67638d05a8b7ac116892512
668843f25253e40b7090588eb63a94a42359304d340c89e62763391a87dee2ef2dedd391c6d5b944e75e9137bc64ce0254416c5c624e5f
3982373ca15dbe69ad9bd69ee9f47c9f69ad458e42b22d4b696f44004ff8ce69febcd1cbaa65e28ac75285722ab12c6155ca43b22506cb
d2863cc25121383ddb22250e3413b1abce966459960951f700e4d22d15574386486768e63149ac6f333b118d48aa1184722a18866369a3
3cba2f0763e608954163539e89887230197f8c879a84f3b9f58138ccbd3c1d4924677eb3b66c502945132847ec1c7266866dc4effe425c
ea43adf5ab659fa32985078e785ae35dffd428eca92ac962894e2c13e19d1162b922a5a200099368a8ffedfb374b75f0c0b12313e1ce9d
4fa971960770ee443f47474928229da4394b04a18be54cf8f521203d1fe9bfed614875f0a0d3cb04e53c0622b1a85fb7c813e39f30a562
e6d537cb3febc37ce7bc4ee7fce1386fd7e58fd1c370aaaf83b21102913e4c438ce555240cc11df5b21e241a431fb4d68facad0118ce79
de73bd788a5862ae52283b14cbe64bfb5a911ae02d28d9ef9ae440720e0de9d9d7a386e114160064227a4e8bd271a5b2e7daf329f3fe4f
ad2d0c64867215d1e39cdcc4ce5916674506d5484edada989c39b73d3b12a36b6f6badf7fcd2188ea8613886e072c17ce53890951d95cb
b1763b60189ecea66589702100d88816ca5a43f019129d7bb46cb1c93afc1ee7e464d3b28455299ad739b5535a02cbf217d203d129a4e9
776b871b964c419494273ad78458cb12d98564338050e63c4cbb0cc24e1ca125d5f1e6426bddfc791f695e362896c82e040ac90996dfab
5323259e10953005e2154b94b815cf8701ae3179c5d708492352514a258d004c2a346b23e88e572c915d483cd24386c02ffe4a688bd7d6
1620829058a2c4ad60a45b1fc87bfa93d4795be9f3d047c629968e359c401623384172fad008530ba97c4cf87c17c70e88c3193a847021
f1dc1ec0b2cce943624bf9e640d6f5d784afdea0badd3daf39f677632b32364a6e4fcbb2c4f0493e7b7500ecb0946f8a504ebd851284b1
c472e7eb72f780f8a16a66f28c5d911163fad8da0244610dc3310497cd08ab5b9452ff64380893f3234a26f53992f2bb6218eee78a6599
1a0f0680879c488addf4bd8c956fb965654143ae5896b02ac5f35c7ad0768653e3823dad071fd5aa9c60590671cd5902a10cb2bae5bdb5
e54028a5527fa397d61620924bb1544afd5bfc44a21925062f3ab9ed82aae9c81af347cae9b4d688691e84a565f99f1ebd318423beb447
6ef991bd640427c7560a7b7a49ec9e0bb1544afd67476f08e90c128377d82138cdd52639b65036622c7e8353670846b14072bce07bf104
a73ab56e5b5b8068e0e01980112c90dcb0b33da464534aa55ad4e748c3361eff82374e3cef06b9ce943464bb8186df0f53ee476b8d5cb1
03f22f78e364d323d57f2a05f904f6b0c42f75f88d7c9583721194ae9432ce836b476f0c817cd35a8b4f6c92db7f460f44574a7d490d95
da53e962ae7b192d285d7c58ca11194128899c17edd085c92851466a4c299c3a03732196986c16c9f9081799115b38337a51ae94a4be13
2d0fc473b662a444e3976bc353ebf782ea884ff03b15e413187908bee75ae8ad87c423ada3bf0c1d4280ac2c0611cadc6b1cd6b143f394
a960f8ed20631d7d57d67196082392c1280f5756513badf5902b7d68c8983a4ff961afc3efdce5add3af52db49ebe87b83e4bf021924c1
6f6e04c5cb11c3d548183e5b3b36186dba21f1f9cf8ad620a1fc6eed88405a8a362422ed8b78ab92e22ab342cd0615caeb47104a22a56c
c94d6a9b68c837922594bdb12ccb09d6655706b12a73fbc7bb1182ecd764deef3dad75cefc6657a87ed2df29d710d367739342e79cab16
2ecbd270626d012db821bd954b4a8f1c4828df8d289453e67cb26923a5d4a7f51ca6115edaae3984b2374ecb728275d983f311d60c1fc9
aaccbcd721565d8510feec779bf30e8965f6242c4867e74e9dbd3b3a2ed9432d21f2faff69eda88069afccb6366becdfb7ac35ef1b86cf
c966c5d7a8de09a7d26f8396f7e5e60f182a24edc842393574c215b697093bfa4853004d2c61af6579f9010cc7abb373a7ce502252709f
3706c9661f458ea327812b534e4c1a53bdf2a9d7b25c80d50775116f75d1f03b977b7daf3e1e08e52fc8d15323d9c9cb4a73f3d5e7fb37
2dcbe9672732f302f7ad1da098015276bd2a5869318ad3aa647e7ec810a15872d2d079f0f60526cbb2faef1025961386e3b5109d2c2377
d5cacc20d30b25c3cd5d0be50c83c3e724e488c934c68cd56bacf98be3b698674d114b78c779111f6252f88214bfac51296592c7fc6eed
88e31042b924b1bd4c8ac14721915c1ddbd7d79a7bbd7d448be5f42b20f990b556b8d953c66c1703dc5f497580c309e51a8a8e58e724fd
87846d57f3b733496239f1ce611c99ea9ebb1218a65cc43a3c1896ddedce9903e2482e85abb546098a42840b65a910bc142c94af4a84d2
58cb10cae3926c59cec0e19387e4e12943e13aafc7b3377b9f5600f549b62c17884ffa2010b131951c153e250aa5b1260b85f20c4209a6
12cb7262082d391a521f3aa652c8a29c1e4c29c144cf2d83b6fc567236f37028a59ea2d0d936828592633ae5b530a1fc27b7e4c5023872
c0154a86e11750a1b377d60eb0e4a9b4d63071b34c426986a9d9392e39317180744f2542f90d8e1ce0a2582ca75f495d21986ebe49ab9c
49ab5658161848888e302b4c48247303cc671e8f9e8b12d4a368ce724de18a885d226df8cd1927dbfbde0ad7ad2f19620dbb54687ef809
05a9bb962d9ed2b2c44f1256e2e4c22a961304738d98b5dfdccb557b0a25736297c7a396e6ed09f5a72f05531edfa6697a3652dbb30cc3
97d0905cdc1c5d075e0b12ca57ccebfa9b878d99c80be3b5a7e13687509ed2dc248432019aebfe41fda9646ed87cf76f4adefb83ac53d1
b05b963354f2f2a85e7211c33aae8a7a2b9a7989c97a31239587d6ce7c30e4cea451b9896086a29e5413cbe9c0719812e62999c267d654
17ca4a0239233a259e641ae784102998ecc3f025147b77b4953e5def77517a7418a1a492a9ffd0757faf2094b769c80da1cca043f29c75
3623115415cb890a9f91a5756eeddc1f27bde2f3163186ec19edb9e30e4d7aaf59d4e99affae20eee710c97228cc2c5628cfe81950f39f
f9ff8cb04291625975186e9d6cdfe9ddbad4c5564afd1fd334fd87d60e1e8ae7f76858fd8cfe4a9754c6f0416bfda8c1790e41e4c285b3
9878db8439749189b19b8ae5b45fc74fd31fb7511b5e8a0e257abd4b45a1ccbf6f55b0044b4108103391215ac946428cd12471797073b1
9c7658a2a2c50f4bc3a1b78dacb351782a6d75d49e88b12a73fbfed6b1258a65f5394b178b79cc53c7eea1a8f9a3cecbf816f37a10ca69
fa48198e1484b21e347ad9e2f9c6fe10c395d8ee62595eb98081c38b38859286ba4fb0fa2919f3c27d31f2323a8944869e1585f884ac4b
58960e4c78d1a05666728810ad4279429eeb1fb3d5489de62384320b33a7f671d196c83fc9438b39e90fd616c174b72c97a0dc2e60c63c
8c4f906e2d9d90d5b7a028f28046531fad1d70f0c473f0a592a00e5d42bb462536d769a9a8f9ce8361782466e29e1a0b39320117bf2f86
ea104d267c62170b3de72f691a4ef4549c48cb720dd2be814a4405531f910c11149b00830b9196e51a337ca237d06b6b2700f9dc59589b
2297d80dc4c7bdb7e3106239636abd9068225f26e0e62386e857c8cde530b7e3ee62608712cb99794e938473a8f003209e37f4b0bf38f8
4f557affcb39e25d08e710739631d022fdf73b4ed401fa70d825954cd53fd7bcd45a0ff922da8d582e21e17c552981ec51998b4e99951d
ff704ce6d3fcd62dfaf315bb92c2ee1d186b2a25909e313578ee8e1403bb4bb15c430fe50be10fa30466417c2b210724fd6e0f1aa6778b
a159590d0954b22e97884cc7e6e21062e94229651ec047071650238caf464b6b46a386674c25707319e6012fa5e10211f1d31d87154b1f
0b6be601e56fdccb1ce83b4a38b1abace1b444f64527f11c76fe2d858671cea2e35e21969560ae6d9dc319ad8bfe22a95d6ad329efe7ee
87e6a175dccc8815cc2143872433e7a0ec2494ef16f54fee1e4d28a79f6165ef4d290c0a2b7bdca8f6d377a5d4aeb3b41be756a35a5a77
688a4c1cb02c99a854a33b06248888801ec0bf1a9c0a56260348a4b15328c4a2a5507e5b542e845046a0b57e450fe00d6abf5a7cdf7b4e
cd85957962ed64829e2951402c0ba044be356a74fb784e02790be55df3a09226b72a67b5faa341c84d7766d1acd496d20ae2412c733199
ce1b85542ceb5f230b38238b042d2f6b1cff48093ae6b6e4144e69ebf431679948c3b9af53ad3532e134a4624ce1216b995358d797022b
b1b86e3d2710cb04c89aac1d9202874d672889c69fcc5721eac16f4d6e289d24470f86e11198a114cd41d514ca9770d8c8c0049a57c868
758d86e5e2057359488ffe8ad7c4d32869b8f2b74b60596ea094fa5279150f2c49e1541851dc931c03eb714eb12d47f41cdf092ccb0130
1600fda8b584f21496e418d0f0f91ee3c57e1e305f2667a0f890b5b560593aa89c3ce0d07357a3432b75b852ff8974fcf82c3f4e2bcf77
8e9ae72c0562b9a2f2b0fb36e223c787b9bebdb8b5d03e213bba586218bea0e2b07b76de4028770005b67315d0bb4373a2e2e910f7d862
5d7f34b02cf92d852587c97b785418fb8e98e99990d5c765e985ceb1e0b9a4851887b72c29a5570da13c8150ee9f8595795678b3d746b0
30398a8fc5ae9d97b662edd09665a5a4a628dc7f50981c83dd2dcc08abaf288c28d2aa14e7083d7259891a8e1cd1f173a03e4cc3f2ae42
11296659f1c1093196e29ca1871c86d370875328cfc98103a13c388b617909d7e8652e99b92e7894602e3274c57026d1197a38cbb2c26a
0c5193d0400e0ca3972e719805e9e5ce2971c692a1d7832f39945856c83178a8b2a8201d8620f6d75aeba665167ae6e2942a94d39186e1
cc1d601e7643284110b20c4b96f7fdd121272667029168240be57414b16416cad758ae08522047c8d38246fbd8325b1109fc8d8641e1af
a50be574846138b350c2db0db2290d2dea252815eb860f95e07ad79625b350de80508212283631dbc2ec3597b82a19f1942100ff291d6f
a84a00bbb52c393bd6084304300e8516a6c8450f34afea123f33afffc51437b3f60cc62ec592313c08e9d440150a4b57b025e205f1ec4e
2ca9de3047194d2c5b0455299c0b44d85a6376356749b5423884f21d8412d486bce4b9f37f3592bf8000bb114b2a519bbc5ac0016ae280
66d04b392b4467802591bb6217c370a594e9709fad1de9602e0874a1c02179b207e7c9080c2f968cc9572194a01b25fd18d11a6dd8c330
1c420986879c352739f7314a598ad1195a2c993ac93b082590000da773d6655fa38cffa022c30ec32935fd1fd68e34e0cc01e2c88d13c6
70bc2e438a259343077194402c990e1f14c8abc8a862597ad15899034453601020d94b2586134b8ea58c18ae8011c8cdb48efe5d87a11c
3cb49e1642090e41ee341147b95a60338c65c9144f89f5b460289452660ef26bea35c328e06724cbb254284f20946034a8cae169ea6523
f6929f21c492c2844a7887256160543293e45e23ab1430217e18ce30fc86e71b0c4fae771cc3713e46b02c8b86df104ab007281c28393b
5187ca90bb45b45852dab5126ef7bd0300f8c87cf17fb4b6802ca45b967f595be2794793e300ec89d7a9f742357f402162e72c73037267
305703f64ace0a363c0fe588b42cc9a903a104c04d722a375897e588b42c0bd77e3fd75a97861a01201a5897ed116759167aefce2194e0
20243b2fe1192f439c65596255e2cd098e04accbb688b22c0be7555e5a5b00d83739d62556f56422cab2845509401a19cf0c56b46522c6
b22cb42a117c0e8e4aaa67bc28c5e19111635916589548a50f0e4dc6b383da531988b02c4bac4a082500d3f3c426f8ddda023611615916
589578430290f70ca1564f22ddc5323713f404a70e0097642c0f86a3271109c3f02ca144a8100057480d3887a32711099665d605c0aa04
e02a19cfd253ad358a9b45d2d5b2544ae5967a805509804d6afab637d616e0a5ab6509ab12005e529f293c4bf174b32c95528fac8d71c0
aa048009a46e8ba79b6509ab12007ea8144b4a850178c523194d2c4f33cb820270183014af439761786e1d70082500fc2013511cbde62c
ffb0b66c935c06148083f221f1b691303b822ec3f0cc2138966701108152eaee344d9f53da0a43f16d9a5b9605437008250011e059a943
8f6178ce10fc9db50500c00659a32080c852b86b905908806452e72d5f585bc0159a8a6541203a00208dd4e9ae87d6167085d69665ce7c
656a6253000e8fd63a37ef02f0d0d41b8ed29d00b403c1e9bc0c31670900a80fd68987692696993f446aca2900c02f529d3c10cb00cd86
e14aa97fa669ba69ed0880610100f990819294b312cf9c9f966289f94a001a83794b3e24cf599e595b0000a0134dc4327375c0336b0b00
0074a29565993c718c383100daa394421a440fadc432673d3800a03d104b0f52e72c315f09000fa9cf12c4d28354b14432520078f89178
94fbd616708148b144e17700d8c0dc3f13d5c51213c600803dd0c2b2845802d00f58964c4814cb6fd6160000e84c0bb14c9d30c67c2500
401c121d3c183600c0041677f0513d910616f203d0173c833c20f92f00004400b10400fcffeddd3b6f5d4dfbdff7993f1eb84880900202
d8400a91f00b10d5a513f50a4475e948764e25aa73e59b2c5389ec5c1820f90a4416869b0024cb20854820450a03e40692224841b24860
3bc504235d4bf7d29e7d5887995973f87e0042cfb3b76e711fd7faad6b66ae41071c2c818231cfd99fa0074b366e07508ad0c972d3b965
351a680048526a97e17d17fd03588dcb704f421f2cfbbe51cc09039024067880b2912c3d49ed6079efdc02600cfa537a42cd12003ae032
1c003ae06009001d70b004800e385802851ab8d491e6db4b70b004ca35e460c95ce725421f2c79e181e91c0df8cd7c6797205902e5da18
f0ccae9c5bf053e88365df7993ac360026648c61aef312410f96c61856e40013d05a6ff1bafb95da6538c912f06348bd122b50b304caf4
65c0b3ba736ec16fa91d2c59f40f4ce794d77e39b6c2050aa3b5b63b143cf77d567cf756e3321c28cf31efa97f3192e59352eaad73c712
9cdd8071fa5ecd35f8eead162359b2220088eb6cc06f3b716ec11f6224cb3da5d477e78ee5de303116f0436b7dda65649c54b95ef06469
8ce9bb7ceac0b905c0d0efdf913d10cac17048e28448718067cfb905c0687307cecbd6bff7caabbb5ef0cb70c5f4212069b6efa53186b1
85353858024007b12ec3af9d5b002023b10e9617ce2d2b0c6c870f00c144b90c57fd2fc5af8d310cf4004846aa074bea960092c2da7000
e820e6c1b2d784583a3d034849cccbf0be6da32e8d31ace601908468074b45d221a9ac0000200049444154dd1240c6a859024007b10f96
5f9d5b56d05a73190e2009512fc31597e20032c56538007430c5c1b2d73a712ec501a4608acbf0de3bcf71290e606ad10f968aba25800c
4d55b37c706e5941f6110180c94c952cb9140790954992e590dd1b592b0e604a534e1dba746e598d3d42004c6692cbf006033d007291d5
a474ad75afed2900c097a993a5dd3ae2bb73c70aa44b00539834591a63ae9c1bd760450f8029a47019de77a0e7dcb90500029bf432bcd1
77a04729f5d118c3e83880685239583e29a5de3a77ac40ed12404ca98c86ef38b7aca1b5de5dfd3700c09f2492a51a76294eba04104d4a
f32cb79d5bd660641c402cc9244b45ba0490b0d456f07c746e5983f66d0062482a592ad2258044a5b8367c48edf2c9b911003c4a2e59aa
81e95229f566489f4c00e822d5ae436f9c5bd6ebd5791d00fa48f260290971e6dcb186d6fa78f5df00806192bc0c6f30d8032015a937ff
bd766e59436b4ddd128077491f2c8d317bce8deb6d68ad8f623f5600654bfa325cfdbda4b1770f4b2ec701f894fcc1520d9f4ac4011380
37b96c5836642a111b9c01f0268b83a54c25ea3dd8a394dad75a6f39b702404f595c8637b81c073095acf60d1fb26e5c319d088007591d
2c8d31b661c69d73c77a763a111b9c01182cabcbf0c6d0cb71a5d4a13186411f00bde57ab0dc1cd138635b122a0074965bcdf227191d3f
71eee8e631f6e30590bf2c936543066e369c3b3a60841c401f591f2cd5b8fa25074c009d6579193e67d07422c59422003d647fb094c19a
33e78e6e36386002e822fbcbf0866c5af6d6b9a39b5763cce6040f1b40268a3958aa91f54b0e98005629a166f9dbc8011b2ec9012c55d4
c1520c6ae72638600258a8b883a54c58ffecdcd1ddc6c8cb7900052a3159da03e6d58811f29f386002682b6a80679ed6da1e343f3977f4
f346d22a808a15992c1bb23be48373473fcf5aebdd091e3e8084147db054bf0e983b765a9073473f37ece703d4ade8cbf0b6314d375a98
8b0954aaf864d99083dcd884f973a45cfa6902a84835074be5ef80a9a48e79ecdc0aa058d55c86b779ba24575c9603f5a82a59363c26cc
e6b27ccfb9074051aa4c960dadf5bd52ea9d73c730a44ca0605526cb864c2bba74ee18a64999a7133c150081559d2c1b3258f39773c738
1f8d31ec550e148283a590553a37ce1de3b15c12280007cb9691fb91af620793b6386802f9e260b940c08e43f6a0b923fb0601c848d503
3ccb48c7f5bb25778f61e7763e32dd08c80f07cb258c31bb239b08aff35d0e9a34e80032c0657807111b011f1a63387802092259762097
e5a33aaf77742e69d3fe1c45f87d003a2259f61070b47ca591bb5602f08064d9839dfa2307ae936c1e34002f48962368aded14a0b7817f
cd832ccb04302192e508c6982da5d476e05fc35a732001244b4f422d97a45e09a48164e9896d9a2107b6c3229e10803f70b0f4ccce9394
83a68f09ed3e1a1403f080836520c6982b39686e8f38e851af041241cd322259dab8dff53752af04d2c1c172025a6b3b8a7ebb6eda1107
4b201d5c864fc0b668b3d38e5ab5cd59752f0290190e961393dae656abbe792d8fc8d7de40003ce0323c03d2fbd2d63b9fe4cf2b1a0803
00629273d18152ea9304fa237600c13caec2332057ea9bb244dc8e0b5db53ab9d99f7bdbd14d9a22010030980d9076f296d6faa57dae91
5d91f66cd5d8187340a8c422542c332641f274c5cc4d3b547425154e0e000080dfa402d9fc6c2c7865665295bc72ee019620581644f6a1
385e7280686b02e72d43ea0050b60e01b2cd6e6b7cc0b90143112c0ba5b5de916ae6871ecff081399c00901f19c1b2c171b763806cd826
4fa7c69863e71e6000826525b4d6f6a071d4e360d3366b0da9df3af70200a2901e4ccd229a77037f2755490443b0ac901c988e57cccdec
83617500f0ac357cbdeb61ff605b1c38b67b7838f700be3fbb044bc89694c73d87cdd799352bd899f80d002e99b2d404c8a1d5c7455e65
2ad4290b37111bc1120eadf58104cdb157c9cbcc5a554e4227806249786ce63dfabc789f77295549468e30298225d68a1034db7e573a25
7872b50d2069012b8fab1024912482257a8b1c34db5e2574de4ae8bc77fe06007836b7e2dac79cc7210892c802c112a349d03c083cccd3
c56c2e78720006b09604c7ddd64facaae332af326a439044760896f04e86858e3cad3af7c906cf7b2a9e405da413c64e42c171de4c16db
5c30fd07b923582238a9061c48d89c6208a98f07099ef7543d81f4c9f1a51d1a7706f6eb8de95a566cd31718c521586212095735bb7890
e0f9f387930310861c27da3f534fb719622673d2afa846a206044b24431a02dbcae6a742de95592b80daf0f9440514b56b55189b9fad4c
03e322afad216dbeeba812c112494b68615068afad10fad4aa8652e140f2649385cdb9b098da3c46dfec77f64286b40991802058223bad
adce721c46f7c99eccfe4fa5d4ffac947a9130fac2a2240c258b5cda3f3b12184bbfb05b672621924a24b006c112459013e29104ced417
08a56026c1b4f96982a9a2529a2fa91caa56206c42e2660515445faea5d50f732281010896285aabbab997c14ad19c354155b5c26aa3bd
b8a9ea8aaa2c46d96cddb4dbfadf4d085404c128a8420201102c519dd62e1a7b052d1402b0980d90ffc17ef59552ff835c60dadb8e8c31
570bff0b0083fd132f1d6a23c35b4fadea90921642ef95529f65ebb4191f0c202bf63bfc5529b56d8cd1ad9f2d63ccff688cf957c6984d
7bbb4c9bb9d05a1bf9b9900b4e002351b1443564fed9456b0ea65dd5b9d7a50fa5cce16cf60aa6ca094c6326f31f6f7d551bb5d6b6c7e4
5fad9bec1ccb03e65702c3102c5134a942dccecd57f37ae22074025e357bfe3701324ac093f9afb77373b14f8c31c7ce5f06b014c11245
92fe97e773cfedd2187310f3f94ab0dd4d788f62600ad7a9eed9bfe462f4412e4669e505ac41b04451ec5ca905fd2daf8d317b293ecf56
f06cf63aaebd5f20f2376bed36955c70ec4a4622eee72a9876facc2e0113588e6089222c0994f624b095fb5ca905e17387d6499848334c
5dcd3ef94b463f0898c012044b646dc9415fd53837aab56bca6e45dbeac18fbbf676a23504c6beb4d6f70bbe4f3660eed00713f81bc112
595a324cd5784f25613d59acb0d50aa15b54438bd034abbf6deda8f444f8194f6b6ddb147d5bf00f253bdd06888d6089ec68ad4f95525f
163cee99540f681312900ccdffef4aa97fbee0b7bc124c477b9040b870bb4df959b46de99931e6c8b9155e49dbb29b25ff2617b5a81ec1
12d99040b3eca43ab38d909d5b118cd6fa76c962a34101bf3594df98ffff4ab63adc71fee33fef8f35fc7fd7daae727ee8d8fb7eeb6b02
0d6d712292cfeae392df48c047d50896c8820cdbfe58f2580995135932efacb13dd5f0eb82a6d721440b736b9e0fa172026bc2e5833166
d50510502cb67444f2b4d67b2b42e5eb9a0a16029293e7eb92dff02855368c201d0f9685ca4b42e534e4a2e9e3925ffe4e6bfdc23691a8
11c112499355dfdf573cc65de6544e6e5578bc91f7100368adaf16b4d16acc6237fcc79f64f5fc9973c72f76aef113e112b52158225952
ed5ad44aa871c644f9e9c97bb0ece46a9d53b9ec4f2a95abb60825542640e653ce963c928dd6c22ba00a044b2449e62f2d5ba8d0600830
111d162bdcc87b8a0ea4adcdb24aa592f636f49a4cc7aacfff5b59e8065481608954ad3b105f32049e9c55554bebcab9050e59a8b6a857
62dba9730b26638cb95a51b5b43ec8c502503c822592237d2a17b5146abb706ec1d4d605c7779c5c3b591bc0a9562669dd31e91b557bd4
806089a4c8817751f3f33f70624d4fc7f7e41b8b199693e0bdeea2eaceb90529e8f2f9a7d28ce2112c911a2a5a79eb127a788f97eb326f
988baa7c7d92a90e40b1089648cdda6a25159bec112c17907ead6c87593e56f3a368044b24a3474b9a45db08221f1bb41f5a686fd18d28
0ec1124523582225848dfc750dfdbcd7aeaeaf09af5dde36180e47c908964849e7132615aff4f47c4f78ff5ceb16ed34a8d8a7a9cf679a
056c2816c1b212765b3dadb5dd5eccacf9b9c864d52ec349e9212c0ed4f742890bab24f53926f1fea15804cb82d980a8b5beb74152b646
ec5211b1bb7d3c4bf84c39bcedd3b626397fd5fe0244c402a884c8c2abae15e7a0ec63b1bd80ed6e3fad9f2bdbca8a3e9a8881605928e9
87f76c9b528f7886768fe794f7b9654bc74448537b0cd7f77bf689aa6552fa7efebded1a66c36253409022c277e9aef1a1f5f34976737a
94bff79278e10019235816480e18ebb684ebcaee9692ea567c5f98043fbdae4dede7a47cc112ddc0ed49d97d2a015aebe301d54a2f9f7f
b9f07f1c5040d890c281e10205be112ccbe47b982c56756448e3e75b86c42737e47d7b726e41dffeac6fed9c68e7564423c7c5de5340c6
ee1c26d39c5e468e48356ea85ec227826561e440e7e36033852107db0daa5fd3b173b806ce2d4bb50a3ea5219fff7da9982132a9d4df0c
f8ad3e367838f5dc4cffdcb905188860599e104d965f63eccd2dbf63e6dcb19eaddc50018b4c42e590d63777c618de2fd7d0b0fd17e132
2e09958f037fa98f63a9f74538322f1f188d60599ebe73ddba8879c019ba08e4ad4c4867583c0209f243fb29327cbb8031e67e4435eb2f
86c5e39051a1a1a1d2becfa95e04301c0e2f08960509344fc65697a29db08c313658be3a7774b321ad92988c1e885d2c252b4f87b65699
c5fc3c65684ce8d8975eb55c5c052295e121c3df8d33e79661425cecbfe3d8091f08966509112ca7b88a1dfb3b6fa8def8272d857e8cfc
87196e5b41a6835c2fff1b6bbd958b2bf61df7a8e9093cb257ab9d52e4e5f32fd5ed4be78ef1a85a6234826521a4ed8eefaddecea6980b
678cb9f270d0dc97561ab4231a49aa942f1ea6595cca7b8b158c317b23aaf68defd2db90eae5489e7a022bdfa1cd18132204ee3bb7003d
112ccbe1fb20e3edea7a0839680e59c833ef07732f876955697e785881fa10e844582a1f4392efa47a49f57e003b2c2cd33e7cf4043e09
745175e8dc3212ad873016c1b21cbe17eda47070d9f150b951adb997f75430d76b054a1f551a251729bcee3dc850a7afd0d054ef09981d
c896882f23e752b65d865ab023f3951f9c3bc62158621482650102cca7ba4e61c8527623d9f2142e9584a4a682c91cb43952a179f11828
95bc77ec4f3c8084069f15a92660deb367b4cb2ecc696d89e8ab47e465844abdef7fff0317e01883605906df43d6c92cb008102e959c34
becb49b6ea9d7ba43a792527d41bcf4d97ed5486ad81db15224cb85472d1d0ec195d75ff4b993ffc249fff310b731689112a9beab6afd5
e60daa96188c609939a93cf85cb473925af36a1b4c8c319b01867c94bc76cf7292bdaaa1922361f2424ea6b63af9c9f94be3d93655844a
0f245cbe0ff4cfff259ffd975a42a654e6efe5f3ff6344ebac55bec69c532cf3e17d5e7c87e8878c4a102cf3e7f3e0354bb879af92797a
beafccdb3eb52a394f254d629779634fad301972f5a7bd38a11f9e4752957ae36941db221bad90d954f28b180e950ba96309cf4d653ee4
b6b7efa51f6f6c5e479a58c483a10896f9f3793049bec7a05c997f74eef0cf5631ce5b27da273939253f6c2e27d2a32648b6e68d85a8cc
b4bdca4995ed050390cafd56e08babc607998fdc0e9a59040d3bea2015f997d685d45f9ea7792c623b1f68b908884e2adb3ef6216f102c
3188fd12f0ca654a16a07cf7f4e8afa57f5e36ecd075a061dcae5e657f67dbd4fa2af6b0af84dc5dd91f7e2fc289739528f3c9ba92615d
df73e6e69d4c15a265cac6fdc4eff94c3eff57d2d83d2a790df6e43b30e571c0fa9cc282c7917b982ff27eaaa08c7c112c3366ab081ee7
57bec9713e9c0cd7dd4e7c825dc556105e2404347fb635f359e7e776dae7b5d9fad377f37b5f5e535ca0537ab06c4815f1dcb9230daff2
796f3efb4fadcf7be35e3efbed9180e673af24346e061eba1e23a90b2af5f70e59bee6484efe19477e089699f27c65fa75a23941de78ae
de623d1b1a7653ad66d4122c1bd2a3925d53e27990cf7f9217e3d236ccc7c5f6ab2c9c043a638e65be7c5d25cf720f954ab681b4f39b42
ec44813f34f3283719224b87ad9ac9e73fc4fed1f8db838ceeec243ec2e3ebfcb041cf5ff445b0cc97af03475113b4ed047639c1bef7dc
7ea376330265fa5a01930b2cbf2e65614eea81f22799ef79eddc314cf28b3a91168265866c1f364f2b7c2fa798741f830d3f3284f3862a
ce28677242ddca2c50563d7c37778115a2ff6b0d5e65518ece749f7b5f81f0033b35a10f82659e7c1d308abf1295162d079c647bb10b8e
b6e5849add67442ebc68f0fcf705d68e7cfe3f53c5efe4443efb9b29acf41e4a36ba38f1f4cfd17a089d112c33232d667cb4d6f85adbae
287327d96dcf3ddf7277dd0a93bba9edbed49584ca9b3c1e6d5c320f79b3153243355bcfcd6b2b4cea925641cb73f1f13e331c8ece0896
f9f1f1057f2861c1ce1836384980d272a2fd5a5935c79e6c0e5b27d3bd5cc3648350d99d84ccadd645566dd345ec45e5c75665b2e4963a
3ece192ce2416704cbfcf818926058638e0dda4d35a7b5f8a1a461f3bbb920b9253b751441fa39122a07908bac83d6677f5b86504bbad0
ba94c567ba55952f727ef93c8f0b79386fa013fa5866c4534526b986beb990f0b297c02e1fabbcb67643c9767e581f133609afa679746b
979b83849b95abd66e4017742ff89b4ca17a76eee82fcb8d341017c132233e9a204b45021ec94977477609d909bc4b4eb39bc96df367cd
07fa488dd097e122edefddaf9acffe4ee0e0396bede6735b4bd5d107adb51d12ff36f29f62271eac45b0cc88d67aec9b7558d2f027ea36
71a86c7c24dc20175aeba791adeaec861ab41ec24accb1cc845c6d8ef140a84429a47a3f75a8b46e64281ec8c1d8cfea5b9992052c45b0
ccc7d80302273f1421c17db1cfa57a0a244daaeb633b00d07a082b31149e010f8b76ce726c740dccd35a5f25bc788aef1992e769210f8b
78b01415cb3c8ca936be72b24309b4d6b789afc8ff22d55420591208bf8e7c7c9c53b01415cb0c8c5cb4c3821d644f4265c8d5f63eddd9
3e897cea9032adf5fd8815fc76cefe8e732baaa7a858a66fe4a29d3b422572272b59730995d607396903291b3312f68e453c58866099be
31db68b1600759f3d01e652aefe4b103499206f26316f2707ec1420c85274c1a0fff18f8085948806cc902031bcc36327f176d43fb2d16
3a205563a65ab1e10616a16299b6a157842cd841b60a0a954a9ec393ecce04a4e870e863f2d05f1905a26299b01157929f6bd9271a6591
00f658e8dbba6d8c61781cc919b19087453c7050b14cd488dd3cae0995c851e1a1d27a64c1031235f47cc3221e380896e91afa45676802
d991f9c42587ca065b402239b290e76ce0e31ab3c0140562283c412316ed9c1863d85a0e59f1b0b3548ee82f8be468ad5f86cc6d66110f
daa858a66948456346a8446e2a0d958afdc591a841235e54e1d146c5324103af1a59b083acc8c9e8bcf277edd218c34919c918b8cb15bb
4de1372a9689d15aef0d08952cd841560895bfed6badf9ee2225432e743ec8142e806099a0214311543c900d19022654feed93548980c9
494bac210b79380fe12786c21332b0ddca5763cca9732b90e667dc86cabf9c3ba0649e348dd491842153b258c40345c532397daff86684
4ae4426b7d4aa85ce9addd5f5c761e02a6d6bb02c9221e282a9669197085f8d118c3101a92a7b5b6ad75f679a73a617f712441e6ff7eea
f15858c4032a96a918b068e79250898cb040a53b7b1ce0bb8d14f49df3ff817df141b04c47df2fb05d4d6ae885871cd8ae0532ffea9237
6c255bad7ccffecb48812ce439e9f950180eaf1cc1327f7f49c07cd15ab39d2392263d1bb72540e14f76379e4dd95e0f989cccf7ed3bb4
cd79a87204cb74ec8d3cd9dae1b36fad90c9fead4892ad82d800653b1af00efd64a7b568b678440a6c98b4732bedb94429f53ca059fa86
eca8854ab1782731728578e471f5eccc0e4d301f132992cfbbfd6cbeabf00db2dfcd1d16e9606af23db4d3aabe787828b4cdaa1c15cbc4
d8938cddf35b2a18764eda47bbd26ec4a37c6bf762964ae63dbb232025f279b79fc9f795bd31b6a3032bbf31293b47bf5599f4112a0f09
95205826ce561a6dfb8656d0fc3a62c8dc56857eb44226070024c1ce2bac6471cf997c971941c024ec5c7c992e653c8c8cddc94592663a
071a0c85674c82e1b187fe80f6e0b047f504299061b9a7017be6a7ecc12e82e03b8629c89cfb0b0fdf295bd43866630eac42c53263b208
e2a055cdfc2cf3b6fab293b39fa59279c1ce1f98920c8fdbcfe061016fc4eff641844ac46417d0c84e4eb67af47d44a8b4a308db729ed9
2454621d2a9685f2b408c80edbd13a0293d25adf0e58999a824386061193cca1bf18b918ce16278e6cef59e71ea003826525e480733ae2
047d62171539b70211c8e7f73693e1f16b630cedbe10854c89ba1879f1659ba09f5255870f04cb4a69ad0f2468f63d5133c70693d15a9f
7a5abd1a027b7c230a1991bae8b98f77db9d1cc7594406ef089618b308e8557a643264826812ed7df999ef01421ad96b928200a22158c2
21bb265c480fccae68c48ea864a5ebf7895ff54bd9a61208c2f69a1c3857fe52c2e493730f1010c1122bb516011df518367f9090c99ec7
08ce6e3f376248708cffc918f3af7987e19bed3529d5c93e5395b8b84712683784955a3b016d4a4ba3f71d76026a3762bfa5113b429285
32ef47eeb53fc47fe28d852fb602df6a5cfead43a87c9545376fa415d016a112292058a217d921a5bd13d0e19ade9976a5e2a384cc2b7a
642204f95c6eca8916c882f49a7ce9d16bf24efaa2363d258f592c86d4102c318aedd32757ca3664be59b325dfa75623f65342267c9396
58f673f87ff1e22245b67556ab71f9cd9a30f92afd509b2d1377996284d4112ce18d0c9bb77702fa28f32d17f9d20a99f4c7843752c1f9
b7bca248859d0ea4b5be9730f963cdc2c8f99d6e68b28fac102c118c9def235bd935d5cc9325f3e0fe9280f92293d601206b764446e698
db30f9b8a23d96bdf8fed8aa4a1eb0921b392358228a058b80b6172c02b24342df5a2193dd4b006443c2e49584c9e725bbe1cc2fbad961
d10d4a42b0c424ec15f99a454036647e6f85cc5dde2900a9913079da0a938b5a5fddcd0d6fb3e806c5225822090b16019db51e970d9937
1232ef65df6800988c9d1bde0a93f3bbe1cc162cba61781b5520582239326c7e34b708a819366ff7c8bca747268058ec1cf0567ba0f9dd
702ee77a4ab2e8065522582279b208a83d6cde2c027ad7ea91794bfb2200bead685cbe68d10dc3dba81ec112d959b008e85a26c93fd388
1dc058d2b8fc69ae71f9cf4537ad20c9a21b60018225b2268b80f6e61601edb41bb1f30e0358471a97dfb71a97bf958bd63f16ddacf967
80ea112c5194058b8094044c1ab103f883342ebf6d352edf9c5b74b3c7a21ba01f82258ab56011d0ad9c4468c40e546aaed7a46d5c7ecf
a21bc01f8225aad15a046487b44e6585e7138dd881b2cdf59abc524a9db6aa92472cba01fc2158a25a365c4a85e24a86c48e69c40e9443
bed336341eb5462f7659740384f30f5e5be0d72220a514733081ccc93417bb80ef5816dbf0bd062222580200b26787bb655e359d208009
31140e00c81ef3248134102c010000e005c1120000005e102c010000e005c1120000005e102c010000e005c1120000005e102c010000e0
05c1120000005e102c010000e005c1120000005e102c010000e005c1120000005e102c010000e005c1120000005e102c010000e005c112
0000005e102c010000e005c1120000005e102c010000e005c1120000005e102c010000e005c1120000005e102c010000e005c112000000
5e102c010000e005c11245d05aff1f5aeb57adf526ef280000d32058226b5aeb2dadf57f564afd774aa9ff4629f57f6bad777857010088
8f60896c69ad8f95528f4aa97fd67a0eff504afdd05a1ff0ce020010d73f78bd911b19eebe574abd5df1d0cf6de5d21873e4dc03000082
a06289ac68adf79452cf6b4265e38bd6fad6b91500000441b04436b4d6ff5e29f5bde7e3fda0b57e726e050000de31148ee4d9053a4aa9
ff4d29f55f0f7cac6fb5d62f4aa92d63cc8b732f0000f0828a2592a6b53e92053a4343a5913f37ec10ba845400001000c112c9d25adb05
3adf463e3e3df7ff1fb5d6bbcedf020000a3112c911cbb9a5b6b6d2b8def023db61b59040400003c22582229d29bf24784c7f49d5e9700
00f8c5e21d244386be43552917b1bd2e378d31a70bee0300003d51b1c4e4220c7daff24daaa4000060248225261571e87b95bf08970000
8cc750382633c1d0f72a365c2a630c0113008081a85822ba8987be57a1720900c008044b44250dcfa71efa5e85700900c0400c85231aad
f5adddbb3b8357dc86cb2763cc85730f0000588a6089e06c4b1fa5d4936cab988b73997349b80400a02386c21194ec70f39c59a86c9cb3
fd230000dd112c118cd6da361eff9ef92b7c43b80400a01b86c2114462ad84c6b2e172db18f394f7d30000202c2a96f04a6bbd95682ba1
b1ee65ae280000588260096fb4d6074aa9c7425f513b47f4deb9150000fc46b084175a6bbb7afabcf057f3adb44c0200000b102c319aed
f9a894daafe495fc20211a0000cc215862b0d67ccab795bd8afbb2831000006821586290c2e75376f18d36440000fc896089de2a994fd9
c50d2bc50100f81b7d2cd14b61fd297db0afc756fe4f030080f1a858a213994ff942a874bc65310f0000bf102cb196cc257ccc74bfef18
f665ce2900005523586225adf5b19d4bb8eaefe0a773adf50e2f0500a066044b2c25cdc0ff5a763f1c344f070a2217d6007a2058c26157
3a4bd3f30fce9d586583f996405998e602f443b0c41f6438f7b9c2a6e7be30df1228c7954c73a1f303d011c112bf4920fac12b321a2722
a000c618db4eec4eda8a01e88060899fb4d6a7343df7eaaaa0e702d4ec56a6b9f09d063a2058a259a4f38557c2ab77ec270e14a19937fd
89692ec07a04cb8ac9229d1716e904f38d2171206fc618bb90f1419e04d35c80350896956a2dd2a1e979580c9f01f96b7f8ff94e032bb0
577885b4d67b4aa9efb5bf0e91fc1c1237c69c56f16c8132b57bfaf29dee40766cdb911f5be5dd941f5b01be8af1fa4901a5f9fdbbadbb
b6a4f3c94c1e4fc3becf7614efde18435fe281b43126cb078e61a4e12f4dcfe3db9621354410e9737e628ca181762564da507b84a7baef
b404b5cd56586b82dbd091af3363cce8b9e832f7d5164c3e3977fa63bb031c1338d7a362591159a47864af610000200049444154c37cca
695ccc5d3103c88b1d02df6f3de22b095559496cc4ca2e1aed1d2ce5399c46eeb76ccf9d375aebe6ff7b09c525628e6525d84967721f58
518a3eec2e4e5a6b937b7701dbca4c9e47ee43c7f3732b73edfc90e5e7c986495b35b69f2509c6536fe2f1453ed7f77621ac736fc5180a
2f9c7ce09f58a4938457630c07a0084a180a971368dba53126bb8b93529e875afc5cac37c69817e7d6442d790e9331c6e865bf5b86deed
68cf3be7cef43c1863b2ab608740c5b26032799a95dfe9d828a06a830864a86fdebe5448b20965720c9ab79f7105f3ceb9e5d7906c4e4e
127aac8b5e4ffbb9399200fc239350a9a4826d3fd78b3ef3552158164a4e3e37b5bf0e09fa421f3c74b0eae4742e4382395447563d8f2f
199e88172ddcd8cfe9394895fd8dad1c3b77c6f7fbf594beca171228bf25f0d886baa97dda13c1b240320cc8f68ce9a26a8975d685463b
0af14316e4a5ac4be0b227e25cf6e25ef67a67f59db643f7763a821d865ef4a3947a6f17a7d8e93bce7fecd7ad04ca7b195ddb0ffcfb62
395f32ea5005826561643f5bda09a5ed53cd071d74d275a1dd874286df9a61c4a4bf172b5acdbc2ba94a658cb17d1c8f644ef847e72ff8
732a813297e1ee3ebed7baa887605910b9ea0bd9c70bfe50b5844f3789ce59ecdb89e2bb5c1ca76ce1bc40dbe3d0b9a50012a6cf023d93
b181d2be175f6df85d54799dff914a6ccc290017ce2d15205816a0b5e77789577da57a5bfb3c1c7867e72c3e155025f9247348537d1ecb
9aa2bfcdbd35d40a2984fdbb050172d7eee0b3a292fc07a9c4fe9e021021647eaa714e3dc13273f2a165e5779ea85ac237dbdbef39a193
d9d0397a1bf23c525ca0b46a3e283b31f975273b1c3521d2eb9c62697bb5eddce1577505048265c6645ed563edaf43c6360aae70609c87
91fffd6322f32e5785b02e7e2438ef72d573e23bedc7652b4c2eab107b21ff7ec870595dfb218265a6e4e0453ba1fc51e1c0223e1a6e97
d2f6e47b4acfa343d58ceff470771228a3bedf122e0f9d3b3008c13243d24e28e73e5ff81b150e2ce26bc8ef7ce25056caf398b76c018f
92ef34f3a7fbfb6c2b9453fd7263cc4584f64a5520586686764245a2c281793ee7924dd953cff7f34825b0adab28f39deec72eca496181
5095abb87d235866847642c5a26a893ff85ea420c3c9d1ab41019ec779267347dfb2b55f6787013e2743e5d2a83f6904cb0c483ba127da
09158d608979ab865b87b899a885cfb573cb383709ac7aef1284a85aae772d43d0a908ba50a81604cbc4c901f449da88a05c6fd98d0773
429c70a7a8c894f23cfafa50630fc39e68b956208265c2a487db233d2aab41d512bf05aae4d80b98a81522993be77b51c4c6943bf4f418
ba6511cf0a090d81c3238265a2647ece8fda5f87ca7c605e16e69c38b78cb73f41753cc4b0f0a78917f37409cb5c2c2ec70aec42112c13
448fcaaa51e1405ba8a1c2d855cbd34041e27cc2ad1fbb0cc76f70b1b8140b650a45b04c0c3d2aabb75fc05ecff0c418f312a86ab9117b
483ce06296d4e7e951b5445508960991033d3d2ac18908bf19638e0355fbf663eec52d55cb9973c778519f474bd7f981b4884355089689
d05adb83d47eedaf037e627538e6859a2211bb254ea8cf76d2554b3a3ea02604cb04488fca0fb5bf0ef8ed1df3b2d0262bab2f03bc289f
22572dedbcba33e78ef1a658f8d667453373a7510d82e584a4f1f90b3d2ab100150efcc11873106828396ae831c61c057a1e294f216138
1cd520584e4416683cd1a3124b7c597c332a17a22a37c5672dc4f3483abc310a815a102c2720bb313c132ab10a2722cc33c6d88bd143e7
8e91627fd6e4797c74ee18ff3c521e72e6fb8c2a102c239303f863554f1a43b13a1c0ed991c7770ba2e8814c765df11d92a33d8f01bbc6
102c5105826544b23290c6e7e88a795958485a10f90c9793849e002139e545902cd04415089691c810cdf72a9e2cbc61381ccb780e9793
2d20f41d9227ea69d909df67d480601981eca6735efc134508ac0ec7523e43d994a1c773b84c79e72a82258a47b00c4c6b7dca6e3a1881
fe77584942d9d7557f2707f23c7cccb94c39bc112c513c826540b245236d6330c606c3675847b64b1cbbca7af2cf99ccb97cefdc510ee6
59a27804cb40d8a2111e112cb196ac527e13685ff16864779e37819aa84f8e0b45948e601980844aae4ce10b272274628c7931c6d83986
d739bf62f23cb6026dff08202082a567ecfb8d00f83ca11763ccde80a1f1be7d198393ed1ffb0e8d3f39b7a4850b45148d60e991844af6
fd86770c9fa12f3b346e8cd14aa9bb9c5f3c3b342ecfa36b159660094c8860e981ddf75b6bfd42a844409c8c3088316657aa7e2be75e0e
d849262aa9c26ee7fe3c3215a237e807d9de3825211e4f8acf332882e5483654ca1532fb7e232482250693aadfe68a763e59ccc9b47b8c
cbf358363c1ead3a3b621421f9a92d3608d9fecb76144e6b6d9452df9cbfe4c7a3fdf7b5d657b2335dece7698b4207765d843ccf50fda6
9be7792bbf2fe55eaba3112c4790ab906742252248763711e4c3b6f39161e5f980799ad9f3b86f3d8f7605f3c2f9cb58c8ee50645be269
adef25f4fcfeb14148fa2fc71a85b3dbd77e9f7f1cf2d82ec64e05923077d504e5d6f37c9630192bec7f90dff7bce0b9dad0795a42e8d4
c618e7467478e17e85ca47e70e209c3776b52cafef7ab2db55e88d094ea4a977b664fb431b36b3be7091e771252bc963fd4e1b766e9c3b
baf938d5907de6a36cdbb66aeddcba84d6fa2860b535a43b99c292252a9603c801855089d8a85ac22ba9fc65ffb992e751d53cb6118e32
1e65eb5b91ce31542a999749b0acc5c8ab54600ce65902e9c8f5fb98f3a847df2a6fce1d11b27d9f08963d102a31312a324019a60ca539
cf43ed1b2c73ed10f0203b5065e91fb93ef0d8ece4df802bc6802e0896403ab2ac58da79da5a6bbbaadeceb5bccf61de76332c3c605eea
a984cb971c829aacddd8caa017eb4a04cb0e089548043bf000182db76ad8d0854e129ab3a95acac2a4ac43a562287c3d42250060012ef4
800508962b48cb12422592c1d68e0080943114be44a43e780080cc70811786bcae5b733f0d3b27f49d34c49f1fcabf9755d43f8792d9d6
735a04cb05089548d86ec62b1d815214bd255f28d29cfd407ede0dfc351b0ba621fcf1ffb5d6ce7fa4949ad926fab22140b62bae73c050
f81c422500600d362be840b6527c9adb42f1db88503986dd9ef28b52eac7dc568a17b21a1b9e102c5b089500800e180a5f40f61fbf6d05
c9f388fb8d0fb56f77d26b05cd22f6eb9e12c152d8ab16422532c0090d981e154b61e7453655495b0d5c304c9d1b5bd57c96907945c8ec
8f60f977a8dc77ee0000e0cff3c566c67b6d7b615f03095d4676a34bbd2a39d4a756c83cc8f329c4577db0245402007aa8b65a2943dd2f
3257f293f317c6b1abbd2f9552874aa98fc618bdecc7deaf94faaa94ba8ef8f4cf9b3999ce3df843d5c192500900e8c9c77494ac3a3bb4
02e50f8fd5dabbb900b9698c3930c65cac6b1764ef37c69c1a63f69aff5e29f5462975e6fc65fff60998ab551b2c0995008001aa99e72c
43def71e03a5ad4abe9730b8ebb3dfa4ddbed118732421735b29f5e0fc25bff619225facca6049a804000c54c5568ed225e5d9536ba026
506ec6e82169f7dc36c6ec48153374c0b443e4f72cf2f95b75c1925089ccb11a159888af1d7752de1946aa944f1ebba49cc50a94f3a48a
b9237332437a278b7caaefdaa16a0b96844a14801d2380e9ec95fcdadbb99452a5f4b5cafbd00e4f3bb7462673327584eae50d43e31505
4b42250060241f15a93be7960448b5ed87c747622b95492d7091eae5a573875fe7b587cb2a8225a1120030866cfb37c55684c149a8bcf1
f87b6629542a17b12bcf099761151f2c099500000f7c0d832735bf5286bf7d864a2bc950d988182eabdc83bce86049a804007852dcc20c
59c9ec3be8be1a63ae9c5b1323e17216f85165d5afd497628325a11200e091af9d66520a1b1701b6a7cca97178e88b85b7d2b6a92a4506
4b42250af6c49b0bc4e5b38d4c2aad86b4d67b01b665543975aeb0fd2e23ecd6e3ab6d53368a0b96727540a844a90896407cbe82e5ab73
cb7442cd83ccea1815639191d63ae939a7be15152c2554567775000008cad7c29d24aa79b260a78a1d843a3a09fcef17ddff745e31c192
50894a5439191c988a84305f6d8652f9fe060b3a29ef2ab4c2e9f2bbbca82ac417112c09950080407c86b0544257b0452b396e6b68b77e
0cddb8bea6ed1eb30f9632778150895ab0a5231097cf40c01ce974311ae449d6c1523adb7f73ee000a2557d600229006d7be86315f6515
720a420ecde65a9923587a926db0945079eedc01942bc93d868182f9dc962fa5d18690c7922c779bc9746e6892b20c96844a548a6a6577
556ea506ef4a9c5f195ace2ba0b978f720bb602913600995a811f32bbb23586214cfabc15562c132e490fc464d0b55baaaa9229a55b094
0fabefcdf2815c305403c4e3b5f29658b0087d915a55dfc60eaaaa8466132c651235a11235a36209c4e373b794d48245e8fdbcbf38b7d4
edaaa6679f45b09450f9e8dc01d463c68a70200e191ddbf0f8cb920a1691fa36866e3a9e0d634c55af45f2c1925009fc44b51288c7f7de
ce2956ac8e9d5bfcfa22e7ef9c8468c3f4d5b9a57049074badf5262754e027be07e961ce6b81e4bcf3c9e333bb4ba87fe56f32e7f3dab9
c3af6c8680032d387aa8ad5aa932a8583e791e8e00724588e96727a7078ba4f8ae56a6fcddb5adfb5e9d5bfd7997d190b8eff75d65dc2c
7e946483a5d69a5009fcf24af3dede387660a81a86c17f92b996a157707f91ded3c992d6523eabd4d676adf3e2930c965a6b7b127debdc
01d4895009442001c8e745891d0a4d7a1a8b5cb41e3a77f8759e78b8f47d8cfd98e2f48758920b96122a43ee630ae4a6aa561563456cce
ccbcd7f25453ad6c33c65c440a97a1170cf516607474bbf611a6a482a5d6fa82500938089609a2fd5359b4d67b9e77da5111fa457a23e1
f263e05ff397d6fa5e16484dca3e06adf58bc7d151db124ed75ca96c24132c6582efbe730750b73b024c6f2cdcc110beab9549ae065f45
2a6d6f6c485af1d7c6b2e1fd596b1d62b14c27322cffecb15279628c611b599144b094f2389dfa0117f32bfb8b510da96a8bb6d2c9f409
dfa365d9542bdbec85ac84a413e74ebfbe69ad4dccb99776918e5429cf9d3b87b12beadf1863921be29fd2e4c1523e547f39770050b99e
9c2656658b0f8ce23d18c8d072b6242cd9eae543e0e7702e01f32ad410b91d11b5bf4329f5c36395f2b33166931125d7a4c152ae127d5d
3900a5c96e282d11312a9654920b21732b7d572b4357fba290eae54ea480f94986c86dc87cb245a72141d3eef66387d9e5df3012287d8e
881eca5c4ae6be2f615f9cc5f784fec5bffa46fd70ee00d0f85ae3ae0d63c98924b41386bfca6017930458b4b35de245a104bdab0a17d9
da21efddd45b47a562928aa5ec1f4aa804562054f617b1d51015cb02c8542cdfa1f2b2d49106a960da80a5a53d51c85d7b527026732837
0995dd450f96ecff0d7472c9cb34489415e1b5f7a92b48888bb72a2e08ed1c52095c36647e2d28645e4a98b423ba47cca1ec6f8a8ae53d
dbad016bb168679818c1b2f42a4d15a41b89ef73d1758d952d3bbad20a99db995d18cf5af326edcf0161729ca8732cd95507e8e44126cc
a3a740f3e5e6d94555ac3ccf984cc77a0cf00cde3364fa2719a5dc933ea1a1bf9bebbcca1cd10b461dc2f947ac5fc4ae3a4067542b0790
13588c131727a4fc8518aeaeb25ab98e54ff2e161dd7644ef4ae8c34ec78da05672623a3f67b7acb7b125f948aa5ecaa430374a003194e
424fd236e67b84d7ed23d58e7c05fc9c14b9121ce82bf81c4bd9b6895009745344ffbb89ecc5f8b584caec3995330f8a5d090ef415b462
19b1820094e20d13c787b10d913d0da5adc2fcca8cc994acfd00cf80ef2d2082cdb19406e8844aa0bb334e4ec3c85cadd0a15231bf325f
52e808112a4ff8defe22dfc366bee496fc2899f37895fb36970dc9373b323f744b76fbb23ff6e2f6aaf61ec4412a963289fed9b903c02a
ccd11a485ac7fc15e15731bf32535aeb9700ed8566c6982de7d6b25eb7262436616a73e442dccf296e8728cf73b315169be73cf433630b
0547cead1508152c630c490125a9f620e443a436432cacca94d6fa4af6a2f62dc990b44cc0d7a18fe0d34912799ed51e2fbc2fde915e95
844aa01ff69d1e482a0d31da0cdd39b72079b26d638890719d59a8dc49216c85263d4a8b7f9e29f31a2ce955090cc2dcca71a2ac0697c6
cac888848cf3408ff8c0b9256d9b957c768b9e9a90036fc152e6388598180d148d21f0d1629de0695c9f9f50f3610f33bc18bcaf613b52
9903cdb6ab13f2122c65a821c6c479a0345f7947878bb81afc8eaa725e649e5d88cfc6758eab9bede7d7eee76de785caee34c592e7f9b1
f4e799aad1c132f0500350b259ed6d293c8855eda55a991119410b31cfeed518136bea4510765ea85dc96e17962cfa91e099fd7c625bb9
ecf03caf9dff10a38d5a152e6d859e02b470006a90d58ad2d4c46c6bc66af07c04de98a39a765301cfefc96d3210aaab04abc287b92754
028364b5a23451b1aa9597ce2d48928ca0850a952735f53095a91fb58ca830cfdda3c1c1525680d356081826b715a52962181cbf4985ed
3ed02b62ab6c35b604ab2248b3e9815f8382252bc08151beb210641c59301863b4e481934e36428da0bdb23f3cd05def6029f35758010e
0cf3c0821d2f62558f78af322073e4428da0ed38b70058aa57b00c3c7f05a841d62b4a5320d5ca18d3705e736c2b531bd9ed2dd4ce4b87
ecdf0ff4d3b762c9901030dc574e525ec4aa56b2cd66e2245486daeded840b0ba0bfcec13260b359a006d70c818fa7b53e8a751ce2fd4a
9b9c934285cacb4a17eb00a3fda3cb3f10b0d92c500b5681fbf12dd2ef39716e413202572aed0a70beafc0406b2b96b2651a8b7580e13e
b30a7c3cad75ac0ae22bd5aa74050e9533568003e3acac584a5fb01be70e005d9dd1087d3cadb55d99fb25d2af2354262ad40e29c286ca
2de75600bdac0c964aa9ffc5b9054057b6b5103b3af8116b1105fbb7274a6bfd14707eed2b6d85003f96064badf57f544afd4be70e005d
31a4e6812cd80955a59ac785406202ee59ddb0a1728be92a801f0be7586aadff0da11218e52327aaf1a4776eac053becdf9e1879ff9f09
95403e9c60295f64564402c37d651b406f62063d5602274476797b0cf888ecb4874d4225e097132c9552ffab730b80ae2e99a3e787ac02
8f3504cefeed099116772177797b60a10e10c6cf6069e7b0d82fb2d6faff534afdb7bcd6c0200ff4bff343aa55b15681b37f7b42a49d50
c81677b64f250b758040fea1b5b6ab2df7798181515e3959f9218b354256abe6b17f7b02222cd251d2fe8b055a4040ff241596192f3230
0ac36afedc47fc5d87ecdf3e3dadf541e0453a4aa63b102a81c09a7643bb8127490325db667e9e1fb2ff7394bdc0651578acfe985822f0
4e3a8d8f2ca803e2f839c752aed80f79cd81de3e52f1f24316eb7c8af4ebecd40586c027647753d25a9bc0a1d2b6137a43a84c4a725386
64eb6a78f27b55b85cb99ff1c2029d1d72c2f2435601c75aaca3d865655a7211f123f083b8a39d5092424e77188a8b4c8ffe683724f34f
ae737f524004270ca3fa21f3eb42ae029ec7bcca8948079297081711763e255528ac25bdbb635ed416cfe96329c3430fb5bf30c00a3654
1e2fbf1b5d49a83c8ff8829d7141300da94a875ea0f32a739e691f95b0c4869e1975f2cc0996ea57b8dc61a538b0d025a1d28f0942e51d
ab82e3b31521a95286ae4a3743df54a3d33779b06c55cf632d16ac8636c62c7cae917a8a0139b9a401ba1f1384ca193badc417b14ff267
f6790f431658f96617cf6d4ef89c6cb0bd71eef0cc18a3633fb7142cac58aa5f2fc88bf4e67b75ee04ea43a8f4648250f9ca629db8ec7b
2c812474a89cc9aa6f42655e366401575452a5bc8f112a6bb634582ac225d020547a22f3ec6287ca2d5606c7d11af68ef11edb053abcb7
e1dd05fa0d5fe4223338696d752f737cdf25f2ba166b65b054122ea564cd9c4bd48850e9890c8bc65cfd6ded103cc26b55821e234c9f6a
aa942cd08923e4e29673392e78271739b75239ff3145a0acb53fe6da60d990f949ac16474d08959e48e88831d7ae6d9b851ce149308855
093aa44a195de82e0afb36fcd963c4d02026173676fac595fc5b462e7242efe8b44e95f3ba972ede5946b65c8bb53b0630155a0a79203d
e2ee27580448a80c4c6b6d57d87f8bf4eb1ea45b09261071115669aa5c34d8b962d9903e9727ce1d4039be122ac7d35aef451a1a6d7b25
5486d55a98132354daf7f33da1725a327253ea74b8cb80abb7df4a05f54a2eb2abd0bb62d99093c677e70e206f8734d01e4f567cc6decd
82853a014db09a9fef6242a405e173414fe98fa94e81da2a7565a7191e97d2dd6070b054d30e7301217c64efef7126ec7f4b9fca402608
94cc6d4e5421fdad175eb04c1c2cdbb21f71e93d14de669fbcac180fd58e0088a1193e25548e200124f4967d8bdc112afdb3f3eae4641b
2b543ec86a6f4265a25a5d62ae337be877f2d9d28b426562b2fffc8f0a960dd9ecffab730790be990c9f32276f0459f51db3aad53893e3
0ffcbc8f9bcdcada888b357ecfa3641a431e64adc576e23dae1fe47365c3e46e469fadeccf45a386c29d7f8ca171e4e58e5032ce04c3a4
6d0b87b4d09f6d20ad94ba8abc6fb20d25f6847fefdc836cc8f0f85502ad7d9454528f86140ab4d64f09ec1b3ee95697be780d960d5a13
20a2ffa294fa67037e1ded84469093c9fd8407e2f70492f122b70c6a10280b25c5a5e388e7ff3b59f4327a1a93e7efc283541eefe5cf9f
ffbb968a7c9060a9febe02fee1dc014cef337b0b0f37f18523fd0c479293ffd5043b9110282b241d6476e567e867ee4142da55c863b75c
301fcb3cc7f6c8eb4cc2e1ad3c8e7ba64f2d172c5836a85e2221afb2c51f078401261ef656329ff2c8b9159dc83eedb1b7d454120af6f8
de0175081e2cd5dfd5cb5be65e6242cca71c2891ef2f43df0348b5e862a2f7ee4e02250b72808a785915be8e3d21c88454568e630a2784
cafe6ca0d45abfc89496a942e583acea24547664f75bb60b116465f7f709debbb30c57e202f0244ac5729eb426893dbf0775a2d2d55342
230ccc85edc88649a5d4e984c755e64f02f8294ac5729e4cbe7f9f780f2ce48d4a574f89542855ab5136a172053bccddaa4cde4c142a2f
e5bddae4bb06404d152cd59fc3e387ce9dc03827ac1ceece2eca4924502aa952d2287b09bb00c7be57ad61ee29da3dbdcaf6a7f6c2ed80
f70a40db2443e18b24b0e214f963d5770f89756c6071d502322de13491e6d3f47e05b05632c1b2417b220c742ddb8c6105092a1709cd71
666e5e4bab8fde17e7ce69d89d4ca84a02e82cb960d92060a2878f3e765e2899d6fa34a1b0d2a87e4b46099247f2934a3bb63b099354fe
910cf9aefc9ce2c4f13e6dc906cb0601132b307cbac2c43d0c57b9b473f356dc5fac09b6bceb6af01ecbc0503282628f057b01e70bcf64
b71c1b466f191d092ff960d92060620ead681690b6331713eee1bdca830c7bd7b15feeaf0acb9e04c914df8f33d96799616e04d7baa8da
4be862f755b6370dba55646db209960d16f9548f2ae59c842b938de203a584c803f949b547ef4caa929c401145e217baabdcc9451743ee
0364172c1b6c135925e652fe392f6f8a7d9ffb282e50b6aa90f6e793f317d243551251c977e42293ef47572c62eb21db60d9900ff12d3b
f914adda39790db9f23f4ea4edcc3ad9074ab970dd950099c36bdeb89420c95c494455c1b9f8ab31e6d4b9158e7f38b764464e5e3f578a
25baf215c355db9732e1451eab6475012061bdf9c9293cb6112431b9886b2066f2e71443eb6cbad151f615cb4518262f42155787f6b36a
57296656919c975cdb2009e64dd57127e3e0d83693cfc8154372488154299f3c9f6b475d2cc9b1f42040d0657e7f474506cb36aa98d979
287d3b46093d4772f0cbf5e2e741e61c0d6edd21afc39673c72ff633b0e9dcaad4bf90dbff85fcb79b855e40ce649eda05d548a44802dc
8da787166cb4431ee7958fe384ddc6d4b9118ee2836523c11d47f0a72277609115db7b89b5d818e35a7ac26db57e725bf1999aeb56cb13
2a91489e9c4f7f78789cd1aa803e46320996dd54132cdb3268cf529bac776069ad14de2d2840c23f0224b2e771f87b926d78b5d6572356
acbfa7c1fa7a5506cb36adf571066d5b72673f648baef4b258ecd1da4a2cf7851e48c35d6b1790eadb67212f1e17ea4cd63e6e44b8a4e5
5d07d507cb364266108b426512ed6864ee4d131a9b1f867531a5d726744a6593f99548865c643f7b7a3cdb537ebeb5d643c20fc1b20382
e5125aeb235981c9b0e630ffaf52eabf9afb2fff1fd9f9e3df75f917572ceed89c6bfdd0feff9bcca345a19aede76e194ec714e4bcf8cd
d3af9e740ad4c042125b097740b0ec402a5ba70416008962fe26820b30aa3775d5f2bee779fdc41873ecdc8a3ffc132fc77ab6f46d5be0
c88ab037f6c325d5030048819d2f766e8729ed109ffcbcd8766bb21a16f0c1f70aee47a9824e425adb9df0c9f08b8ae54899ee9002a04e
cd70fa29ab5bd157e0750867c698c94226fca16239922de3db95cdb69a2915cd6dd939000052b32117c13f16543617cd670662f9d2fa4c
5ef179cc1715cbc06415dd81ecb4c28a630039b893aa260b15f09b84bdc7095e910759e7c01ce20c102c2750c8967e00ea3193fd9bb3dd
c8007e8c6c30eed3acb5608d16400921582642c2e69e844d569f0348d9ab543459215b9909ab965d35f3882f089cd32058268ec089c2dd
cd3dbd27f9196bbed7a992ffcf08411854342ba2b53e902e0439b996b0c9f48ec008961993fe9a6c3388d81e24fcd955c52fcd9f25ae32
96563d9bf21d6bc2ea16f3a5d7ba93cd1058795ea84cc3651b55f74008960593e0d9deae908a2716994938fcfdc35682c3cced2bcf36a1
bf70022f945c78dd163212602f860e38f68d47b0447beb44aa32e599b5f69ebea7829486ca471baea59ac909bc10b65d956d1754d053ba
b46d049d5bd109c112dec8c9b26dd95edf63d970f4dfcbbcd37f5ef93bf8daec1d6dffe4649dbfd6bceadd4456df86f420219345160528
60787c9efd7ceed2e2a81f8225b2c00e47bff782263c564c6bbd27a173afc085483309992caec89c4c09b92aa81a4f05b3078225922527
d1d38a86e45f5b7dd938b9a233a914ed1554e17c95f96e7c0f32274581d3423e9bef994eb41ec1124909bc176d4aae5b211e9f3caa0000
13d049444154926116782753530e0aa8f2db90b9c7707919e473799c7135f390b65aab112c31291932392d7c88fb9266bd48818c021c65
7c527f904a2655a34264ba13dd67aae9cb112c119d1c482e0a5d0d4b884436e4bb782027f6dce66c5e4bc8a4e25f182938340bd8529d4f
bccd7cf7c5089688a2d030f920fdf9181641115a27f4e3cce6367f35c69c3ab7a238092d60bb36c6ec39b7826089700a0c9397b26d1d57
a9a886cc7bcea5a23993f9980c955764c2799b6fa898bbfec9b90518c1563cb4d6575a6b7bc5f29871a8b40b064e64b843cb0fbb32a03a
76c71c63cca6fd0ed813a952ea2ce1d7c056597fd8e38f34ed8607f27a2efa79d25a5f2ce8611c959d7a648cd96d8ed5b6822dc7f0d0a8
582e40c5125e14b0f3c2ab2c223ae50a14e84602c569e2dbc53e4815938bc281a450d0c583316627b1c77e117071e8095b95baa8586230
db3b4f6bfd22079d1c43e565ab22b92995194225d091548a7612af66dad0fb28d5369a5c87f54e6b9dd4c245696cbeeddc81600896e845
6bbda3b5be9730799ed94a525bb9f8c8d036e09fbd2833c61c4d301cd9c7b9044c16dc8593dcf42739ce7f74ee4010044b7422f3686c98
fc91f8b0d7bc139960ada5b2421b202002bb4abb3537f3a32cac49c5be04cc7b59090f8f52ac0ccbb1ffceb903de112cb1949d3fd51aea
cea581f96cae2ac9f03630311932df4a3064da8be467099849cd0d4c4cdf4056cbc2292adf0b102ce16855276f3219eabe6ecd95dca22a
09a46b2e647e4e64b8fc9dac267f21607ab191da5ccb00ee984ab518c1123fc9dcc9a78caa9367ad216e567c0219b2dbe2b586cb539893
b941c0f4e6839c5392986a20efa7cff99f2c045b8260593959d9ddcc9d4c7da78d93d610f71143dc40399a3999b2bafc72e22746c0fcd3
d0eae35b996a3069089370ebb382fa9562c67204cb0acd35313f4ffc15688749fa85018593d5e50752c57c3ff17c4c02a61fe753b57b92
dff9ec715ad725db87ae46b0ac48d32a48be649f127ee6ed616ec2245029bb35636b3ee69455cc26602633b41b99af6a5f13309f42864c
299e9c06289e9c485f4cacc0ce3b15902beddbc417e2dcc9ee180c6f03584a76fbb99af878667be2eed672bc92d7fcc6b9c3af3b394fd9
9ffbaeafadd67acbbe17f2b317f073f19e3de8bb2158164cae08531eea9e4998e4cb0aa097d6bcb929fbea5ed652c1eab1ad6369ceec9c
fe4a9ffb200c8517a8b52027d550f9b5d51a885009a037998bb933f130f97e455b45a6b68b526897cd42d1b29fa67f04cb82241e28ef5a
f32699f80cc09bd6629fc3895ed5f3d217f8c88afd14b7e9f4692643de9ab994c3112c0ba0b53e4e3450da03d0a17c49ab998f04601ac6
988b090366b3c0a7d85198f6369dadd7f9c1f98b79b96bedd6c6289a07ccb1cc58c27328ab9ad80e204d131f23ed4575755bfe69adf75a
0b6952ec8d6ca74d5cb0435b3804cb0c251c284f680f042035131e33eda8cd0ecdb4ff26d305eccf56eb4f9f0bb0ec6b7edfac2eb77f52
e4888b609991482d1ffa7a95ea24c30700923661c0ac66f53840b0cc40a27d2819ee069025adb51da2de8ffcd8b908471558bc9330d93d
e05ef6f14e25545ecaeaee1d4225801cb55691c75c78d22ceea96ede25ea42c53251135d51afc2500e80e248a3f5a7c817efccbd44b1a8
5826a6d58b32955079424f2f00a59246eb365c7e8ef8146d887db4fb593bf70099a362998889ae9a57a9b2550680ba4d305a3493ea2553
8b50042a96099003d97322a1b269684ea804501d199d7913718719dbebf159fa3f02d923584ec8b60f4a68d8fb924009007f0c8f7f8df8
727cd75a5f39b7029961287c22b2dadb6753d8a1589403004bc834a5db88c76b5b29dd62681cb9a26219596b71ced4a1f28e453900b09a
542f7722ee3fbe2143e3bbce3d4006a8584632c155ef323436078001e4387e1f710fec3363cc91732b90302a9611c8a4ece78943a51d5e
794f6373001846aa97766feb93482fe11799360564838a65605a6b5ba5fc30f1c3a075100078a4b5de92ea658c6e1eccbb4436089681c8
fede3f267e18761e25f374002090c8c5838fc6985be75620210c850720bb294c192a9b616f4225000424c7d9586d896eec0250e7562021
542c3d9a6062f7225f8d316c13060011451e1ab75bed1e3bb7020920587a22ad216e267c08acf606808945ec514c0f622489a1700fb4d6
c71387cacfacf60680e949cfcbb3080f645fe6770249a16239d2c43be83cc8410c009010990b791ee111cda40512900482e540329ff229
d27c9a456c95927d6501205111bb83102e910c86c20790f994cf13854a3b97f20da11200d2668cb1235a6fa45347486fb5d64f7c1c9002
82654f5aeba309e7539e30971200f221bbf5d811ae59e0074db84412180aef41fa537e99e057dbab5d1b2839680040a622355367581c93
225876a4b5b643cf9f26f8d5ec9e03008588742e215c62320c85772057995384caaf844a0028873166cff6a00cfc84ded28a0853a162b9
86cc599962279df732f11b005018adf585ed4519f859d1441dd151b15c61a250399355df844a00289404bed095cb7dd9c00388868ae512
1385ca6b192601005420d29ccb4363cc85732b1000c172818942a56d25c49525005426d26a71a657210a82e59c894225bbe80040c5229c
7b6cdbba2dfa202334e658b6c8bedfb143e57b422500546f27f00e3d76a738568a233882a590a18877ce1d616d33340100904a62e8f672
ef64a30f201882e5df6d1f42cf6f697b9595dfeca40300f8490a0d87815f8d2f5a6b16892298eae7584a2b86bf9c3bc2619e0b0060a948
3d2edf701e420855074bb96afbeedc110ea11200b05684c53cb4b74310d5064badb5dd47f5d1b9231c422500a09348e728fa5bc2bb9a83
e58bac928b81500900e8456b7da494fa16f85563481c5e55b97847763a88152a15a11200d09731c6aee07e08fcc2b14a1c5e55172ce50a
30f4f6596ddb844a00c04007815f38bb9f78e83647a8485543e113ccabfc688ca1212d006030e93df925e02b3833c66c39b70203d4162c
636ed7c8a4680080175aebd0276bce59f0a29aa170b9e28b152a2ff98202003cfa1af8c564ae25bca8a262a9b5b67bb0fe70ee08e3c118
b313fb390200ca1661d4edab2c180206ab2558de47da079cb642008020b4d67621cf79c857d718a39d1b811e8a1f0a9755e03142a5b547
a80400842053ac5e43beb872ce04867f864aae586aad379552cfce1d619c1963f84202008289d034fdd518b3e9dc0a74547ac5f2d8b925
8c0742250020b40873203764c81d18a4d860293d2b43f6fd6a63237f00402c97817f0fc11283955cb18c55adb4abe89e9c5b0100082374
3bbb0fd24d05e8adc86029d5ca7de70eff66b4660000c4243bbacd02ff4a46e23048a915cb58d54abe7800802984ae5ab26e008314172c
23562bedee3af7cead000084771bf837d8453cbbceadc01a25562c635d657135070098840c8707ed69c9a81c86283158c658097e462374
00c0c4ae02ff7a82257a2b2a58c6eabd45cf4a004002420f87bf657538fa2aad621923589e38b70000105f8c79feccb3442fc56ce9288b
761e9d3bfc7bc330380020055aebd027f16b630c43e2e8aca48a658c6ae525a112009090bbc00f85a170f45252b08c714545337400404a
420f87db79969bceadc01245044b19067fe7dce1d71d7d2b01008989318a46d5129d9552b18c51ad0cbdcb0100007d855e19ae58c0833e
08961d196308960000002b94122c3f38b7f8157a72340000a9a26289ceb20f9691f632a55a0900481173ff9194122a9631261587de360b
0080de22b5c00b3d2a888294102cb79c5bfcbaa377250000c07a542cd78bb1e20e0000207b2504cbd0257ae6af00009214699d01d05949
3bef8442b0040000e820eb6019e34acd18f3e4dc080000000715cbd5e85f09004819db2d2229b907cbd0154b5683030052b6c9bb839450
b15c8df995008094c558bc33736e01962058020090af18154bd61aa03382e56af4b00400a4ec1def0e5242b004002043117b5852644167
044b0000f2142b5832148ece08960000e4896089e4102c0100c88cd67a33c296c63f1963180a4767044b0000f2b317e911b351087a2158
0200909f83488f986a257ac93d58f28107005425e63038e759f445c57235f6600500a4e628d6e3617e25faca3a5846f8c0b3072b002035
b186c12f9d5b8035a85802009009adb55db4f336d2a3bd706e01d6282158865cb116ab471800005dc41a067f65181c4394102cef9d5bfc
618e25002009b28563ac453ba7ce2d40072504cb17e7167f36623f19000096385e7c73100c836390128265d0527dc44dfe01005848e656
c6aa565e1b63d8c61183641f2c23cc0121580200a61673689a61700c56caaa7016f000008aa4b53e8ab812fc8e453b18a3946019f24b10
6be80100008731e6d418a3ed8f52eabd52ea4c293573fea21f31e771a240f6839afdb392799037ce1dfe7ce40a0e00902aadf58e344e1f
d3e7d2562b19a5c32845044bf5eb4bf5127015f7893186ab380040567a06cef7c698902dfc50819276de09d91a21d6f6590000786383a2
31e6c818b3d50ca7cb90fa67d9b2f1557ed725a1123e9454b164381c00006042c5542c25f4859accaca85a020000ac56d250b80a3c1cbe
efdc02000080df08963d482f310000002c5054b0942da8ae9d3bfc21580200002c515ac55205de8aeaadd69ab9960000000b14b32abc4d
6b6d5b26bc73eef06366db364cf4d4000000925562c55205de928aaa250000c00245562c55f8aae5ab3166d3b9150000a062a5562c55e0
aae586d69a2d1e0100005a8aad58aa5f554bdb34fd8373873fdbb2121d0000a07a25572c5584f6402157a003000064a5e860291bea5f3a
77f8f389853c000000bf143d14ae7e0d87db4536cfce1d7ebd31c6bcc47e6e00000029297d285c49e03b74eef02be856920000003928be
62d908dc7ec83a34c61030010040b56a0a9676b79c47e70ebf58250e0000aa55fc507843025fe821f15be7160000804a54132cd5af7069
87aaaf9d3bfcb1db3d5ec57e5e00000029a86628bc4d6b6d17f46c3877f873628c61671e00005095aa2a962d7bce2d7efda5b50efd3b00
0000925265b034c6d8b9905f9d3bfcfa2e0b86000000aa50e5507843e6437e72eef08be6e90000a00a55074bf52b5cdad5e26f9d3bfc79
35c66cc67e5e000000b1d53ac7b26dc7b9c5af0d09af00000045ab3e58ca30f57be70ebfde122e010040e9aa0f96ea57b8bc8fd03cdd86
4b1aa803008062112c85344f0fbd52fc03e1120000948a60d9628c39b5cdcd9d3bfc225c0200802255bf2a7c11adb5ad5eee2fb8cba799
31863e970000a018542c1730c61c28a52edd7bbc62410f0000280a15cb1522552e5f95525b3451070000b92358ae21f3213facfe5b5e6c
1b63a8600200806c3114be86316637c2b0b8f5a8b5de756e050000c804c1b28348732ead1badf591732b0000400608961d49b83c8bf0ab
bec9dc4e000080ac30c7b227adf5b152eaaf08bfeac118137a1f730000006f08960368ad6df5f23cc2afb22bc67758d403000072c050f8
00b2fde3c708bf6a4316f530ef120000248f8ae5085a6bbb73ce63a45f77272bd40100009244b01c496bbda994bab73be944f8750c8d03
00806431143e92dd3147f6fcbe8bf0eb181a070000c922587a22c3d427917e9d6d49742fd552000080243014ee99ec9e7313f1577e36c6
5c39b70200004446b00c402a894f32741d033d2f0100c0e4180a0f40e65d6e469a7769bdd35a1be9af09000030092a9681c9429b6f117f
e54c568ebf38f70000000444b08c6082a171ebcc18c3ea710000100dc13222adb55d64f329f2af7d6f8cb9776e050000f08c6019d904ab
c6ad07a5d42ec3e30000202416ef44668cb935c668097bb1bc534a3d6bad2f6a799d0100407c04cb89487ba0cf917ffb3eabc701004028
0c8527c0eea22355c5d8987f090000bca162998089aa97d60fadf58bd69ae6ea000060342a9689d15adf2aa53e4cf0a85e65810f154c00
003008c132415241bc8ddcf7b2c10a720000300843e109b25543d912f2648247d7ac2067881c0000f442c53271b26bcfed448b7b1443e4
0000a02b2a9689b343d2b2b8e7bd84bcd8365a8b7cf66a7f3f0000c07204cb4cb486c70f277ac436607ea70f2600005886a1f04cc92e3a
fb133ffa4b630c21130000fc44b0cc5802f32f1b7625f901f3300100a81bc1b20009054cebd018c39ee40000548860599089fb5fcea31f
2600009561f14e415a0b7ca65a41ded6f4c3b48b7d8e9c7b01004071a858162cb10aa635534aed31171300803251b12c5862154cebadf4
c4b455cc5badf596f337000040b6a8585624c10a66e35a5695331f1300808c51b1ac48ab82b92d8b6b52f1a9351ff34a56b9030080cc50
b1ac58626d8a16b9534a1d31271300803c50b1ac58b30fb93146db5d74127c253eb4e664beb09524000069a362893f48783bcfe055b996
6ae693730f00009804c1120bc9429f2b59c99d3abbe2fdd4feb000080080e9102cb196d6da6ed1b89fd12b45d004006002044b749670bb
a22e2e2568b21008008040089618446b6d2b825f327ef5ec2e401754350100f087608951a465d195ace0ce5d13362f58140400407f044b
782343e51709f7c51cea5ac2f315d54d00009623582208adf5ae84cc1c56950ff520734e2f98bb090000c112115412324f8c31c7cead00
0054846089a8b4d65b12324b9893d9f68661720040edd8d21151d94531c6985dd946f28dadf449dfc99c9d112a0100a0628984c890f971
86d5cc6d5691030040c5120931c6dc36d54ca9681e4a0ba0945d122a0100f8858a25b2a1b53e908a664a8b80deb3221c00805f0896c896
0c9d1f29a53e4df41cee6c85d5b91500804a112c510cd905682f6255f3a31dbe776e0500a052044b14ad15368f3cef08f4608cd9716e05
00a062044b5449b69f3c92d0b931e035f86c8cb9726e0500a062044ba04502a75d24b4bba2c23933c66c39b7020050398225d091844e5b
e1b44dde2f78dd000068514afdff0a065c21cc50e5d10000000049454e44ae426082
END FIGURE
